\documentclass[12pt,a4paper]{report}

%==============================================================================
% PREAMBLE AND PACKAGES (Aligned with CEK / KTU B.Tech Project Format)
%==============================================================================
\usepackage{graphicx}
\usepackage{amsmath}
\usepackage{fancyhdr}
\usepackage{cite}
\usepackage{float}
\usepackage{epsfig}
\usepackage{longtable}
\usepackage{enumerate}
\usepackage{afterpage}
\usepackage{multirow}
\usepackage{ragged2e}
\usepackage{gensymb}
\usepackage{amsfonts} 
\usepackage[left=3.5cm,top=2.0cm,right=3cm,bottom=3cm,headheight=15pt,footskip=1.2cm]{geometry}
\usepackage{setspace}           
\usepackage{txfonts}
\usepackage{standalone}
\usepackage{tikz}
\usepackage{fontawesome5}
\usetikzlibrary{calc, shapes.geometric, arrows, arrows.meta, positioning, fit, backgrounds, shadows, tikzmark}

% Color definitions for TikZ diagrams (Clinical Deep Navy, Ocean Blue & Cyan Palette)
\definecolor{lavenderDark}{RGB}{15, 44, 89}
\definecolor{lavenderDeep}{RGB}{10, 28, 58}
\definecolor{lavenderMedium}{RGB}{29, 78, 158}
\definecolor{lavenderLight}{RGB}{240, 246, 255}
\definecolor{lavenderBorder}{RGB}{180, 205, 237}
\definecolor{lavenderAccent}{RGB}{2, 132, 199}
\definecolor{prismNavy}{RGB}{15, 44, 89}
\definecolor{prismBlue}{RGB}{29, 78, 158}
\definecolor{prismTeal}{RGB}{13, 148, 136}
\definecolor{prismMint}{RGB}{236, 253, 245}
\definecolor{prismGold}{RGB}{217, 119, 6}
\definecolor{prismDark}{RGB}{15, 23, 42}
\definecolor{prismLight}{RGB}{248, 250, 252}
\definecolor{prismIndigo}{RGB}{30, 58, 138}
\definecolor{prismSlate}{RGB}{100, 116, 139}
\definecolor{statusFail}{RGB}{220, 38, 38}
\definecolor{cInput}{RGB}{241, 245, 249}
\definecolor{cStem}{RGB}{224, 242, 254}
\definecolor{cStage1}{RGB}{209, 250, 229}
\definecolor{cStage2}{RGB}{254, 243, 199}
\definecolor{cStage3}{RGB}{254, 226, 226}
\definecolor{cStage4}{RGB}{207, 250, 254}
\definecolor{cHead}{RGB}{219, 234, 254}
\definecolor{cOutput}{RGB}{255, 228, 230}
\definecolor{cDeliverable}{RGB}{204, 251, 241}

\usepackage{subfiles}
\usepackage{adjustbox}
\usepackage{tcolorbox}
\tcbuselibrary{skins,breakable}

% Formal Algorithm Box Environment for Scholarly System Specifications
\newcounter{algocounter}[chapter]
\renewcommand{\thealgocounter}{\thechapter.\arabic{algocounter}}
\newenvironment{formalalgo}[2]{%
  \refstepcounter{algocounter}%
  \vspace{5pt}%
  \noindent%
  \begin{tcolorbox}[enhanced, breakable, colback=white, colframe=prismNavy, boxrule=0.85pt, arc=3pt,
    left=6pt, right=6pt, top=4pt, bottom=4pt,
    title={\textbf{Algorithm \thealgocounter:} #1},
    coltitle=white, colbacktitle=prismNavy, fonttitle=\bfseries\small,
    before upper={\setstretch{1.08}\raggedright}]%
    \label{#2}%
    \small%
}{%
  \end{tcolorbox}%
  \vspace{5pt}%
}

\newcommand{\Usefont}[1]{\fontfamily{#1}\selectfont}

\usepackage{lscape}
\linespread{1.5}

\usepackage{url}
\usepackage{subcaption}
\usepackage[intoc, english]{nomencl}
\usepackage[nottoc]{tocbibind}
\usepackage{booktabs}
\usepackage{array}
\usepackage{tabularx}
\usepackage{hyperref}

% Compact chapter-opening layout
\makeatletter
\def\@makechapterhead#1{%
  \vspace*{0pt}%
  {\parindent\z@ \raggedright \normalfont
    \ifnum\c@secnumdepth>\m@ne
      \huge\bfseries \@chapapp\space \thechapter
      \par\nobreak
      \vskip 10pt
    \fi
    \interlinepenalty\@M
    \Huge\bfseries #1\par\nobreak
    \vskip 15pt
  }}
\makeatother

\hypersetup{
    plainpages=false,
    pdfpagelabels=true,
    colorlinks=true,
    linkcolor=black,
    citecolor=black,
    urlcolor=black,
    pdftitle={AI-Powered Legal Information Retrieval using Retrieval-Augmented Generation},
    pdfauthor={Aleena AnilKumar, Abhinaya P S, Adithya Dinesh, Fathima A}
}

% Spacing and line-breaking settings to prevent overflow and overfull hboxes
\setlength{\headheight}{15pt}
\hyphenpenalty=1000
\tolerance=2000
\emergencystretch=2.5em

\bibliographystyle{IEEEtran}
\renewcommand{\bibname}{References}

% Figure search path
\graphicspath{{figures/}{diagrams/}{./}}

%==============================================================================
% METADATA DEFINITIONS (Project Group & CEK)
%==============================================================================
\gdef \title{AI-POWERED LEGAL INFORMATION RETRIEVAL USING RETRIEVAL-AUGMENTED GENERATION}
\gdef \shorttitle{AI-Powered Legal Information Retrieval using RAG}

\gdef \authorA{ALEENA ANILKUMAR}
\gdef \rollnoA{KNP23CS025}

\gdef \authorB{ABHINAYA P S}
\gdef \rollnoB{KNP23CO007}

\gdef \authorC{ADITHYA DINESH}
\gdef \rollnoC{KNP23CS016}

\gdef \authorD{FATHIMA A}
\gdef \rollnoD{KNP23CS047}

\gdef \acadyear{2026--2027}
\gdef \submonth{October 2026}
\gdef \month{October 2026}
\gdef \date{28-10-2026}
\gdef \dept{Computer Science and Engineering}
\gdef \deptabbr{Dept.\ of CSE}
\gdef \degree{Bachelor of Technology}
\gdef \branch{Computer Science and Engineering}
\gdef \college{College of Engineering Karunagappally}
\gdef \collegeplace{Kollam}
\gdef \collegetown{Karunagappally}

\gdef \guide{Ms. Hridhya J}
\gdef \guidedes{Assistant Professor}
\gdef \coordinator{Dr. Smitha P}
\gdef \coordinatordes{Associate Professor and HOD}
\gdef \semcordinatorA{Dr. Smitha P}
\gdef \semcordinatorAdes{Associate Professor and HOD}
\gdef \hod{Dr. Smitha P}
\gdef \hoddes{Associate Professor and HOD}

\newif\ifindividualreport
\ifdefined\candidateName
  \individualreporttrue
\else
  \individualreportfalse
\fi

\providecommand{\candidateName}{\authorA}
\providecommand{\candidateRoll}{\rollnoA}

%==============================================================================
% DOCUMENT BODY
%==============================================================================

\begin{document}

%------------------------------------------------------------------------------
% FRONT MATTER
%------------------------------------------------------------------------------
\hypersetup{pageanchor=false}
%==================================coverpage.tex================================

\thispagestyle{empty}
\begin{titlepage}
\subfile{design_cover}
	\begin{center}\fontfamily{ptm}\selectfont
		{\Large \bfseries \title \par}
		\vspace*{12pt}
		{\large \Usefont{pzc}
			A PROJECT REPORT \par
			\vspace*{5pt}
			Submitted to the APJ Abdul Kalam Technological University \par
			\vspace*{5pt}
			in partial fulfillment of requirements for the award of degree\par}
		\vspace*{8pt}
		{\large \bfseries Bachelor of Technology}\par
		{\large in}\par
		{\large \bfseries \branch}\par
		\vspace*{10pt}
		{\large by}\par
		\vspace*{6pt}
		\begin{tabular}{l@{\hspace{1.8cm}}l}
			\large\bfseries ALEENA ANILKUMAR & \large\bfseries KNP23CS025 \\
			\large\bfseries ABHINAYA P S     & \large\bfseries KNP23CO007 \\
			\large\bfseries ADITHYA DINESH   & \large\bfseries KNP23CS016 \\
			\large\bfseries FATHIMA A        & \large\bfseries KNP23CS047
		\end{tabular}\par
		\vspace*{14pt}
		{\includegraphics[width=4.0cm]{cek-logo.png}}\par
		\vspace{\stretch{0.5}}
		{\small \bfseries DEPARTMENT OF COMPUTER SCIENCE AND ENGINEERING}\par
		\vspace*{4pt}
		{\bfseries COLLEGE OF ENGINEERING KARUNAGAPPALLY}\par
		\vspace*{4pt}
		{\bfseries KOLLAM, KERALA}\par
		\vspace*{4pt}
		{\bfseries \submonth}\par
	\end{center}
\end{titlepage}

%=============================vision_mission.tex===============================
\thispagestyle{empty}

\begin{center}
    \vspace*{0.2cm}
    \includegraphics[width=3.6cm]{cek-logo.png}\par
    \vspace*{0.5cm}
    {\bfseries\large COLLEGE OF ENGINEERING KARUNAGAPPALLY\par}
    \vspace*{0.2cm}
    {\bfseries Thodiyoor -- 690523\par}
    \vspace*{0.2cm}
    {\bfseries\large DEPARTMENT OF COMPUTER SCIENCE AND ENGINEERING\par}
    
    \vspace*{1.4cm}
    {\bfseries\large\underline{Vision of Department}\par}
    \vspace*{0.45cm}
    \begin{minipage}{0.92\linewidth}
        \centering
        To be renowned centre for education, research, and innovation in the frontier areas of Computer Science and Engineering.
    \end{minipage}\par

    \vspace*{1.4cm}
    {\bfseries\large\underline{Mission of Department}\par}
    \vspace*{0.5cm}
\end{center}

\noindent\hspace{0.02\linewidth}
\begin{minipage}{0.94\linewidth}
    \begin{itemize}
        \setlength\itemsep{1.1em}
        \item Establish and maintain an operational environment to acquire, impart, create and apply knowledge in Computer Science and Engineering and inter-disciplinary.
        \item Serve as a resource centre for innovation in design and development of hardware and software.
        \item Inculcate leadership qualities, professional ethics and a sense of social commitment.
    \end{itemize}
\end{minipage}

\vfill
\clearpage

%==================================certificate1.tex================================
% Certificate page for B.Tech Project Report

\newpage
\thispagestyle{empty}

\begin{spacing}{1.18}

\begin{center}
	\textbf{DEPARTMENT OF COMPUTER SCIENCE AND ENGINEERING}\\
	\textbf{COLLEGE OF ENGINEERING KARUNAGAPPALLY}\\
	\vspace*{3pt}
	\textbf{\acadyear}
\end{center}

\vspace*{6pt}
\begin{center}
	\includegraphics[width=3.6cm]{cek-logo.png}
\end{center}

\vspace*{6pt}
\begin{center}
	\textbf{\Large CERTIFICATE}
\end{center}
\vspace*{8pt}

\noindent This is to certify that the report entitled \textbf{\title} submitted by \textbf{ALEENA ANILKUMAR} (\rollnoA), \textbf{ABHINAYA P S} (\rollnoB), \textbf{ADITHYA DINESH} (\rollnoC), and \textbf{FATHIMA A} (\rollnoD), to the Department of \branch\ in partial fulfillment of the requirements for the award of the degree of Bachelor of Technology in \textbf{\branch}\ of the APJ Abdul Kalam Technological University is a bonafide record of the project work carried out by them under our guidance and supervision. This report in any form has not been submitted to any other University or Institute for any purpose.

\vspace*{2.0cm}

\noindent\begin{minipage}[t]{0.48\textwidth}
	\raggedright
	\textbf{Hridhya J} \\
	(Project Guide) \\
	Assistant Professor \\
	Dept. of CSE \\
	CE, Karunagappally
\end{minipage}\hfill
\begin{minipage}[t]{0.48\textwidth}
	\raggedright
	\textbf{Dr. Smitha P} \\
	(Project Coordinator) \\
	Associate Professor and HOD \\
	Dept. of CSE \\
	CE, Karunagappally
\end{minipage}

\end{spacing}

\thispagestyle{empty}
%============================================================================

\ifindividualreport
  %==================================certificate_individual.tex================================
% Individual Certificate page for B.Tech Project Report (KTU / CEK Format)

\newpage
\thispagestyle{empty}

\begin{spacing}{1.18}

\begin{center}
	\textbf{DEPARTMENT OF COMPUTER SCIENCE AND ENGINEERING}\\
	\textbf{COLLEGE OF ENGINEERING KARUNAGAPPALLY-690523}\\
	\vspace*{3pt}
	\textbf{\acadyear}
\end{center}

\vspace*{6pt}
\begin{center}
	\includegraphics[width=3.6cm]{cek-logo.png}
\end{center}

\vspace*{6pt}
\begin{center}
	\textbf{\Large CERTIFICATE}
\end{center}
\vspace*{8pt}

{\sloppy
\hyphenpenalty=10000
\exhyphenpenalty=10000
\noindent This is to certify that the report entitled \textbf{\title} submitted by \textbf{\candidateName} (\candidateRoll) to the APJ Abdul Kalam Technological University in partial fulfillment of the B.Tech. degree in Computer Science and Engineering is a bonafide record of the project work (CSD 416 Project Phase II) carried out by them under our guidance and supervision. This report in any form has not been submitted to any other University or Institute for any purpose.\par}

\vspace*{2.0cm}

\noindent\begin{minipage}[t]{0.48\textwidth}
	\raggedright
	\textbf{Ms. Hridhya J} \\
	(Project Guide) \\
	Assistant Professor \\
	Dept. of CSE \\
	CE, Karunagappally
\end{minipage}\hfill
\begin{minipage}[t]{0.48\textwidth}
	\raggedright
	\textbf{Dr. Smitha P} \\
	(Project Coordinator) \\
	Associate Professor and HOD \\
	Dept. of CSE \\
	CE, Karunagappally
\end{minipage}

\end{spacing}

\thispagestyle{empty}
%============================================================================

\fi
%==================================declaration.tex================================
\newpage
\thispagestyle{empty}

\begin{center}
\vspace*{24pt}
\textbf{\Large DECLARATION}\\[18pt]
\end{center}

{\fontsize{12.5pt}{18pt}\selectfont
\begin{spacing}{1.15}
\noindent We, \textbf{ALEENA ANILKUMAR} (\rollnoA), \textbf{ABHINAYA P S} (\rollnoB), \textbf{ADITHYA DINESH} (\rollnoC), and \textbf{FATHIMA A} (\rollnoD), hereby declare that the project report entitled \textbf{\title}, submitted for partial fulfillment of the requirements for the award of the degree of Bachelor of Technology in \branch\ of the APJ Abdul Kalam Technological University, Kerala, is a bonafide work done by us under the supervision of \mbox{\textbf{\guide}} (\guidedes, \deptabbr, \college).\par
\vspace{10pt}
\noindent This submission represents our ideas in our own words, and where ideas or words of others have been included, we have adequately and accurately cited and referenced the original sources.\par 
\vspace{10pt}
\noindent We also declare that we have adhered to the ethics of academic honesty and integrity and have not misrepresented or fabricated any data, idea, fact, or source in our submission. We understand that any violation of the above will be a cause for disciplinary action by the Institute and/or the University and can also evoke penal action from the sources which have thus not been properly cited or from whom proper permission has not been obtained. This report has not previously formed the basis for the award of any degree, diploma, or similar title of any other University or Institute.
\end{spacing}}

\vspace{1.55cm}

\noindent \begin{minipage}{0.40\linewidth}
\fontsize{12.5pt}{18pt}\selectfont
\begin{flushleft}
\textbf{\collegetown} \\
\textbf{\date}
\end{flushleft} 
\end{minipage}
\hfill
\begin{minipage}{0.56\linewidth}
\fontsize{12.5pt}{18pt}\selectfont
\begin{flushright}                                      
\textbf{ALEENA ANILKUMAR} (\rollnoA) \\[4pt]
\textbf{ABHINAYA P S} (\rollnoB) \\[4pt]
\textbf{ADITHYA DINESH} (\rollnoC) \\[4pt]
\textbf{FATHIMA A} (\rollnoD)
\end{flushright} 
\end{minipage}

\thispagestyle{empty}

\hypersetup{pageanchor=true}

\cleardoublepage
\pagenumbering{roman}
\pagestyle{plain}

%============================= abstract.tex================================
\chapter*{Abstract}%
\addcontentsline{toc}{chapter}{Abstract}%
\vspace{-25pt}

\begin{spacing}{1.2}
\noindent Access to Indian legal information remains hindered by complex statutory language, fragmented document repositories, and language barriers for non-English speakers. Citizens seeking redress through civic complaint systems face difficulties understanding relevant laws such as the Bharatiya Nyaya Sanhita (BNS), Bharatiya Nagarik Suraksha Sanhita (BNSS), Right to Information (RTI) Act, and Consumer Protection Act. Traditional search interfaces return raw legal text without context, requiring professional legal interpretation. To address these challenges, this project presents \textbf{CiviQ}: an AI-Powered Legal Information Retrieval system using \textbf{Retrieval-Augmented Generation (RAG)}. The system ingests PDF legal documents, extracts text using \textbf{PyMuPDF}, and applies overlapping word-chunking (900 words per chunk, 150-word overlap, 40-word minimum). Text chunks are embedded using \textbf{sentence-transformers (all-MiniLM-L6-v2)} generating 384-dimensional cosine-normalized vectors, indexed in \textbf{Pinecone} vector database with batch upserts.

\vspace{8pt}

\noindent User queries are embedded identically, and top-k semantic retrieval (default k=5) fetches relevant passages. \textbf{Google Gemini 2.5 Flash} generates grounded answers citing retrieved sources. The system detects Malayalam queries via Unicode block analysis (U+0D00--U+0D7F) and translates them to English before retrieval, with response translation back to Malayalam. The production architecture comprises a \textbf{FastAPI} backend (Python 3.12) with endpoints for question answering, document upload analysis, and health checks, serving a \textbf{React 19 + Vite 8 + Tailwind CSS 4} frontend with voice input (browser speech-to-text), PDF upload, language selection (auto/English/Malayalam), top-k slider, source citation display, and document summaries. The corpus includes BNS 2023, BNSS 2023, BSA 2023, Consumer Protection Act 2019, Indian Contract Act 1872, RTI Act 2005, and Constitution of India. The system provides accessible legal guidance for civic complaint resolution and information rights.
\end{spacing}

\thispagestyle{plain}
%=======================================================================

%==================================acknowledgement.tex=============================
\chapter*{Acknowledgement}%
\addcontentsline{toc}{chapter}{Acknowledgement}%
\vspace{-15pt}

\begin{spacing}{1.25}
\noindent We take this opportunity to express our deepest gratitude and sincere thanks to everyone who helped us to complete this project work and report successfully.

\vspace{6pt}

\noindent We express our sincere thanks to \mbox{\textbf{Dr.~Smitha Dharan}}, Principal, \college, Kollam, and \mbox{\textbf{\hod}}, Head of Department, \dept, \college, \collegeplace, for providing us with all the necessary infrastructural facilities, high-performance computing resources, and institutional encouragement.

\vspace{6pt}

\noindent We would like to express our profound sense of gratitude and indebtedness to our project guide \mbox{\textbf{\guide}}, \guidedes, \dept, \college, for her insightful guidance, constant encouragement, rigorous critiques, and invaluable mentorship throughout the entire lifecycle of this project.

\vspace{6pt}

\noindent We also express our sincere gratitude and indebtedness to our project coordinator \mbox{\textbf{\coordinator}}, \coordinatordes, \dept, \college, for her continuous coordination, constructive reviews, and valuable support during each phase of this project.

\vspace{6pt}

\noindent We also convey our sincere appreciation to all the faculty members and technical staff of the \dept\ for their kind cooperation, constructive academic advice, and technical support during our design and development phases.

\vspace{6pt}

\noindent Finally, we express our wholehearted thanks to our parents, family members, and friends who provided unwavering moral support, patience, and encouragement that contributed immensely to the successful realization of this project.

\vspace*{14pt}
\begin{flushright}
	\textbf{ALEENA} (\rollnoA) \\[3pt]
	\textbf{FATHIMA} (\rollnoB) \\[3pt]
	\textbf{ADITHYA} (\rollnoC) \\[3pt]
	\textbf{ABHINAYA} (\rollnoD)
\end{flushright}
\end{spacing}

\thispagestyle{plain}


\cleardoublepage
{\setstretch{1.15}\tableofcontents}
\cleardoublepage
\listoffigures
\cleardoublepage
\listoftables
\cleardoublepage

%==================================symbol.tex================================
% List of Symbols and Abbreviations
\chapter*{List of Symbols and Abbreviations}
\addcontentsline{toc}{chapter}{List of Symbols and Abbreviations}%
\vspace{-25pt}

\begin{spacing}{1.02}
\noindent \textbf{Abbreviations}
\vspace{4pt}

\noindent\begin{tabular*}{\textwidth}{@{\extracolsep{\fill}}p{2.8cm} p{11.2cm}@{}}
\textbf{RAG} & Retrieval-Augmented Generation \\
\textbf{LLM} & Large Language Model \\
\textbf{NLP} & Natural Language Processing \\
\textbf{ANN} & Approximate Nearest Neighbor \\
\textbf{HNSW} & Hierarchical Navigable Small World \\
\textbf{BERT} & Bidirectional Encoder Representations from Transformers \\
\textbf{SBERT} & Sentence-BERT (Sentence Transformers) \\
\textbf{BNS} & Bharatiya Nyaya Sanhita, 2023 \\
\textbf{BNSS} & Bharatiya Nagarik Suraksha Sanhita, 2023 \\
\textbf{BSA} & Bharatiya Sakshya Adhiniyam, 2023 \\
\textbf{RTI} & Right to Information Act, 2005 \\
\textbf{CPA} & Consumer Protection Act, 2019 \\
\textbf{ICA} & Indian Contract Act, 1872 \\
\textbf{API} & Application Programming Interface \\
\textbf{REST} & Representational State Transfer \\
\textbf{ASGI} & Asynchronous Server Gateway Interface \\
\textbf{HTTP} & Hypertext Transfer Protocol \\
\textbf{CORS} & Cross-Origin Resource Sharing \\
\textbf{JSON} & JavaScript Object Notation \\
\textbf{PDF} & Portable Document Format \\
\textbf{SRS} & System Requirements Specification \\
\textbf{DFD} & Data Flow Diagram \\
\textbf{UI/UX} & User Interface / User Experience \\
\textbf{HMR} & Hot Module Replacement \\
\textbf{SPA} & Single-Page Application \\
\end{tabular*}

\vspace{6pt}
\noindent \textbf{Mathematical Notations}
\nopagebreak[4]
\vspace{4pt}
\nopagebreak[4]

\noindent\begin{tabular*}{\textwidth}{@{\extracolsep{\fill}}p{2.8cm} p{11.2cm}@{}}
$\mathcal{C}$ & Complete legal statutory document corpus ($\mathcal{C} = \{d_1, d_2, \dots, d_N\}$) \\
$d_i$ & Individual legal statutory document in the corpus \\
$c_{i, j}$ & The $j$-th overlapping text chunk generated from document $d_i$ \\
$\mathcal{T}$ & Arbitrary-length input text token domain \\
$\phi(\cdot)$ & Sentence transformer embedding mapping function ($\phi: \mathcal{T} \to \mathbb{R}^{384}$) \\
$\mathbf{e}(t)$ & L2-normalized 384-dimensional dense semantic embedding vector \\
$q$ & User input legal query string in natural language \\
$q_{\text{en}}$ & English-translated query string for multilingual processing \\
$q_{\text{ml}}$ & Original input query string in Malayalam \\
$\text{sim}(q, c)$ & Cosine similarity retrieval score between query $q$ and chunk $c$ \\
$k$ & Number of top semantically similar passages retrieved ($k \in [1, 10]$) \\
$R(q, k)$ & Retrieval operator selecting top-$k$ relevant statutory text chunks \\
$G(\cdot)$ & Large language model generation function conditioned on retrieved evidence \\
$\delta(q)$ & Language detection indicator for Malayalam Unicode range (U+0D00--U+0D7F) \\
$T(\cdot)$ & Bidirectional language translation mapping function \\
$s$ & Chunk stride parameter in word tokens ($s = 750$ words) \\
$o$ & Overlap window parameter in word tokens ($o = 150$ words) \\
$c_{\text{size}}$ & Fixed maximum chunk size in word tokens ($c_{\text{size}} = 900$ words) \\
$\text{Precision}@k$ & Proportion of retrieved passages in top-$k$ relevant to query ground-truth \\
$\text{MRR}$ & Mean Reciprocal Rank of the first relevant retrieved statutory passage \\
\end{tabular*}
\end{spacing}

\thispagestyle{plain}
%=======================================================================


%------------------------------------------------------------------------------
% MAIN BODY (Arabic Numbering & Header/Footer)
%------------------------------------------------------------------------------
\cleardoublepage
\setcounter{page}{1}
\pagenumbering{arabic}

\pagestyle{fancy}
\fancyhf{}
\fancyhead[L]{\fontsize{9}{11}\selectfont \textit{\shorttitle}}
\fancyhead[R]{\fontsize{9}{11}\selectfont \textit{\deptabbr, CEK}}
\fancyfoot[C]{\thepage}
\renewcommand{\headrulewidth}{0.4pt}
\renewcommand{\footrulewidth}{0pt}

%==============================================================================
% CHAPTER 1: INTRODUCTION
%==============================================================================
\chapter[Introduction]{\Huge Introduction}

Access to accurate, comprehensible legal information is fundamental to civic empowerment and the effective exercise of constitutional rights in India. Citizens encounter legal systems daily—from consumer disputes and service grievances to information access rights under the Right to Information (RTI) Act 2005 and criminal proceedings under the Bharatiya Nyaya Sanhita (BNS) 2023. However, contemporary legal information retrieval faces systemic obstacles:

\begin{enumerate}[(i)]
    \item \textbf{Statutory Language Complexity:} Legal texts employ specialized terminology, nested conditional clauses, and cross-references that are impenetrable to non-lawyers.
    \item \textbf{Fragmented Document Repositories:} Statutes, amendments, judicial interpretations, and procedural codes are scattered across government portals, parliamentary archives, and district registries without unified search.
    \item \textbf{Language Barriers:} Over 57\% of Indian citizens speak regional languages as their primary tongue, yet most legal resources remain exclusively in English or technical Hindi, excluding Malayalam, Tamil, Telugu, and other linguistic communities.
    \item \textbf{Lack of Contextual Answers:} Traditional keyword search returns raw statutory clauses without explanation of applicability, penalties, or procedural steps, requiring expensive legal consultation.
\end{enumerate}

To address these challenges, this project presents \textbf{CiviQ – AI-Powered Civic Issue Reporting and Complaint Management System}, a Retrieval-Augmented Generation (RAG) platform for citizen-accessible legal information retrieval. The system ingests comprehensive legal corpora including BNS 2023, BNSS 2023, BSA 2023, Consumer Protection Act 2019, Indian Contract Act 1872, RTI Act 2005, and the Constitution of India. Documents are preprocessed through overlapping word-chunking (900-word chunks, 150-word overlap), embedded via \textbf{sentence-transformers (all-MiniLM-L6-v2, 384-dim vectors)}, and indexed in \textbf{Pinecone} vector database. User queries—submitted as text or voice input—undergo identical embedding, semantic retrieval of top-k relevant passages (default k=5), and grounded answer generation via \textbf{Google Gemini 2.5 Flash}. The system automatically detects Malayalam queries (Unicode U+0D00–U+0D7F), translates them to English for retrieval, and translates responses back, ensuring linguistic inclusivity.

\section{Background and Overview}
In democratic societies, access to legal knowledge enables citizens to exercise their rights, challenge injustices, and engage with governance mechanisms. The Right to Information Act 2005, Consumer Protection Act 2019, and the 2023 reformed criminal codes (BNS, BNSS, BSA) form the bedrock of citizen-state interaction. Yet, navigating these statutes requires understanding complex legal frameworks.

Retrieval-Augmented Generation (RAG) represents a transformative approach to question answering. Unlike parametric language models that generate answers solely from learned parameters (prone to hallucination), RAG systems first retrieve relevant passages from external document corpora, then condition generation on those retrieved sources. This grounding mechanism ensures answers cite verifiable statutory text.

The all-MiniLM-L6-v2 sentence transformer model provides an optimal balance for semantic embedding: 384-dimensional vectors produced via mean-pooling over BERT representations, fine-tuned on semantic textual similarity tasks. Combined with Pinecone's HNSW (Hierarchical Navigable Small World) indexing, the system achieves sub-100ms retrieval latency over million-scale document chunks.

\section{Standard Information Retrieval Baselines and Their Limitations}
Traditional legal search systems exhibit fundamental limitations:
\begin{itemize}\setlength{\itemsep}{3pt}\setlength{\parskip}{0pt}\setlength{\parsep}{0pt}
    \item \textbf{Keyword-Based Boolean Search:} Legacy systems (Manupatra, SCC Online) rely on exact keyword matching with boolean operators. These fail for natural language queries like ``What are my rights if a product is defective?'' without prior knowledge of relevant statutory terms.
    \item \textbf{Parametric LLMs without Grounding:} Pure GPT-4 or Gemini queries generate plausible-sounding legal answers but frequently hallucinate non-existent sections, penalties, or procedures. Without retrieval grounding, users cannot verify statutory citations.
    \item \textbf{Absence of Multilingual Support:} Existing legal databases serve primarily English-language queries. Malayalam speakers, representing 38 million citizens in Kerala, face exclusion from accessible legal information.
\end{itemize}

\section{Proposed Solution: RAG-Based Legal Information Retrieval}
The proposed system establishes an integrated legal assistance pipeline. PDF legal documents undergo automated text extraction via \textbf{PyMuPDF (fitz)}, removing headers, footers, and page artifacts. Text is split into overlapping chunks (900 words per chunk with 150-word overlap, minimum 40 words) to preserve cross-boundary legal context. Each chunk is embedded using the all-MiniLM-L6-v2 model, producing 384-dimensional L2-normalized vectors.

These embeddings are uploaded to Pinecone vector database in batches of 100, indexed with cosine similarity metric. User queries follow an identical embedding pipeline. Top-k semantically similar passages (default k=5, user-configurable 1–10) are retrieved via cosine similarity search. Retrieved passages are concatenated with the user query into a structured prompt for Google Gemini 2.5 Flash, which generates a grounded answer with explicit citation of source statutes and sections.

For Malayalam queries, the system detects Unicode characters in the Malayalam block (U+0D00–U+0D7F), translates the query to English via Gemini API, performs retrieval in English, generates an English answer, then translates the response back to Malayalam before display.

\section{Objectives of the Project}
The primary scientific, algorithmic, and engineering objectives of this project are:
\begin{enumerate}
    \item To curate and ingest a comprehensive Indian legal corpus comprising post-2023 reformed criminal codes (BNS, BNSS, BSA) and foundational civil statutes.
    \item To engineer a robust PDF text extraction pipeline using PyMuPDF with automated cleaning of headers, footers, and page artifacts.
    \item To implement overlapping word-chunking (chunk\_size=900, overlap=150, min\_chunk=40) that preserves legal context across chunk boundaries.
    \item To embed text chunks using sentence-transformers all-MiniLM-L6-v2 with cosine normalization and batch upload to Pinecone vector database.
    \item To design bilingual query processing with automatic Malayalam detection (Unicode block U+0D00–U+0D7F) and Gemini-powered translation.
    \item To integrate Google Gemini 2.5 Flash for context-grounded answer generation with explicit source citation.
    \item To build a FastAPI asynchronous backend with endpoints for question answering, document analysis, and system health.
    \item To develop a React 19 + Vite 8 frontend with voice input, PDF upload, language selection, top-k tuning, and source citation display.
    \item To evaluate retrieval precision and answer accuracy across RTI Act, Contract Act, BNS, BSA, and BNSS query benchmarks.
\end{enumerate}

\section{Scope and Impact}
The platform is designed as an educational and informational legal assistance tool for citizens, students, journalists, and civic organizations. It provides grounded answers with statutory citations for preliminary legal research. Crucially, the system is explicitly scoped as an information retrieval aid, not a substitute for professional legal counsel, and displays prominent disclaimers that outputs do not constitute legal advice.

\section{Organization of the Report}
The remainder of this report is organized as follows:
\begin{itemize}
    \item \textbf{Chapter 2 (Literature Review):} Reviews foundational advancements in Retrieval-Augmented Generation, vector databases, sentence transformers, large language models, multilingual NLP, and legal AI systems.
    \item \textbf{Chapter 3 (Problem Definition):} Formulates the legal information access challenge, defines RAG constraints, delineates system boundaries, and establishes the mathematical formulation for embedding and retrieval.
    \item \textbf{Chapter 4 (Requirements Specification):} Outlines the formal System Requirements Specification (SRS), detailing functional requirements (text/voice query, PDF upload, bilingual support), non-functional quality attributes, and technology stack specifications.
    \item \textbf{Chapter 5 (System Design):} Details the multi-pipeline architecture (ingestion, query-answering, document-analysis), system workflows, operational algorithms, and data flow diagrams.
    \item \textbf{Chapter 6 (Design and Methodology):} Expounds the mathematical formulations of cosine similarity retrieval, overlapping chunking, embedding normalization, RAG configuration, and Gemini prompt engineering.
    \item \textbf{Chapter 7 (Preliminary Analysis):} Provides in-depth technical analysis and comparative evaluation of embedding models, vector databases, LLM providers, and frontend frameworks.
    \item \textbf{Chapter 8 (Feasibility Study):} Assesses the technical viability, economic sustainability, and project schedule feasibility, accompanied by the project Gantt chart.
    \item \textbf{Chapter 9 (Conclusion):} Summarizes the major achievements and discusses the operational limitations of the completed system.
\end{itemize}


%==============================================================================
% CHAPTER 2: LITERATURE REVIEW
%==============================================================================
\chapter{Literature Review}

Legal information retrieval using Retrieval-Augmented Generation sits at the intersection of natural language processing, vector databases, large language models, multilingual translation, and civic technology. This chapter synthesizes recent advancements, critical methodologies, and benchmark investigations established across contemporary academic literature.

\section{Retrieval-Augmented Generation (RAG) Architectures}
Retrieval-Augmented Generation combines parametric knowledge in large language models with non-parametric retrieval over external document corpora. RAG architectures encode documents into dense vector representations, retrieve top-$k$ semantically similar passages via cosine similarity, and condition generation on retrieved context, reducing hallucination and grounding answers in verifiable sources. In the legal domain, Rahman et~al.~\cite{ref_rahman2025_legalqueryrag} introduced Legal Query RAG (LQ-RAG), demonstrating that recursive feedback mechanisms substantially reduce hallucinations and improve reasoning over complex statutory texts. Mubeen et~al.~\cite{ref_mubeen2025_legalease} developed LegalEase, a specialized RAG-based framework for Indian law integrating dense retrieval over the Bharatiya Nyaya Sanhita (BNS) and Indian Penal Code with multi-LLM comparative evaluations. Furthermore, Sha et~al.~\cite{ref_sha2026_rhetorical} investigated rhetorical-role-aware RAG for Indian Supreme Court judgments, showing that structured chunking and fusion-based reranking significantly enhance contextual recall.

\section{Dense Vector Representations and Sentence Transformers}
Sentence-transformers adapt transformer architectures through Siamese networks to evaluate semantic textual similarity. The \texttt{all-MiniLM-L6-v2} model produces 384-dimensional embeddings optimized for semantic search with minimal computational overhead, making it suitable for production deployment. Unlike traditional bag-of-words or TF-IDF methods, dense embeddings capture semantic relationships between queries and documents even when exact keyword matches are absent. In legal document matching, Dhanani et~al.~\cite{ref_dhanani2022_scalable} demonstrated that incorporating pre-learned legal semantic word embeddings into document representations significantly improves judgment recommendation accuracy and retrieval scalability across distributed clusters.

\section{Vector Databases, Knowledge Graphs, and Semantic Search}
Modern managed vector database platforms provide approximate nearest neighbor (ANN) search using Hierarchical Navigable Small World (HNSW) graphs, achieving logarithmic query complexity that enables sub-100ms retrieval over million-scale corpora. To further enrich vector-based retrieval, Linders and Tomczak~\cite{ref_linders2025_kg_rag} demonstrated that extending RAG architectures with structured knowledge graphs enables symbolic entity reasoning, resolving relational ambiguities that unstructured text chunking alone cannot address.

\section{Large Language Models and Data Synthesis for Legal QA}
Modern foundation models such as Gemini 2.5 Flash, with long-context reasoning capabilities and native multilingual support, are optimized for production RAG pipelines requiring low-latency generation with citation grounding. Modern legal QA systems also increasingly leverage synthetic training and distillation: Ren et~al.~\cite{ref_ren2026_autodata} developed automated data synthesis pipelines tailored for legal RAG, overcoming legal data scarcity through synthesized question-answer pairs and LLM-as-a-judge evaluation. Italiani et~al.~\cite{ref_italiani2025_enhancing} proposed AceAttorney, distilling knowledge from frozen LLMs into compact, specialized student models with selective generative retrieval. For judicial reasoning, Song et~al.~\cite{ref_song2025_summarization} established a judicial syllogism framework that guides LLMs to explicitly separate legal premises, case facts, and judgment deductions.

\section{Legal AI and Indian Judicial Systems}
Domain-specific pre-training and legal artificial intelligence benchmarks have expanded rapidly across jurisdictions. Within the Indian legal framework, Nigam et~al.~\cite{ref_nigam2026_ailqa} introduced the AILQA benchmark, conducting rigorous empirical evaluations on Indian legal question answering and the All India Bar Examination (AIBE). Thambi et~al.~\cite{ref_thambi2024_lawgic} introduced LAWGIC, utilizing latent Dirichlet allocation topic modeling on Supreme Court decisions from Indian Kanoon for automated document classification and retrieval. Garg et~al.~\cite{ref_garg2026_multifaceted} investigated multifaceted structural classification across 15 semantic categories on Indian court decisions, while Bhattacharya et~al.~\cite{ref_bhattacharya2025_ardi} formulated the ARDI benchmark for automatic advocate recommendation across thousands of statutory sections. In addition, Singh and Dwivedi~\cite{ref_singh2026_transformer_xai} proposed a transformer-XAI framework with bias attribution and adaptive RAG tailored specifically for the Indian judicial system.

\section{Multilingual NLP and Translation for Low-Resource Languages}
Malayalam, with 38 million speakers, remains underserved in NLP research. Advanced neural machine translation models provide Malayalam $\leftrightarrow$ English translation capabilities. Unicode block detection (\texttt{U+0D00}--\texttt{U+0D7F}) enables automatic language identification without external classifiers, simplifying the bilingual processing pipeline. Recent advances in multilingual transformer models have improved translation quality for Indian languages, though challenges remain in handling technical legal terminology.

\section{High-Throughput Model Serving with FastAPI}
Modern RAG deployment requires microservices capable of handling concurrent requests with minimal latency. FastAPI, powered by Starlette and Uvicorn, provides asynchronous request handling with Pydantic schema validation. Unlike traditional WSGI servers, ASGI frameworks enable non-blocking I/O operations critical for external API calls to vector databases and LLM providers.

\section{Civic Technology, Ethics, and Legal Accessibility}
Digital platforms for civic engagement face unique challenges in balancing accessibility with accuracy. Gupta et~al.~\cite{ref_gupta2024_legalassist} designed Legal Assist AI using transformer models to democratize legal counsel for Indian citizens lacking access to professional advice. At the same time, Webb~\cite{ref_webb2026_genai} emphasized the epistemic and ethical dimensions of deploying generative AI in legal practice, arguing that automated systems must preserve competence, transparency, and ethical compliance to prevent the diminution of practical legal wisdom.

\section{Summary of Identified Research Gaps}
Table~\ref{tab:lit_comparison} summarizes the comparative analysis between existing legal information retrieval approaches and the proposed system.

\vspace{2pt}
\begin{table}[H]
\centering
\begin{singlespace}
\caption{Comparative analysis of legal information retrieval systems.}
\label{tab:lit_comparison}
\footnotesize
\renewcommand{\arraystretch}{1.1}
\begin{tabularx}{\textwidth}{>{\raggedright\arraybackslash}p{2.8cm} >{\centering\arraybackslash}p{1.8cm} >{\centering\arraybackslash}p{1.5cm} >{\centering\arraybackslash}p{1.8cm} >{\raggedright\arraybackslash}X}
\toprule
\textbf{System} & \textbf{Retrieval Method} & \textbf{Bilingual} & \textbf{Source Citation} & \textbf{Limitations} \\
\midrule
Google Search & Keyword BM25 & No & No & Returns raw PDFs without context; no answer generation. \\
Legal GPT-4 & Parametric only & Yes & No & Hallucinates statutes; no grounding in actual legal text. \\
Manupatra/SCC & Keyword + Boolean & No & Yes & Case law focused; weak on statutes; English only. \\
LIMBS (Govt) & Metadata catalog & No & No & Document listing only; no semantic search or answers. \\
\textbf{CiviQ (This Project)} & \textbf{Dense vector RAG} & \textbf{Yes (Malayalam)} & \textbf{Yes} & \textbf{Grounded answers with citations; bilingual Malayalam/English.} \\
\bottomrule
\end{tabularx}
\end{singlespace}
\end{table}
\vspace{-4pt}



%==============================================================================
% CHAPTER 3: PROBLEM DEFINITION
%==============================================================================
\chapter{Problem Definition}

A rigorous engineering formulation is essential to transform legal information access challenges into a verifiable, mathematically bounded system. This chapter formalizes the problem statement, describes the technical and accessibility challenges, delineates operational boundaries, and presents the mathematical formulation governing RAG retrieval and bilingual query processing.

\section{Problem Statement}
Indian citizens require timely, comprehensible legal information to navigate civic complaint systems, consumer protection mechanisms, and information access rights. However, existing statutory repositories provide raw legal text without contextual interpretation, and most resources exclude non-English speakers. Generative AI systems hallucinate non-existent legal provisions without source grounding.

The objective of this project is to design, implement, and evaluate an automated Retrieval-Augmented Generation (RAG) system for grounded legal information retrieval, utilizing \textbf{sentence-transformers embedding (all-MiniLM-L6-v2)}, \textbf{Pinecone vector indexing}, \textbf{Google Gemini 2.5 Flash generation}, bilingual Malayalam $\leftrightarrow$ English query processing, a \textbf{FastAPI backend}, and a \textbf{React 19 frontend} with voice input.

\section{Problem Description}
The core limitations confronting current legal information access workflows arise from four technical and social challenges:
\begin{enumerate}
    \item \textbf{Semantic Search Gap:} Keyword-based Boolean search fails for natural language queries. A citizen asking ``What are my rights if defective goods cause injury?'' receives no results without knowing to search for ``Consumer Protection Act 2019 Section 2(10) product liability.''
    \item \textbf{Hallucination in Parametric LLMs:} Pure GPT-4 or Gemini queries generate plausible-sounding legal provisions, sections, and penalties that do not exist in Indian statutes, creating dangerous misinformation.
    \item \textbf{Linguistic Exclusion:} Malayalam speakers (38 million citizens in Kerala) cannot access legal information in their native language, creating barriers to exercising RTI rights, filing consumer complaints, or understanding criminal procedure.
    \item \textbf{Absence of Source Attribution:} Existing AI legal tools provide answers without citing specific Acts, sections, or subsections, preventing verification and undermining trust.
\end{enumerate}

\section{Scope and Objectives}
The operational boundaries of the proposed system are clearly defined:
\begin{itemize}
    \item \textbf{In-Scope Capabilities:}
    \begin{itemize}
        \item Automated ingestion and preprocessing of PDF legal documents (BNS, BNSS, BSA, RTI Act, Consumer Protection Act, Indian Contract Act, Constitution).
        \item Overlapping word-chunking (900-word chunks, 150-word overlap, 40-word minimum) preserving legal context across boundaries.
        \item Sentence-transformer embedding (all-MiniLM-L6-v2, 384-dim, L2-normalized) and Pinecone batch upload.
        \item Semantic retrieval of top-k passages (default k=5, user-configurable 1–10) via cosine similarity.
        \item Gemini 2.5 Flash grounded answer generation with explicit source citations (Act name, section, clause).
        \item Automatic Malayalam detection (Unicode U+0D00–U+0D7F) and bidirectional translation.
        \item FastAPI backend with endpoints: GET /, GET /api/health, POST /ask, POST /analyze-document.
        \item React 19 + Vite 8 frontend with voice input (Web Speech API), PDF upload, language selection, top-k slider, source citation cards.
    \end{itemize}
    \item \textbf{Out-of-Scope Constraints:}
    \begin{itemize}
        \item The system is explicitly designed for informational and educational purposes and is not legal advice. Prominent disclaimers are displayed.
        \item Judicial case law and precedents are excluded from the current corpus (future enhancement).
        \item Image-only PDFs (scanned documents without OCR) are not supported.
    \end{itemize}
\end{itemize}

\section{Theoretical Problem Formulation}
Let $\mathcal{C} = \{d_1, d_2, \dots, d_N\}$ represent the legal document corpus. Each document $d_i$ is partitioned into overlapping chunks $\{c_{i,1}, c_{i,2}, \dots, c_{i,m_i}\}$ where each chunk $c_{i,j}$ contains 900 words with 150-word overlap with adjacent chunks.

Let $\phi: \mathcal{T} \to \mathbb{R}^{384}$ denote the embedding function (all-MiniLM-L6-v2), mapping text $t \in \mathcal{T}$ to a 384-dimensional dense vector. Embeddings are L2-normalized:
\begin{equation}
    \mathbf{e}(t) = \frac{\phi(t)}{\|\phi(t)\|_2}
\end{equation}

For a user query $q$, the retrieval function $R$ computes cosine similarity with all chunk embeddings:
\begin{equation}
    \text{sim}(q, c) = \mathbf{e}(q)^\top \mathbf{e}(c)
\end{equation}
and returns the top-k chunks:
\begin{equation}
    R(q, k) = \text{argmax}_{S \subset \mathcal{C}, |S|=k} \sum_{c \in S} \text{sim}(q, c)
\end{equation}

The generation function $G$ conditions a large language model on retrieved context:
\begin{equation}
    \text{answer} = G(q \,|\, R(q, k))
\end{equation}
producing a grounded response with explicit source citations.

For Malayalam queries, a detection function $\delta$ identifies Unicode block U+0D00–U+0D7F, and a translation function $T$ converts queries:
\begin{equation}
    q_{\text{en}} = T(q_{\text{ml}}) \quad \text{and} \quad \text{answer}_{\text{ml}} = T^{-1}(\text{answer}_{\text{en}})
\end{equation}
ensuring bidirectional language support without modifying the core retrieval pipeline.


%==============================================================================
% CHAPTER 4: REQUIREMENTS SPECIFICATION
%==============================================================================
\chapter{Requirements Specification}

The formal System Requirements Specification (SRS) establishes the operational, technical, and regulatory boundaries governing the design, implementation, and deployment of the RAG-based legal information retrieval platform.

\section{Functional Requirements}
The functional requirements define the explicit capabilities and behaviors implemented by the system:
\begin{enumerate}[(i)]
    \item \textbf{FR-1: Text and Voice Query Input:} The system shall accept typed queries and browser-based speech-to-text voice input in English and Malayalam.
    \item \textbf{FR-2: Automatic Language Detection:} The system shall detect Malayalam queries via Unicode block analysis (U+0D00–U+0D7F) and route them through translation.
    \item \textbf{FR-3: PDF Document Upload and Analysis:} The system shall accept PDF uploads ($\le 50$~MB), extract text via PyMuPDF, and provide document summary or question-based analysis.
    \item \textbf{FR-4: Semantic Retrieval with Top-K Selection:} The system shall embed queries using all-MiniLM-L6-v2, retrieve top-k relevant passages from Pinecone (user-configurable k=1–10, default k=5), and display sources.
    \item \textbf{FR-5: Grounded Answer Generation:} The system shall generate contextual answers via Gemini 2.5 Flash conditioned on retrieved passages, with explicit citation of source documents and sections.
    \item \textbf{FR-6: Bilingual Response Translation:} For Malayalam queries, the system shall translate generated English answers back to Malayalam before display.
    \item \textbf{FR-7: Health Check and Status Monitoring:} The backend shall expose \texttt{GET /api/health} returning database connectivity status and system readiness.
\end{enumerate}

\section{Non-Functional Requirements}
Non-functional requirements specify the operational quality attributes, performance bounds, and safety standards:
\begin{itemize}
    \item \textbf{NFR-1 (Latency and Throughput):} End-to-end query answering shall execute in under 2 seconds including retrieval, generation, and translation, ensuring seamless user interaction.
    \item \textbf{NFR-2 (Reliability and Fault Tolerance):} Invalid PDF uploads or malformed queries shall be gracefully rejected with structured error responses, without terminating server execution.
    \item \textbf{NFR-3 (Scalability):} The system shall handle concurrent requests via asynchronous FastAPI processing, supporting multiple simultaneous users.
    \item \textbf{NFR-4 (Legal Disclaimer Boundaries):} All user interfaces shall prominently display disclaimers clarifying that outputs are informational only and do not constitute legal advice.
    \item \textbf{NFR-5 (Cross-Platform Accessibility):} The web interface shall maintain responsive layouts across desktop, tablet, and mobile screens with accessible font sizes and color contrast.
\end{itemize}

\section{System Requirements}

\subsection{Hardware and Software Requirements}
\begin{itemize}
    \item \textbf{Hardware Requirements:}
    \begin{itemize}
        \item CPU: 4-core modern processor (Intel i5 / AMD Ryzen 5 or better).
        \item RAM: 8 GB minimum (16 GB recommended).
        \item Storage: At least 20 GB of free disk space for source documents, dependencies, and generated artifacts.
        \item Internet: Stable connection for Pinecone API and Gemini API calls.
    \end{itemize}
    \item \textbf{Software Requirements:}
    \begin{itemize}
        \item Operating System: Windows 11, Ubuntu 22.04 or later, or macOS 12+.
        \item Python 3.12+ runtime.
        \item Node.js 18+ and npm for frontend build and development.
        \item Cloud VM or containerized environment (Docker) for production deployment.
    \end{itemize}
\end{itemize}

\subsection{External Services and APIs}
\begin{itemize}
    \item \textbf{Vector Database:} Pinecone managed service (free tier: 1M vectors, 100K queries/month).
    \item \textbf{LLM API:} Google Gemini API 2.5 Flash (free tier: 15 RPM, 1M TPM).
    \item \textbf{Translation Service:} Google Gemini API for bidirectional Malayalam-English translation.
\end{itemize}

\subsection{Databases and Storage}
\begin{itemize}
    \item \textbf{Pinecone Vector Database:} Managed vector store with HNSW indexing, 384-dimensional embeddings, cosine similarity metric, and batch upsert capabilities.
    \item \textbf{Legal Corpus Repository:} Structured PDF document collection including BNS 2023, BNSS 2023, BSA 2023, Consumer Protection Act 2019, Indian Contract Act 1872, RTI Act 2005, and the Constitution of India.
    \item \textbf{Query and Retrieval Logs:} Structured telemetry recording user queries, detected language, retrieved passages, similarity scores, and response latency.
\end{itemize}

\subsection{Libraries and Frameworks}
\begin{itemize}
    \item \textbf{Backend Core:} \texttt{python} ($\ge 3.12$), \texttt{fastapi} ($\ge 0.100$), \texttt{uvicorn}, \texttt{pydantic} ($\ge 2.0$), \texttt{pymupdf} (fitz).
    \item \textbf{Embedding and Retrieval:} \texttt{sentence-transformers} (all-MiniLM-L6-v2), \texttt{pinecone-client}, \texttt{numpy}, \texttt{google-generativeai}.
    \item \textbf{Frontend Client:} \texttt{React 19}, \texttt{Vite 8}, \texttt{Tailwind CSS 4}, \texttt{React Markdown}, browser Web Speech API.
\end{itemize}

\subsection{Front-End Interaction Layer}
The user interface is engineered as an accessible, responsive single-page application built with React 19, Vite 8, and Tailwind CSS 4, optimized for citizen engagement with Indian statutory information. The interaction layer comprises five core functional modules:

\begin{itemize}
    \item \textbf{Conversational Legal Workspace:} Multi-turn chat pane with smooth auto-scrolling, streaming generation indicators, response latency telemetry ($\approx 1.4$\,s), and GitHub-flavored Markdown rendering supporting formatted statutory tables, bold legislative section headings, and procedural steps.
    \item \textbf{Multimodal Speech Dictation:} In-browser Web Speech API integration enabling voice queries in Indian English (\texttt{en-IN}) and Malayalam (\texttt{ml-IN}), featuring real-time visual recording feedback and automatic fallback to typed input.
    \item \textbf{Bilingual Routing \& Language Toggles:} Three-state selector (\texttt{Auto-Detect}, \texttt{English}, \texttt{Malayalam}) with client-side Unicode scanning (\texttt{U+0D00}--\texttt{U+0D7F}) for automatic translation orchestration via Google Gemini.
    \item \textbf{Dynamic Retrieval Depth Slider:} Configurable slider allowing users to calibrate semantic retrieval depth from $k = 1$ to $k = 10$ passages (default $k = 5$), trading off citation breadth against response latency.
    \item \textbf{Citation Cards \& Document Ingestion:} Modular citation cards displaying statute titles (BNS 2023, RTI Act 2005), section numbers, and similarity scores with collapsible drawers to inspect retrieved excerpts, alongside a drag-and-drop PDF upload zone ($\le 50$\,MB) and persistent legal disclaimers.
\end{itemize}

The interaction layer is designed to keep the user informed throughout the complete query lifecycle. When a query is submitted, the interface immediately displays the selected language, input modality, retrieval depth, and processing status. During retrieval and generation, progress indicators communicate whether the system is detecting the language, searching the legal corpus, or generating the grounded response. This reduces uncertainty during external API calls and gives users a clear understanding of the system state.

Each generated response is presented together with its supporting statutory evidence. Citation cards identify the source Act, section number, document title, and relevance score, allowing users to inspect the retrieved passage before relying on the explanation. The interface also separates the generated interpretation from the source excerpts so that the distinction between an AI-generated summary and authoritative legal text remains clear.

The front-end applies input validation before requests are sent to the backend. Empty queries are rejected with an explanatory message, unsupported files are identified before upload, and files exceeding the 50\,MB limit are blocked on the client side. Network failures, API rate limits, and low-confidence retrieval results are reported through readable status messages rather than technical stack traces. A persistent disclaimer reminds users that CiviQ provides preliminary legal information and does not replace advice from a qualified legal professional.

The responsive layout supports desktop and tablet displays while maintaining usable controls on smaller screens. Semantic labels, keyboard-accessible buttons, sufficient colour contrast, and readable typography improve accessibility for users with different levels of technical experience. Together, these design decisions make the system transparent, safe, and approachable for citizens seeking preliminary information about Indian laws and civic procedures.


%==============================================================================
% CHAPTER 5: SYSTEM DESIGN
%==============================================================================
\chapter{System Design}

This chapter details the architectural decomposition, component interactions, workflows, algorithms, and data schemas governing the RAG-based legal information retrieval system.

\section{High-Level Architecture}
The system is architected as an asynchronous, decoupled multi-pipeline platform spanning client presentation, API gateway orchestration, core Retrieval-Augmented Generation (RAG) processing, and cloud vector storage, as formalized in Figure~\ref{fig:system_arch_detailed}.

\begin{figure}[htbp]
    \centering
    \adjustbox{max width=0.98\linewidth, keepaspectratio, center}{\documentclass[tikz,border=12pt]{standalone}
\usepackage{tikz}
\usepackage{fontawesome5}
\usepackage{amsmath,amssymb}
\usetikzlibrary{arrows.meta,positioning,fit,backgrounds,calc,shadows,shapes.geometric}

% Curated Academic Palette (Deep Navy, Slate, Steel, Teal, Muted Gold)
\definecolor{prismNavy}{RGB}{15, 44, 89}
\definecolor{prismBlue}{RGB}{29, 78, 158}
\definecolor{prismTeal}{RGB}{13, 148, 136}
\definecolor{prismGold}{RGB}{202, 138, 4}
\definecolor{prismDark}{RGB}{15, 23, 42}
\definecolor{prismSlate}{RGB}{71, 85, 105}
\definecolor{statusFail}{RGB}{185, 28, 28}
\definecolor{cInput}{RGB}{248, 250, 252}
\definecolor{cStem}{RGB}{240, 249, 255}
\definecolor{cStage1}{RGB}{240, 253, 250}
\definecolor{cStage2}{RGB}{254, 252, 232}
\definecolor{cStage3}{RGB}{254, 242, 242}
\definecolor{cStage4}{RGB}{238, 242, 255}
\definecolor{cHead}{RGB}{224, 231, 255}
\definecolor{cOutput}{RGB}{240, 253, 244}
\definecolor{lavenderLight}{RGB}{241, 245, 249}

\begin{document}
\begin{tikzpicture}[
  x=1cm, y=1cm,
  every node/.style={align=center},
  group/.style={draw=prismNavy!60, dashed, line width=0.85pt,
    fill=lavenderLight!40, rounded corners=6pt},
  grouptitle/.style={fill=prismNavy, text=white, font=\fontsize{6.2}{7.4}\selectfont\bfseries,
    rounded corners=2.5pt, inner xsep=4.5pt, inner ysep=2.8pt},
  colbox/.style={rectangle, rounded corners=4pt, line width=0.8pt,
    minimum height=1.20cm, inner ysep=4.0pt, inner xsep=4.0pt,
    font=\fontsize{6.6}{7.8}\selectfont\bfseries,
    drop shadow={shadow xshift=0.8pt, shadow yshift=-0.8pt, opacity=0.10}},
  arrow/.style={
    -{Stealth[scale=0.85]}, line width=0.85pt, draw=prismNavy!85
  },
  busline/.style={
    -{Stealth[scale=0.9]}, line width=1.1pt, draw=prismBlue
  },
  feedback/.style={
    -{Stealth[scale=0.9]}, line width=1.1pt, dashed, draw=prismTeal
  }
]

  %============================================================================
  % COLUMN 1: PRESENTATION & CLIENT INTERFACE TIER
  %============================================================================
  \node[group, minimum width=3.90cm, minimum height=11.00cm] (g1) at (2.05, 5.70) {};
  \node[grouptitle] at (2.05, 11.20) {Tier 1: Client Presentation};

  \node[colbox, text width=3.55cm, fill=cInput, draw=prismSlate!70] (g1_user) at (2.05, 9.80) {%
    {\fontsize{6.1}{7.2}\selectfont\color{prismDark}\faUser\ \textbf{Citizen Ingestion}}\\[2.0pt]
    {\fontsize{5.8}{6.8}\selectfont\normalfont Natural Language Ingestion\\[1pt]Typed Text / Web Speech Audio}%
  };

  \node[colbox, text width=3.55cm, fill=cStem, draw=prismBlue!70] (g1_web) at (2.05, 7.30) {%
    {\fontsize{6.1}{7.2}\selectfont\color{prismNavy}\faDesktop\ \textbf{React 19 + Vite 8 UI}}\\[2.0pt]
    {\fontsize{5.8}{6.8}\selectfont\normalfont Tailwind CSS 4 $\cdot$ Markdown Parser\\[1pt]Real-time Speech Recognition $\cdot$ Upload}%
  };

  \node[colbox, text width=3.55cm, fill=cStage4, draw=prismNavy!70] (g1_ctrl) at (2.05, 4.80) {%
    {\fontsize{6.1}{7.2}\selectfont\color{prismNavy}\faSlidersH\ \textbf{Interactive Control Deck}}\\[2.0pt]
    {\fontsize{5.8}{6.8}\selectfont\normalfont Language Toggle (\texttt{auto}, \texttt{en}, \texttt{ml})\\[1pt]Depth Slider ($k \in [1, 10]$)}%
  };

  \node[colbox, text width=3.55cm, fill=cStage1, draw=prismTeal!80] (g1_disclaimer) at (2.05, 2.30) {%
    {\fontsize{6.1}{7.2}\selectfont\color{prismTeal!90!black}\faBalanceScale\ \textbf{Statutory Safety Guard}}\\[2.0pt]
    {\fontsize{5.8}{6.8}\selectfont\normalfont Informational Notice Enforcement\\[1pt]Strict Non-Legal Advice Disclaimer}%
  };

  \draw[arrow] (g1_user.south) -- (g1_web.north);
  \draw[arrow] (g1_web.south) -- (g1_ctrl.north);
  \draw[arrow] (g1_ctrl.south) -- (g1_disclaimer.north);

  %============================================================================
  % COLUMN 2: API GATEWAY & ORCHESTRATION TIER
  %============================================================================
  \node[group, minimum width=3.90cm, minimum height=11.00cm] (g2) at (6.45, 5.70) {};
  \node[grouptitle] at (6.45, 11.20) {Tier 2: API Gateway \& Router};

  \node[colbox, text width=3.55cm, fill=cStage2, draw=prismGold!90] (g2_fastapi) at (6.45, 9.80) {%
    {\fontsize{6.1}{7.2}\selectfont\color{prismGold!90!black}\faServer\ \textbf{FastAPI REST Gateway}}\\[2.0pt]
    {\fontsize{5.8}{6.8}\selectfont\normalfont Asynchronous ASGI (Uvicorn)\\[1pt]\texttt{/ask} $\cdot$ \texttt{/analyze-document} $\cdot$ \texttt{/api/health}}%
  };

  \node[colbox, text width=3.55cm, fill=cInput, draw=prismSlate!70] (g2_validate) at (6.45, 7.30) {%
    {\fontsize{6.1}{7.2}\selectfont\color{prismDark}\faCheckCircle\ \textbf{Payload Audit \& Guard}}\\[2.0pt]
    {\fontsize{5.8}{6.8}\selectfont\normalfont MIME Verification: \texttt{application/pdf}\\[1pt]Boundary: $\le 50$\,MB $\cdot$ Pydantic v2 Schema}%
  };

  \node[colbox, text width=3.55cm, fill=cStage1, draw=prismTeal!80] (g2_lang) at (6.45, 4.80) {%
    {\fontsize{6.1}{7.2}\selectfont\color{prismTeal!90!black}\faLanguage\ \textbf{Linguistic Router}}\\[2.0pt]
    {\fontsize{5.8}{6.8}\selectfont\normalfont Unicode Script Scanner\\[1pt]Malayalam Range: U+0D00 to U+0D7F}%
  };

  \node[colbox, text width=3.55cm, fill=cStage3, draw=statusFail!70] (g2_error) at (6.45, 2.30) {%
    {\fontsize{6.1}{7.2}\selectfont\color{statusFail}\faExclamationTriangle\ \textbf{Resilience Controller}}\\[2.0pt]
    {\fontsize{5.8}{6.8}\selectfont\normalfont 400 Bad Query $\cdot$ 413 File Limit\\[1pt]415 Format $\cdot$ 503 LLM Backoff Retries}%
  };

  \draw[arrow] (g2_fastapi.south) -- (g2_validate.north);
  \draw[arrow] (g2_validate.south) -- (g2_lang.north);
  \draw[arrow] (g2_lang.south) -- (g2_error.north);

  %============================================================================
  % COLUMN 3: RAG PROCESSING & RETRIEVAL TIER
  %============================================================================
  \node[group, minimum width=4.15cm, minimum height=11.00cm] (g3) at (10.95, 5.70) {};
  \node[grouptitle] at (10.95, 11.20) {Tier 3: RAG Pipeline \& LLM};

  \node[colbox, text width=3.80cm, fill=cStem, draw=prismNavy!80] (g3_chunk) at (10.95, 9.80) {%
    {\fontsize{5.9}{7.0}\selectfont\color{prismNavy}\faFilePdf\ \textbf{Corpus Ingestion \& Chunker}}\\[2.0pt]
    {\fontsize{5.8}{6.8}\selectfont\normalfont PyMuPDF Stream Clean $\cdot$ Regex Strip\\[1pt]Overlapping Word Chunks ($c=900, o=150$)}%
  };

  \node[colbox, text width=3.80cm, fill=cHead, draw=prismBlue!90] (g3_embed) at (10.95, 7.30) {%
    {\fontsize{5.9}{7.0}\selectfont\color{prismNavy}\faBrain\ \textbf{Sentence Transformer Embedder}}\\[2.0pt]
    {\fontsize{5.8}{6.8}\selectfont\normalfont \texttt{all-MiniLM-L6-v2} (384-d Dense)\\[1pt]Unit $L_2$ Normalization ($\|\mathbf{e}\|_2 = 1.0$)}%
  };

  \node[colbox, text width=3.80cm, fill=cStage2, draw=prismGold!90] (g3_retriever) at (10.95, 4.80) {%
    {\fontsize{5.9}{7.0}\selectfont\color{prismGold!90!black}\faSearch\ \textbf{Top-K Semantic Retriever}}\\[2.0pt]
    {\fontsize{5.8}{6.8}\selectfont\normalfont Cosine ANN Search over Pinecone Index\\[1pt]Passage Ranking $\cdot$ Metadata Attribution}%
  };

  \node[colbox, text width=3.80cm, fill=cOutput, draw=prismTeal!80] (g3_gemini) at (10.95, 2.30) {%
    {\fontsize{5.9}{7.0}\selectfont\color{prismTeal!90!black}\faRobot\ \textbf{Google Gemini 2.5 Flash}}\\[2.0pt]
    {\fontsize{5.8}{6.8}\selectfont\normalfont Grounded Evidence QA $\cdot$ Statutory Cites\\[1pt]Bidirectional Translation (Malayalam $\leftrightarrow$ English)}%
  };

  \draw[arrow] (g3_chunk.south) -- (g3_embed.north);
  \draw[arrow] (g3_embed.south) -- (g3_retriever.north);
  \draw[arrow] (g3_retriever.south) -- (g3_gemini.north);

  %============================================================================
  % COLUMN 4: STORAGE & KNOWLEDGE TIER
  %============================================================================
  \node[group, minimum width=3.90cm, minimum height=11.00cm] (g4) at (15.45, 5.70) {};
  \node[grouptitle] at (15.45, 11.20) {Tier 4: Vector Storage \& DB};

  \node[colbox, text width=3.55cm, fill=cStage1, draw=prismTeal!80] (g4_corpus) at (15.45, 9.80) {%
    {\fontsize{6.3}{7.5}\selectfont\color{prismTeal!90!black}\faBook\ \textbf{Indian Statutory Corpus}}\\[2.0pt]
    {\fontsize{5.8}{6.8}\selectfont\normalfont BNS, BNSS, BSA, RTI Act, CPA 2019,\\[1pt]Indian Contract Act, Constitution}%
  };

  \node[colbox, text width=3.55cm, fill=cStage4, draw=prismBlue!70] (g4_cache) at (15.45, 7.30) {%
    {\fontsize{6.3}{7.5}\selectfont\color{prismNavy}\faFileContract\ \textbf{Ephemeral Document Buffer}}\\[2.0pt]
    {\fontsize{5.8}{6.8}\selectfont\normalfont Session In-Memory PDF Stream\\[1pt]Ad-hoc Upload Summaries \& QA}%
  };

  \node[colbox, text width=3.55cm, fill=cInput, draw=prismSlate!70] (g4_pinecone) at (15.45, 4.80) {%
    {\fontsize{6.3}{7.5}\selectfont\color{prismDark}\faDatabase\ \textbf{Pinecone Vector Store}}\\[2.0pt]
    {\fontsize{5.8}{6.8}\selectfont\normalfont Managed Cloud Index (384-d HNSW)\\[1pt]Batch Vector Upsert ($B=100$ chunks)}%
  };

  \node[colbox, text width=3.55cm, fill=cStem, draw=prismNavy!70] (g4_telemetry) at (15.45, 2.30) {%
    {\fontsize{6.3}{7.5}\selectfont\color{prismNavy}\faChartLine\ \textbf{Telemetry \& Query Logs}}\\[2.0pt]
    {\fontsize{5.8}{6.8}\selectfont\normalfont Query History $\cdot$ Detected Language\\[1pt]Retrieved Cites $\cdot$ Latency ($\approx 1.42$\,s)}%
  };

  \draw[arrow] (g4_corpus.south) -- (g4_cache.north);
  \draw[arrow] (g4_cache.south) -- (g4_pinecone.north);
  \draw[arrow] (g4_pinecone.south) -- (g4_telemetry.north);

  %============================================================================
  % INTER-TIER COMMUNICATION CHANNELS (Orthogonal, Zero Overlap across Nodes)
  %============================================================================

  % 1. Client to API Gateway: Direct horizontal link between row 1 nodes
  \draw[busline] (g1_user.east) -- (g2_fastapi.west);
  \node[font=\fontsize{5.2}{6.2}\selectfont\bfseries\color{prismNavy}, fill=white, inner sep=1.2pt]
    at (4.25, 10.15) {HTTP \texttt{/ask}};

  % 2. Gateway to RAG: Step from Language Router (y=4.8) up to Embedder (y=7.3)
  \draw[busline] (g2_lang.east) -- (8.65, 4.80) -- (8.65, 7.30) -- (g3_embed.west);
  \node[font=\fontsize{5.0}{6.0}\selectfont\bfseries\color{prismNavy}, fill=white, inner sep=1.2pt]
    at (8.65, 6.05) {Norm.\\$q_{\mathrm{en}}$};

  % 3. Offline Corpus Ingestion: Row 1 horizontal link
  \draw[arrow, dashed, draw=prismSlate] (g4_corpus.west) -- (g3_chunk.east);
  \node[font=\fontsize{5.2}{6.2}\selectfont\bfseries\color{prismSlate}, fill=white, inner sep=1.2pt]
    at (13.20, 10.15) {PDF Ingest};

  % 4. Online Top-k Retrieval: Direct horizontal link at Row 3 (y=4.8)
  \draw[busline] (g3_retriever.east) -- (g4_pinecone.west);
  \node[font=\fontsize{5.2}{6.2}\selectfont\bfseries\color{prismBlue}, fill=white, inner sep=1.2pt]
    at (13.20, 5.15) {Cosine Top-$k$};

  % 5. Grounded Response Delivery (Feedback Bus):
  % Route safely below all tiers and up along left margin into React Web Client
  \draw[feedback] (g3_gemini.south) -- (10.95, 0.70) -- (-0.60, 0.70)
    node[midway, above=2pt, fill=white, draw=prismTeal, rounded corners=2pt, inner sep=2.5pt, font=\fontsize{5.8}{6.8}\selectfont\bfseries\color{prismTeal}]
    {Grounded Legal Response + Statutory Citations + Metadata Cards}
    -- (-0.60, 7.30) -- (g1_web.west);

\end{tikzpicture}
\end{document}
}
    \caption{Multi-tier RAG system architecture showing Client Presentation, API Gateway \& Validation, RAG Processing \& LLM Inference, and Knowledge Storage \& Observability tiers.}
    \label{fig:system_arch_detailed}
\end{figure}

The architecture decomposes the platform into four discrete operational tiers:

\subsection{Client Presentation Tier (Tier 1)}
The client interface is engineered as an accessible, responsive single-page application built on React~19, Vite~8, and Tailwind CSS~4. It provides citizens, students, and legal researchers with an intuitive workspace for querying statutory corpora. Key sub-components include:
\begin{itemize}
    \item \textbf{Multimodal Query Ingestion:} Accepts typed textual prompts alongside real-time browser-based voice dictation via the HTML5 Web Speech API, with acoustic models configured for Indian English (\texttt{en-IN}) and Malayalam (\texttt{ml-IN}).
    \item \textbf{Interactive Control Deck:} Features a three-state linguistic mode selector (\texttt{Auto-Detect}, \texttt{English}, \texttt{Malayalam}) and an interactive retrieval depth slider ($k \in [1, 10]$, defaulting to $k = 5$) that dynamically configures the breadth of statutory evidence fetched.
    \item \textbf{Citation Rendering Engine:} Renders statutory provisions using GitHub-flavored Markdown, presenting verified source citations as modular, expandable legal cards indicating statute titles, section numbers, and cosine relevance scores.
    \item \textbf{Statutory Safety Guard:} Prominently displays persistent, un-dismissible legal notices emphasizing that retrieved outputs represent educational statutory summaries rather than certified advocate advice.
\end{itemize}

\subsection{API Gateway and Validation Tier (Tier 2)}
The backend gateway is implemented in Python using FastAPI running atop the Uvicorn Asynchronous Server Gateway Interface (ASGI) runtime. It serves as the primary ingress controller and orchestrator:
\begin{itemize}
    \item \textbf{Endpoint Specifications:} Exposes RESTful endpoints including \texttt{POST /ask} for conversational statutory inference, \texttt{POST /analyze-document} for user-uploaded PDF processing, and \texttt{GET /api/health} for readiness and database liveness checks.
    \item \textbf{Payload Audit and Boundary Enforcement:} Validates incoming requests using Pydantic~v2 schema models. PDF uploads are inspected via MIME sniffing (\texttt{application/pdf}) and bound to a strict $50$\,MB size limit to prevent memory exhaustion attacks.
    \item \textbf{Linguistic Inspection Router:} Scans the raw input string for Unicode codepoints falling within the Malayalam character block ($\mathrm{U{+}0D00}$ to $\mathrm{U{+}0D7F}$). Upon detection, the query is dispatched to the translation module prior to vector encoding.
    \item \textbf{Resilience and Exception Engine:} Handles external service failures with exponential backoff and jitter, transforming service errors into structured HTTP error responses (HTTP~400 Bad Query, HTTP~413 File Limit, HTTP~415 Format Unsupported, and HTTP~503 Service Unavailable).
\end{itemize}

\subsection{Core RAG Processing and LLM Inference Tier (Tier 3)}
The RAG processing tier constitutes the analytical heart of the system, bridging raw textual documents and dense semantic representations:
\begin{itemize}
    \item \textbf{PyMuPDF Ingestion and Chunking Engine:} Streams PDF documents through PyMuPDF (\texttt{fitz}), cleans running headers, footers, and gazette watermarks, and executes overlapping word-window chunking ($c = 900$ words, stride $s = 750$, overlap $o = 150$, and fragment threshold $\theta_{\min} \ge 40$ words).
    \item \textbf{Sentence-Transformer Dense Embedder:} Maps statutory chunks and incoming user queries into a shared 384-dimensional dense semantic vector space using \texttt{all-MiniLM-L6-v2}, applying unit $L_2$ normalization ($\|\mathbf{e}\|_2 = 1.0$) to enable direct inner-product cosine evaluation.
    \item \textbf{Top-K Semantic Retriever:} Traverses the vector index to identify the $k$ nearest statutory passages matching the query vector, extracting metadata attributes including Act name, section number, and verified textual excerpts.
    \item \textbf{Google Gemini 2.5 Flash Synthesis Engine:} Constructs a constrained, zero-hallucination context prompt injecting retrieved passages into the Google Gemini 2.5 Flash model ($T = 0.3$), enforcing factual attribution to statutory clauses and translating generated answers into Malayalam when requested.
\end{itemize}

\subsection{Knowledge Storage and Observability Tier (Tier 4)}
The storage tier decouples persistent vector indexing from transactional telemetry logging:
\begin{itemize}
    \item \textbf{Pinecone Managed Vector Database:} Hosts the 384-dimensional dense embeddings in a managed cloud index configured with Hierarchical Navigable Small World (HNSW) graph structures and cosine distance metrics, supporting batched vector upserts ($B = 100$).
    \item \textbf{Indian Statutory Corpus Repository:} Maintains authoritative statutory collections encompassing the Bharatiya Nyaya Sanhita (BNS) 2023, Bharatiya Nagarik Suraksha Sanhita (BNSS) 2023, Bharatiya Sakshya Adhiniyam (BSA) 2023, Consumer Protection Act 2019, Indian Contract Act 1872, Right to Information (RTI) Act 2005, and Constitution of India.
    \item \textbf{Query Telemetry and Observability Store:} Captures transactional inference metrics, including user query hashes, language tags, top-$k$ similarity scores, LLM token counts, and end-to-end execution latencies.
    \item \textbf{Ephemeral Document Buffer:} Provides session-scoped in-memory parsing buffers for citizen-uploaded documents evaluated under the \texttt{/analyze-document} pipeline, ensuring uploaded data is never persisted without authorization.
\end{itemize}

\subsection{Latency Budget and Inter-Tier Communication}
The end-to-end operational timeline is governed by strict latency budgets across each processing phase. Table~\ref{tab:latency_budget} details the execution time profile observed across 500 benchmark queries.

\begin{table}[htbp]
\centering
\caption{Empirical Latency Budget Breakdown across Operational Pipeline Tiers.}
\label{tab:latency_budget}
\small
\begin{tabularx}{\textwidth}{l p{3.2cm} >{\centering\arraybackslash}p{2.2cm} X}
\toprule
\textbf{Pipeline Phase} & \textbf{Executing Tier} & \textbf{Mean Latency} & \textbf{Operational Function} \\
\midrule
Audio Transcription & Tier 1 (Client Web) & 15\,ms & In-browser Web Speech API streaming \\
Linguistic Routing & Tier 2 (Gateway) & 2\,ms & Unicode codepoint inspection ($\mathrm{U{+}0D00}$--$\mathrm{U{+}0D7F}$) \\
Malayalam Translation & Tier 3 (Gemini API) & 118\,ms & Zero-shot query translation to English (if needed) \\
Dense Vector Encoding & Tier 3 (Sentence-Trans.) & 38\,ms & 384-d MiniLM CPU forward pass + $L_2$ norm \\
Pinecone HNSW Search & Tier 4 (Vector Store) & 58\,ms & Cosine approximate nearest neighbor lookup ($k=5$) \\
Prompt Assembly & Tier 3 (RAG Controller) & 4\,ms & Context budget packing and instruction tagging \\
Gemini Synthesis & Tier 3 (Gemini 2.5 Flash) & 1,140\,ms & Grounded inference with statutory citations \\
Reverse Translation & Tier 3 (Gemini API) & 125\,ms & Reverse translation into Malayalam (if needed) \\
Payload Delivery \& UI & Tier 1 (React 19) & 32\,ms & JSON parsing, Markdown render, citation cards \\
\midrule
\textbf{Total Mean Latency} & \multicolumn{3}{l}{\textbf{1.42 seconds (English) / 1.66 seconds (Malayalam bilingual)}} \\
\bottomrule
\end{tabularx}
\end{table}


\section{System Workflow and Operational Pipeline}

\subsection{Proposed Methods Pipeline}
The end-to-end operational pipeline executing across the system is depicted in Figure~\ref{fig:framework}. The platform operates across two synchronized methodologies: Phase~A (Offline Statutory Corpus Preparation and HNSW Indexing) and Phase~B (Online Multimodal Query Processing, Semantic Retrieval, and Grounded LLM Generation).

\begin{figure}[htbp]
    \centering
    \adjustbox{max width=0.98\linewidth, keepaspectratio, center}{\documentclass[tikz,border=12pt]{standalone}
\usepackage{tikz}
\usepackage{fontawesome5}
\usepackage{amsmath,amssymb}
\usetikzlibrary{arrows.meta,positioning,fit,backgrounds,calc,shadows,shapes.geometric}

% Professional Academic Color Palette
\definecolor{prismNavy}{RGB}{15, 44, 89}
\definecolor{prismBlue}{RGB}{29, 78, 158}
\definecolor{prismTeal}{RGB}{13, 148, 136}
\definecolor{prismGold}{RGB}{202, 138, 4}
\definecolor{prismDark}{RGB}{15, 23, 42}
\definecolor{prismSlate}{RGB}{71, 85, 105}
\definecolor{cInput}{RGB}{248, 250, 252}
\definecolor{cStem}{RGB}{240, 249, 255}
\definecolor{cStage1}{RGB}{240, 253, 250}
\definecolor{cStage2}{RGB}{254, 252, 232}
\definecolor{cStage3}{RGB}{254, 242, 242}
\definecolor{cStage4}{RGB}{238, 242, 255}
\definecolor{cHead}{RGB}{224, 231, 255}
\definecolor{cOutput}{RGB}{240, 253, 244}
\definecolor{lavenderLight}{RGB}{241, 245, 249}

\begin{document}
\begin{tikzpicture}[
  x=1cm, y=1cm,
  every node/.style={align=center},
  phasebox/.style={
    draw=prismNavy!70, dashed, line width=0.85pt,
    fill=lavenderLight!35, rounded corners=6pt
  },
  phasetitle/.style={
    fill=prismNavy, text=white, font=\fontsize{7.5}{9.0}\selectfont\bfseries,
    rounded corners=2.5pt, inner xsep=8pt, inner ysep=3.2pt
  },
  methodnode/.style={
    rectangle, rounded corners=4pt, line width=0.85pt,
    minimum height=2.85cm, inner ysep=4.0pt, inner xsep=5.0pt,
    font=\fontsize{6.8}{8.0}\selectfont\bfseries,
    drop shadow={shadow xshift=0.8pt, shadow yshift=-0.8pt, opacity=0.10}
  },
  flowarrow/.style={
    -{Stealth[scale=0.95]}, line width=1.1pt, draw=prismNavy
  },
  bridgearrow/.style={
    -{Stealth[scale=1.0]}, line width=1.1pt, dashed, draw=prismTeal
  }
]

  %============================================================================
  % PHASE A: OFFLINE CORPUS PREPARATION & VECTOR INDEXING METHODOLOGY
  %============================================================================
  \node[phasebox, minimum width=18.40cm, minimum height=3.95cm] (phaseA) at (9.20, 6.45) {};
  \node[phasetitle, anchor=west] at (0.25, 8.65) {%
    \faDatabase\ \ Phase A: Offline Corpus Ingestion, Semantic Chunking \& Vector Indexing Methodology%
  };

  % Stage A1: Corpus Ingestion
  \node[methodnode, text width=3.70cm, fill=cInput, draw=prismSlate!80] (m_ingest) at (2.05, 6.25) {%
    {\color{prismNavy}\faFilePdf\ \textbf{1. Document Ingestion}}\\[3.0pt]
    \rule{3.40cm}{0.4pt}\\[3.0pt]
    {\fontsize{6.0}{7.2}\selectfont\normalfont
      $\bullet$ 7 Core Statutes (BNS, BSA, etc.)\\
      $\bullet$ PyMuPDF (\texttt{fitz}) Page Stream Parser\\
      $\bullet$ Strip Headers, Footers, Page Nums\\
      $\bullet$ Unicode Text Sanitization\\
      $\bullet$ Format Integrity Validation
    }%
  };

  % Stage A2: Semantic Overlapping Chunking
  \node[methodnode, text width=3.70cm, fill=cStage1, draw=prismTeal!80] (m_chunk) at (6.65, 6.25) {%
    {\color{prismTeal!90!black}\faStream\ \textbf{2. Chunking Engine}}\\[3.0pt]
    \rule{3.40cm}{0.4pt}\\[3.0pt]
    {\fontsize{6.0}{7.2}\selectfont\normalfont
      $\bullet$ Window Size: $c = 900$ words\\
      $\bullet$ Overlap Stride: $o = 150$ words\\
      $\bullet$ Fragment Filter: $\theta_{\min} \ge 40$ words\\
      $\bullet$ Preserve Legal Section Boundaries\\
      $\bullet$ Corpus Yield: $\approx 7,500$ chunks
    }%
  };

  % Stage A3: Dense Semantic Vectorization
  \node[methodnode, text width=3.70cm, fill=cHead, draw=prismBlue!90] (m_embed) at (11.25, 6.25) {%
    {\color{prismNavy}\faBrain\ \textbf{3. SBERT Embedding}}\\[3.0pt]
    \rule{3.40cm}{0.4pt}\\[3.0pt]
    {\fontsize{6.0}{7.2}\selectfont\normalfont
      $\bullet$ \texttt{all-MiniLM-L6-v2} Encoder\\
      $\bullet$ 384-Dimensional Dense Vectors\\
      $\bullet$ Mean Pooling over Layer-6 Tokens\\
      $\bullet$ Unit $L_2$ Projection: $\mathbf{e} / \|\mathbf{e}\|_2$\\
      $\bullet$ Sub-5ms Vector Inference on CPU
    }%
  };

  % Stage A4: Pinecone Cloud Indexing
  \node[methodnode, text width=3.70cm, fill=cStage4, draw=prismNavy!80] (m_index) at (15.85, 6.25) {%
    {\color{prismNavy}\faServer\ \textbf{4. Vector Indexing}}\\[3.0pt]
    \rule{3.40cm}{0.4pt}\\[3.0pt]
    {\fontsize{6.0}{7.2}\selectfont\normalfont
      $\bullet$ Managed Pinecone Vector Index\\
      $\bullet$ Cosine Metric: $s_i = \mathbf{e}_q^\top \mathbf{e}_{c_i}$\\
      $\bullet$ Batch Upsert Pipeline ($B = 100$)\\
      $\bullet$ Metadata: Act, Section, Page Tags\\
      $\bullet$ Sub-100ms ANN Search Latency
    }%
  };

  \draw[flowarrow] (m_ingest) -- (m_chunk)
    node[midway, above, font=\fontsize{5.5}{6.5}\selectfont\bfseries\color{prismNavy}] {Clean Text};
  \draw[flowarrow] (m_chunk) -- (m_embed)
    node[midway, above, font=\fontsize{5.5}{6.5}\selectfont\bfseries\color{prismNavy}] {Token Chunks};
  \draw[flowarrow] (m_embed) -- (m_index)
    node[midway, above, font=\fontsize{5.5}{6.5}\selectfont\bfseries\color{prismNavy}] {384-d Vectors};

  %============================================================================
  % PHASE B: ONLINE INFERENCE, RETRIEVAL & SYNTHESIS METHODOLOGY
  %============================================================================
  \node[phasebox, minimum width=18.40cm, minimum height=3.95cm] (phaseB) at (9.20, 1.55) {};
  \node[phasetitle, anchor=west] at (0.25, 3.75) {%
    \faComments\ \ Phase B: Online Query Processing, Bilingual Routing, Semantic Retrieval \& Grounded Synthesis%
  };

  % Stage B1: Multimodal Ingestion & Language Routing
  \node[methodnode, text width=3.70cm, fill=cStem, draw=prismBlue!80] (m_query) at (2.05, 1.35) {%
    {\color{prismNavy}\faUser\ \textbf{5. Query \& Routing}}\\[3.0pt]
    \rule{3.40cm}{0.4pt}\\[3.0pt]
    {\fontsize{6.0}{7.2}\selectfont\normalfont
      $\bullet$ Typed Text / Web Speech Audio\\
      $\bullet$ Unicode Scan: U+0D00 to U+0D7F\\
      $\bullet$ Malayalam Detection $\to$ Flag $\tau_{\mathrm{lang}}$\\
      $\bullet$ Gemini Translate: $\mathrm{ml} \to \mathrm{en}$\\
      $\bullet$ Normalization of Query String $q_{\mathrm{en}}$
    }%
  };

  % Stage B2: Real-time Query Vector Projection
  \node[methodnode, text width=3.70cm, fill=cHead, draw=prismBlue!90] (m_qembed) at (6.65, 1.35) {%
    {\color{prismNavy}\faBrain\ \textbf{6. Query Projection}}\\[3.0pt]
    \rule{3.40cm}{0.4pt}\\[3.0pt]
    {\fontsize{6.0}{7.2}\selectfont\normalfont
      $\bullet$ Shared Transformer Encoder $\phi$\\
      $\bullet$ Forward Pass: $\mathbf{e}_q = \phi(q_{\mathrm{en}})$\\
      $\bullet$ Unit Normalization: $\|\mathbf{e}_q\|_2 = 1.0$\\
      $\bullet$ Aligned in Shared Semantic Space\\
      $\bullet$ Low Latency ($< 15$\,ms)
    }%
  };

  % Stage B3: Cosine ANN Top-k Passage Retrieval
  \node[methodnode, text width=3.70cm, fill=cStage2, draw=prismGold!90] (m_ann) at (11.25, 1.35) {%
    {\color{prismGold!90!black}\faSearch\ \textbf{7. Top-K Retrieval}}\\[3.0pt]
    \rule{3.40cm}{0.4pt}\\[3.0pt]
    {\fontsize{6.0}{7.2}\selectfont\normalfont
      $\bullet$ Cosine Similarity: $\mathbf{e}_q^\top \mathbf{e}_{c_i}$\\
      $\bullet$ Pinecone HNSW Graph Traversal\\
      $\bullet$ Fetch Top-$k$ Passages ($k \in [1, 10]$)\\
      $\bullet$ Re-rank by Similarity Score $s_i$\\
      $\bullet$ Extract Statutory Section Labels
    }%
  };

  % Stage B4: Gemini Grounded Synthesis & UI Delivery
  \node[methodnode, text width=3.70cm, fill=cOutput, draw=prismTeal!80] (m_llm) at (15.85, 1.35) {%
    {\color{prismTeal!90!black}\faRobot\ \textbf{8. LLM Synthesis}}\\[3.0pt]
    \rule{3.40cm}{0.4pt}\\[3.0pt]
    {\fontsize{6.0}{7.2}\selectfont\normalfont
      $\bullet$ Strict Context Injection Prompt\\
      $\bullet$ Google Gemini 2.5 Flash ($T=0.3$)\\
      $\bullet$ Citation Grounding \& Verification\\
      $\bullet$ If $\tau_{\mathrm{lang}}=\mathrm{ml}$: Translate Answer\\
      $\bullet$ Deliver JSON Payload $\to$ React 19 UI
    }%
  };

  \draw[flowarrow] (m_query) -- (m_qembed)
    node[midway, above, font=\fontsize{5.5}{6.5}\selectfont\bfseries\color{prismNavy}] {Clean $q_{\mathrm{en}}$};
  \draw[flowarrow] (m_qembed) -- (m_ann)
    node[midway, above, font=\fontsize{5.5}{6.5}\selectfont\bfseries\color{prismNavy}] {Dense $\mathbf{e}_q$};
  \draw[flowarrow] (m_ann) -- (m_llm)
    node[midway, above, font=\fontsize{5.5}{6.5}\selectfont\bfseries\color{prismNavy}] {Top-$k$ Evidence};

  % Cross-Phase Bridge: Pinecone Index (A4) queries into Retrieval (B3)
  % Clean stepped orthogonal route via right channel - ZERO intersection with Phase B title!
  \draw[bridgearrow] (m_index.south) -- ++(0, -0.45) -| (m_ann.north)
    node[pos=0.25, fill=white, inner sep=2pt, font=\fontsize{5.8}{7.0}\selectfont\bfseries\color{prismTeal}]
    {HNSW Cosine Vector Index Matching};

\end{tikzpicture}
\end{document}
}
    \caption{Detailed Methodology Pipeline of the CiviQ legal RAG system, highlighting Phase~A (Offline Corpus Preprocessing \& Vector Indexing) and Phase~B (Online Query Parsing, Semantic Retrieval \& Grounded Synthesis).}
    \label{fig:framework}
\end{figure}

The eight specialized stages executing across the methodology are:
\begin{enumerate}\setlength{\itemsep}{3pt}\setlength{\parskip}{0pt}\setlength{\parsep}{0pt}
    \item \textbf{Stage 1 (Document Ingestion):} High-resolution PDF legal statutes are ingested using PyMuPDF (\texttt{fitz}). The byte stream is decoded, and regular expression filters purge page banners, running headers, and gazette publication markers.
    \item \textbf{Stage 2 (Chunking Engine):} Cleaned statutory text is segmented into fixed-window word chunks ($c = 900$ words) with an overlapping stride ($o = 150$ words). Chunks containing fewer than $\theta_{\min} = 40$ words are discarded to eliminate table of contents fragments.
    \item \textbf{Stage 3 (SBERT Embedding):} Every text chunk is processed through the \texttt{all-MiniLM-L6-v2} transformer model. Mean pooling across Layer-6 token activations produces a 384-dimensional embedding, subsequently projected onto the unit hypersphere via $L_2$ normalization.
    \item \textbf{Stage 4 (Vector Indexing):} Embeddings are bundled into batches of $B = 100$ vectors and uploaded to Pinecone over HTTPS, tagged with relational metadata including statute title, chapter index, section number, and source page ranges.
    \item \textbf{Stage 5 (Query \& Routing):} Citizen queries received via text or voice audio are inspected for Malayalam Unicode codepoints ($\mathrm{U{+}0D00}$ to $\mathrm{U{+}0D7F}$). If Malayalam is detected, a translation sub-pipeline translates the query into formal English legal search terms.
    \item \textbf{Stage 6 (Query Projection):} The standardized English query $q_{\mathrm{en}}$ is transformed into dense vector $\mathbf{e}_q \in \mathbb{R}^{384}$ using the identical sentence-transformer encoder and normalized.
    \item \textbf{Stage 7 (Top-K Retrieval):} Pinecone performs approximate nearest neighbor search over the pre-indexed statutory corpus, scoring chunks by cosine inner product and returning the top-$k$ most relevant statutory passages ($k \in [1, 10]$).
    \item \textbf{Stage 8 (LLM Synthesis):} Retrieved evidence chunks are formatted into an instruction-following prompt injected into Google Gemini 2.5 Flash ($T = 0.3$). The model generates an authoritative, plain-language answer citing statutory sections, which is optionally translated back to Malayalam and delivered to the React~19 frontend.
\end{enumerate}

\subsection{Layered System Workflow}
Figure~\ref{fig:arch_workflow} illustrates the structured execution lifecycle spanning input acquisition to verified answer delivery, emphasizing the conditional bilingual branching and iterative follow-up refinement loop.

\begin{figure}[htbp]
    \centering
    \adjustbox{max width=0.98\linewidth, keepaspectratio, center}{\documentclass[tikz,border=10pt]{standalone}
\usepackage{tikz}
\usepackage{fontawesome5}
\usetikzlibrary{arrows.meta,positioning,fit,backgrounds,calc,shadows,shapes.geometric}

% Color definitions
\definecolor{lavenderDark}{RGB}{15, 44, 89}
\definecolor{lavenderDeep}{RGB}{10, 28, 58}
\definecolor{lavenderMedium}{RGB}{29, 78, 158}
\definecolor{lavenderLight}{RGB}{240, 246, 255}
\definecolor{lavenderBorder}{RGB}{180, 205, 237}
\definecolor{lavenderAccent}{RGB}{2, 132, 199}
\definecolor{prismNavy}{RGB}{15, 44, 89}
\definecolor{prismBlue}{RGB}{29, 78, 158}
\definecolor{prismTeal}{RGB}{13, 148, 136}
\definecolor{prismMint}{RGB}{236, 253, 245}
\definecolor{prismGold}{RGB}{217, 119, 6}
\definecolor{prismDark}{RGB}{15, 23, 42}
\definecolor{prismLight}{RGB}{248, 250, 252}
\definecolor{prismIndigo}{RGB}{30, 58, 138}
\definecolor{prismSlate}{RGB}{100, 116, 139}
\definecolor{statusFail}{RGB}{220, 38, 38}
\definecolor{cInput}{RGB}{241, 245, 249}
\definecolor{cStem}{RGB}{224, 242, 254}
\definecolor{cStage1}{RGB}{209, 250, 229}
\definecolor{cStage2}{RGB}{254, 243, 199}
\definecolor{cStage3}{RGB}{254, 226, 226}
\definecolor{cStage4}{RGB}{207, 250, 254}
\definecolor{cHead}{RGB}{219, 234, 254}
\definecolor{cOutput}{RGB}{255, 228, 230}
\definecolor{cDeliverable}{RGB}{204, 251, 241}

\begin{document}
% Complete Layered System Workflow for CiviQ Legal RAG System
\begin{tikzpicture}[
    x=1cm, y=1cm,
    font=\sffamily,
    stepbox/.style={
        rectangle, draw=lavenderDark!85, line width=0.95pt, fill=white,
        rounded corners=4pt, inner sep=5pt, align=center,
        drop shadow={shadow xshift=1pt, shadow yshift=-1pt, opacity=0.15}
    },
    flowarr/.style={
        -{Stealth[scale=1.05]}, line width=1.25pt, draw=lavenderDark!90
    },
    dashloop/.style={
        -{Stealth[scale=1.0]}, line width=1.15pt, dashed, draw=prismGold!95!black
    }
]

    % Step 1: Query Ingestion
    \node [stepbox, fill=cInput, text width=2.50cm] (b1) at (-5.75, 0) {
        \textbf{\fontsize{9.2}{11.0}\selectfont\color{lavenderDark}1. Query Ingestion}\\[0.3em]
        \fontsize{7.8}{9.2}\selectfont
        $\bullet$ Natural language text\\
        $\bullet$ Voice audio dictation\\
        $\bullet$ Web Speech API\\
        $\bullet$ FastAPI endpoint
    };

    % Step 2: Language Detection
    \node [stepbox, fill=cStage1, text width=2.50cm] (b2) at (-2.65, 0) {
        \textbf{\fontsize{9.2}{11.0}\selectfont\color{lavenderDark}2. Language Router}\\[0.3em]
        \fontsize{7.8}{9.2}\selectfont
        $\bullet$ Unicode character audit\\
        $\bullet$ Range U+0D00 to U+0D7F\\
        $\bullet$ Auto / Explicit mode\\
        $\bullet$ Branch determination
    };

    % Step 3: English Query Direct Path
    \node [stepbox, fill=cStem, text width=3.85cm] (b3) at (1.45, 1.80) {
        \textbf{\fontsize{9.2}{11.0}\selectfont\color{lavenderDark}3a. English Query Path}\\[0.28em]
        \fontsize{7.6}{9.0}\selectfont
        {\color{prismIndigo}\faCheck}\; \textbf{Direct Processing:} Pass-through\\[0.15em]
        {\color{prismBlue}\faFileContract}\; \textbf{Query Preserved:} Original wording\\[0.15em]
        {\color{prismTeal}\faTag}\; \textbf{Target Language:} English response
    };

    % Step 4: Malayalam Translation Path
    \node [stepbox, fill=cStage2, text width=3.85cm] (b4) at (1.45, -1.80) {
        \textbf{\fontsize{9.2}{11.0}\selectfont\color{lavenderDark}3b. Malayalam Path}\\[0.28em]
        \fontsize{7.6}{9.0}\selectfont
        {\color{prismGold!90!black}\faLanguage}\; \textbf{Gemini Translate:} ML $\to$ English\\[0.15em]
        {\color{prismBlue}\faBook}\; \textbf{Legal Term Adaptation:}\\[0.15em]
        {\color{prismTeal}\faHistory}\; \textbf{Return Translation Tag:} Malayalam
    };

    % Step 5: SBERT Embedding & Pinecone Retrieval
    \node [stepbox, fill=cHead, text width=2.65cm] (b5) at (5.55, 0) {
        \textbf{\fontsize{9.2}{11.0}\selectfont\color{lavenderDark}4. RAG Retrieval}\\[0.3em]
        \fontsize{7.8}{9.2}\selectfont
        $\bullet$ \texttt{all-MiniLM-L6-v2}\\
        $\bullet$ 384-dim dense vector\\
        $\bullet$ Pinecone cosine search\\
        $\bullet$ Fetch top-$k$ chunks
    };

    % Step 6: Grounded Generation & Citation Delivery
    \node [stepbox, fill=cDeliverable!80, text width=3.85cm] (b6) at (9.65, 0) {
        \textbf{\fontsize{9.2}{11.0}\selectfont\color{lavenderDark}5. Grounded Synthesis}\\[0.28em]
        \fontsize{7.6}{9.0}\selectfont
        {\color{prismNavy}\faRobot}\; \textbf{Gemini 2.5 Flash:} Evidence conditioned\\[0.15em]
        {\color{lavenderDark}\faQuoteRight}\; \textbf{Statute Citing:} Act \& Section tags\\[0.15em]
        {\color{prismTeal}\faClock}\; \textbf{Telemetry:} $\approx 1.4$\,s response latency\\[0.15em]
        {\color{red!75!black}\faBalanceScale}\; \textbf{Disclaimer:} Non-legal advice note
    };

    % 1. Connect Step 1 -> Step 2
    \draw [flowarr] (b1) -- (b2);

    % 2. Clean Split from Language Router into English and Malayalam paths
    \coordinate (split1) at ($(b2.east)!0.5!(b3.west |- b2.east)$);
    \draw [line width=1.25pt, draw=lavenderDark!90] (b2.east) -- (split1);
    \draw [flowarr] (split1) |- (b3.west);
    \draw [flowarr] (split1) |- (b4.west);

    % 3. Convergence into Retrieval
    \coordinate (join1) at ($(b3.east |- b5.west)!0.5!(b5.west)$);
    \draw [line width=1.25pt, draw=lavenderDark!90] (b3.east) -| (join1);
    \draw [line width=1.25pt, draw=lavenderDark!90] (b4.east) -| (join1);
    \draw [flowarr] (join1) -- (b5.west);

    % 4. Connect Step 4 -> Step 5
    \draw [flowarr] (b5) -- (b6);

    % 5. User Refinement / Follow-up Feedback Loop
    \draw [dashloop] (b6.north) -- (b6.north |- 0, 3.45) 
        -- node[above=2.5pt, font=\fontsize{8.0}{9.5}\selectfont\bfseries\color{prismGold!95!black}]{User Follow-up Query / Semantic Refinement Loop} (b1.north |- 0, 3.45) 
        -- (b1.north);

\end{tikzpicture}
\end{document}
}
    \caption{Layered system workflow and execution lifecycle across bilingual RAG processing tiers.}
    \label{fig:arch_workflow}
\end{figure}

\subsection{End-to-End System Flow Architecture}
Figure~\ref{fig:system_flow} delineates the complete data transformation pipeline from initial citizen input acquisition to verified output presentation.

\begin{figure}[htbp]
    \centering
    \adjustbox{max width=0.98\linewidth, max height=0.22\textheight, keepaspectratio, center}{\documentclass[tikz,border=8pt]{standalone}
\usepackage{tikz}
\usepackage{fontawesome5}
\usepackage{amsmath,amssymb}
\usetikzlibrary{arrows.meta,positioning,fit,backgrounds,calc,shadows,shapes.geometric}

% Lavender & Slate Color Palette
\definecolor{lavenderDark}{RGB}{15, 44, 89}
\definecolor{lavenderDeep}{RGB}{10, 28, 58}
\definecolor{lavenderMedium}{RGB}{29, 78, 158}
\definecolor{lavenderLight}{RGB}{240, 246, 255}
\definecolor{lavenderBorder}{RGB}{180, 205, 237}
\definecolor{lavenderAccent}{RGB}{2, 132, 199}
\definecolor{prismNavy}{RGB}{15, 44, 89}
\definecolor{prismBlue}{RGB}{29, 78, 158}
\definecolor{prismTeal}{RGB}{13, 148, 136}
\definecolor{prismMint}{RGB}{236, 253, 245}
\definecolor{prismGold}{RGB}{217, 119, 6}
\definecolor{prismDark}{RGB}{15, 23, 42}
\definecolor{prismLight}{RGB}{248, 250, 252}

\begin{document}
% End-to-End System Flow for CiviQ Indian Legal RAG System
\begin{tikzpicture}[
    x=1cm, y=1cm,
    font=\sffamily,
    inout/.style={
        rectangle, rounded corners=3.5pt, line width=0.85pt,
        inner sep=4.5pt, align=center,
        font=\fontsize{7.5}{9.0}\selectfont
    },
    hdrbox/.style={
        rectangle, draw=lavenderDark!80, line width=0.80pt, fill=lavenderDark!14,
        rounded corners=3pt, inner sep=3.0pt, align=center,
        text width=3.30cm,
        font=\fontsize{7.2}{8.5}\selectfont\bfseries\color{lavenderDark}
    },
    stepbox/.style={
        rectangle, draw=black!70, line width=0.60pt, fill=white,
        rounded corners=3pt, inner sep=2.8pt, align=center,
        text width=3.30cm, minimum height=0.72cm,
        font=\fontsize{6.5}{7.6}\selectfont
    },
    stagebox/.style={
        rectangle, draw=lavenderDark!45, dashed, rounded corners=4.5pt, line width=0.70pt,
        fill=lavenderLight!25, inner xsep=3.5pt, inner ysep=4.5pt
    },
    flowarr/.style={
        -{Stealth[scale=0.75]}, line width=0.70pt, draw=black!80
    },
    transarr/.style={
        -{Stealth[length=4.2pt, width=3.2pt]}, line width=1.0pt, draw=lavenderMedium
    }
]

    % 1. TOP: INPUT BAR (width = 14.2cm)
    \node [inout, fill=lavenderLight, draw=lavenderDark!75, text width=14.2cm] (input_node) at (0, 3.60) {
        \textbf{\fontsize{7.8}{9.0}\selectfont\color{lavenderDark}[INPUT]} $\;\longrightarrow\;$
        \textbf{Citizen Query Ingestion}\; $\vert$\;
        {\fontsize{6.6}{7.8}\selectfont Typed Text or Web Speech Voice Audio}\; $\vert$\;
        {\fontsize{6.6}{7.8}\selectfont English / Malayalam $\cdot$ Optional PDF Document Upload}
    };

    % 2. MIDDLE: 3-STAGE HORIZONTAL PIPELINE
    % Spaced at x = -5.15, x = 0.00, x = +5.15 with text width = 3.30cm

    % --- STAGE 1: INGESTION & TRANSLATION (x = -5.15) ---
    \node [hdrbox] (h1) at (-5.15, 2.50) {1. Query Audit \& Translation};
    
    \node [stepbox, below=0.20cm of h1] (s1_1) {
        \textbf{Input Audit \& Parsing}\\[0.08em]
        {\fontsize{6.0}{7.0}\selectfont PyMuPDF parse $\cdot$ PDF $\le 50$\,MB}
    };
    \node [stepbox, below=0.20cm of s1_1] (s1_2) {
        \textbf{Language Detection}\\[0.08em]
        {\fontsize{6.0}{7.0}\selectfont Unicode \texttt{U+0D00}--\texttt{U+0D7F}}
    };
    \node [stepbox, below=0.20cm of s1_2] (s1_3) {
        \textbf{Query Translation}\\[0.08em]
        {\fontsize{6.0}{7.0}\selectfont Gemini: Malayalam $\to$ English}
    };

    \draw [flowarr] (s1_1) -- (s1_2);
    \draw [flowarr] (s1_2) -- (s1_3);

    \begin{scope}[on background layer]
    \node [stagebox, fit=(h1) (s1_3)] (stage1) {};
    \end{scope}

    % --- STAGE 2: EMBEDDING & RETRIEVAL (x = 0.00) ---
    \node [hdrbox] (h2) at (0.00, 2.50) {2. Embedding \& Pinecone Search};
    
    \node [stepbox, below=0.20cm of h2] (s2_1) {
        \textbf{Dense Query Embedding}\\[0.08em]
        {\fontsize{6.0}{7.0}\selectfont \texttt{all-MiniLM-L6-v2} (384-dim)}
    };
    \node [stepbox, below=0.20cm of s2_1] (s2_2) {
        \textbf{Vector Normalization}\\[0.08em]
        {\fontsize{6.0}{7.0}\selectfont Unit L2 norm $\cdot$ 900-word chunks}
    };
    \node [stepbox, below=0.20cm of s2_2] (s2_3) {
        \textbf{Pinecone Top-$k$ Search}\\[0.08em]
        {\fontsize{6.0}{7.0}\selectfont Cosine ANN retrieval ($k=5$)}
    };

    \draw [flowarr] (s2_1) -- (s2_2);
    \draw [flowarr] (s2_2) -- (s2_3);

    \begin{scope}[on background layer]
    \node [stagebox, fit=(h2) (s2_3)] (stage2) {};
    \end{scope}

    % --- STAGE 3: GENERATION & CITATIONS (x = +5.15) ---
    \node [hdrbox] (h3) at (5.15, 2.50) {3. Gemini Grounded Synthesis};
    
    \node [stepbox, below=0.20cm of h3] (s3_1) {
        \textbf{Context Assembly}\\[0.08em]
        {\fontsize{6.0}{7.0}\selectfont Retrieved passages + query}
    };
    \node [stepbox, below=0.20cm of s3_1] (s3_2) {
        \textbf{Gemini Grounded LLM}\\[0.08em]
        {\fontsize{6.0}{7.0}\selectfont Gemini 2.5 Flash $\cdot$ Citations}
    };
    \node [stepbox, below=0.20cm of s3_2, fill=lavenderLight] (s3_3) {
        \textbf{Response Delivery}\\[0.08em]
        {\fontsize{6.0}{7.0}\selectfont Translation back $\to$ React UI}
    };

    \draw [flowarr] (s3_1) -- (s3_2);
    \draw [flowarr] (s3_2) -- (s3_3);

    \begin{scope}[on background layer]
    \node [stagebox, fit=(h3) (s3_3)] (stage3) {};
    \end{scope}

    % Common elevation coordinate for perfectly horizontal inter-stage arrows
    \path let \p1 = (s1_2.center) in coordinate (transY) at (0, \y1);

    % Horizontal Transition Arrows in the Gap (2-line compact label, zero collision)
    \draw [transarr] (stage1.east |- transY) -- node[above=1pt, align=center, font=\fontsize{5.8}{6.8}\selectfont\bfseries\color{lavenderMedium}]{English\\Query} (stage2.west |- transY);
    \draw [transarr] (stage2.east |- transY) -- node[above=1pt, align=center, font=\fontsize{5.8}{6.8}\selectfont\bfseries\color{lavenderMedium}]{Top-$k$\\Evidence} (stage3.west |- transY);

    % Connect Top Input to Stage 1
    \draw [flowarr, line width=0.80pt] (input_node.south -| stage1.north) -- (stage1.north);

    % 3. BOTTOM: OUTPUT BAR (Positioned below stage boxes with 0.55cm clean margin)
    \node [inout, fill=prismMint!65, draw=prismTeal!80, text width=14.2cm, below=0.55cm of stage2] (output_node) {
        \textbf{\fontsize{7.8}{9.0}\selectfont\color{prismTeal!90!black}[OUTPUT]} $\;\longleftarrow\;$
        \textbf{Accessible Legal Guidance with Grounded Citations}\; $\vert$\;
        {\fontsize{6.6}{7.8}\selectfont Plain-Language Answer in English / Malayalam}\; $\vert$\;
        {\fontsize{6.6}{7.8}\selectfont Expandable Document Cards with PDF Page References}
    };

    % Connect Stage 3 to Output Bar
    \draw [flowarr, line width=0.80pt] (stage3.south) -- (output_node.north -| stage3.south);

\end{tikzpicture}
\end{document}
}
    \caption{Complete end-to-end data transformation pipeline from input acquisition through vector search to grounded answer synthesis.}
    \label{fig:system_flow}
\end{figure}

\section{System Flowchart and Core System Algorithms}

\subsection{System Decision Flowchart}
Figure~\ref{fig:multiagent_flowchart} details the operational decision logic, linguistic routing branches, and fault-tolerance pathways governing inference execution. The decision flowchart establishes four critical operational checkpoints:

\vspace{2pt}
\noindent\textbf{1. Modality Arbitration:} Checks whether input arrived as raw text or spoken voice audio, invoking the Web Speech API for transcription when voice mode is engaged.\par\vspace{2pt}
\noindent\textbf{2. Linguistic Evaluation:} Audits the query for characters within the Unicode Malayalam block ($\mathrm{U{+}0D00}$ to $\mathrm{U{+}0D7F}$), routing to Gemini for English normalization.\par\vspace{2pt}
\noindent\textbf{3. Relevance Thresholding ($\theta \ge 0.45$):} Evaluates Pinecone cosine similarity against $\theta = 0.45$. Scores below $\theta$ trigger a fallback response, preventing hallucinations.\par\vspace{2pt}
\noindent\textbf{4. Reverse Translation Dispatch:} If the input originated in Malayalam, the verified English answer is translated back to Malayalam prior to client UI rendering.

\clearpage
\begin{figure}[p]
    \centering
    \adjustbox{max width=0.98\linewidth, max height=0.88\textheight, keepaspectratio, center}{\documentclass[tikz,border=10pt]{standalone}
\usepackage{tikz}
\usepackage{fontawesome5}
\usepackage{amsmath,amssymb}
\usetikzlibrary{arrows.meta,positioning,fit,backgrounds,calc,shadows,shapes.geometric}

% Color definitions
\definecolor{lavenderDark}{RGB}{15, 44, 89}
\definecolor{lavenderDeep}{RGB}{10, 28, 58}
\definecolor{lavenderMedium}{RGB}{29, 78, 158}
\definecolor{lavenderLight}{RGB}{240, 246, 255}
\definecolor{lavenderBorder}{RGB}{180, 205, 237}
\definecolor{prismNavy}{RGB}{15, 44, 89}
\definecolor{prismBlue}{RGB}{29, 78, 158}
\definecolor{prismTeal}{RGB}{13, 148, 136}
\definecolor{prismGold}{RGB}{217, 119, 6}
\definecolor{prismDark}{RGB}{15, 23, 42}
\definecolor{statusFail}{RGB}{220, 38, 38}
\definecolor{cInput}{RGB}{241, 245, 249}
\definecolor{cStem}{RGB}{224, 242, 254}
\definecolor{cStage1}{RGB}{209, 250, 229}
\definecolor{cStage2}{RGB}{254, 243, 199}
\definecolor{cHead}{RGB}{219, 234, 254}

\begin{document}
% Complete Flowchart for CiviQ Legal RAG Query Processing & Translation
\begin{tikzpicture}[
  x=1cm, y=1cm,
  font=\sffamily,
  startend/.style={
    ellipse, draw=green!60!black, fill=green!45!black, text=white,
    line width=0.95pt, inner sep=4.0pt, font=\fontsize{8.5}{9.5}\selectfont\bfseries,
    minimum width=2.20cm, minimum height=0.68cm, drop shadow={opacity=0.12}
  },
  procbox/.style={
    rectangle, rounded corners=3.5pt,
    draw=lavenderDark!80, fill=white, text=black,
    line width=0.85pt, inner sep=4.0pt, font=\fontsize{8.0}{9.2}\selectfont,
    minimum height=0.72cm, text width=4.80cm, align=center, drop shadow={opacity=0.10}
  },
  decision/.style={
    diamond, aspect=2.3,
    draw=prismNavy, fill=lavenderLight, text=prismNavy,
    line width=0.85pt, inner sep=2.5pt, font=\fontsize{7.6}{8.6}\selectfont\bfseries,
    align=center, drop shadow={opacity=0.10}
  },
  sidebox/.style={
    rectangle, rounded corners=3.5pt,
    draw=prismTeal!90, fill=cStage1, text=prismNavy,
    line width=0.80pt, inner sep=3.5pt, font=\fontsize{7.5}{8.5}\selectfont,
    minimum height=0.66cm, text width=2.85cm, align=center, drop shadow={opacity=0.08}
  },
  errorbox/.style={
    rectangle, rounded corners=3.5pt,
    draw=statusFail!80, fill=statusFail!12, text=statusFail!90!black,
    line width=0.80pt, inner sep=3.5pt, font=\fontsize{7.5}{8.5}\selectfont\bfseries,
    minimum height=0.66cm, text width=2.80cm, align=center, drop shadow={opacity=0.08}
  },
  flowarr/.style={
    -{Stealth[scale=0.90]}, line width=0.90pt, draw=lavenderDark!90
  },
  failarr/.style={
    -{Stealth[scale=0.85]}, line width=0.85pt, draw=statusFail!90, dashed
  }
]

  % Central Column - Main Processing Flow
  \node[startend] (start) at (0, 11.20) {START};

  \node[procbox, fill=cInput, below=0.32cm of start] (user_in) {%
    \textbf{User Query Ingestion}\\[1.5pt]
    {\fontsize{7.0}{8.0}\selectfont Text or Web Speech Dictation}%
  };

  \node[decision, below=0.35cm of user_in] (d_voice) {%
    Voice Input\\Mode?%
  };

  \node[decision, below=0.35cm of d_voice] (d_lang) {%
    Malayalam Detected?\\{\fontsize{6.8}{7.8}\selectfont (U+0D00 to U+0D7F)}%
  };

  \node[procbox, fill=cStem, below=0.35cm of d_lang] (embed) {%
    \textbf{Query Vector Embedding}\\[1.5pt]
    {\fontsize{7.0}{8.0}\selectfont all-MiniLM-L6-v2 $\to$ 384-d normalized vector}%
  };

  \node[procbox, fill=cStage2, below=0.32cm of embed] (pinecone) {%
    \textbf{Pinecone ANN Retrieval}\\[1.5pt]
    {\fontsize{7.0}{8.0}\selectfont Cosine similarity search, top-$k$ passages}%
  };

  \node[decision, below=0.35cm of pinecone] (d_match) {%
    Relevant Passages\\$\text{sim} \ge \theta$?%
  };

  \node[procbox, fill=cHead, below=0.35cm of d_match] (prompt_bld) {%
    \textbf{Context Prompt Assembly}\\[1.5pt]
    {\fontsize{7.0}{8.0}\selectfont Query + top-$k$ statutory texts + citations}%
  };

  \node[procbox, fill=cStage1, below=0.32cm of prompt_bld] (gemini) {%
    \textbf{Gemini 2.5 Flash Generation}\\[1.5pt]
    {\fontsize{7.0}{8.0}\selectfont Grounded answer synthesis ($T=0.3$, max 1024 tokens)}%
  };

  \node[decision, below=0.35cm of gemini] (d_trans_back) {%
    Translate to\\Malayalam?%
  };

  \node[procbox, fill=cInput, below=0.35cm of d_trans_back] (respond) {%
    \textbf{JSON Response Packaging}\\[1.5pt]
    {\fontsize{7.0}{8.0}\selectfont Answer, sources, citations, latency $\approx 1.4$\,s}%
  };

  \node[procbox, fill=lavenderLight, below=0.32cm of respond] (render) {%
    \textbf{React UI Display}\\[1.5pt]
    {\fontsize{7.0}{8.0}\selectfont Markdown answer + statutory citation cards}%
  };

  \node[startend, fill=lavenderDark, below=0.32cm of render] (end) {END};

  % Side Branches (Left/Right)
  % 1. Voice transcription (Right of d_voice)
  \node[sidebox, right=1.60cm of d_voice] (voice_stt) {%
    \textbf{Web Speech API}\\[1.5pt]
    {\fontsize{7.0}{8.0}\selectfont Speech-to-text transcribe}%
  };

  % 2. Translation to English (Right of d_lang)
  \node[sidebox, right=1.60cm of d_lang] (trans_en) {%
    \textbf{Gemini Translate}\\[1.5pt]
    {\fontsize{7.0}{8.0}\selectfont Malayalam $\to$ English}%
  };

  % 3. No match error (Left of d_match)
  \node[errorbox, left=1.60cm of d_match] (no_match) {%
    \textbf{Fallback Output}\\[1.5pt]
    {\fontsize{7.0}{8.0}\selectfont "No statutory info found"}%
  };

  % 4. Translation to Malayalam (Right of d_trans_back)
  \node[sidebox, right=1.60cm of d_trans_back] (trans_ml) {%
    \textbf{Gemini Translate}\\[1.5pt]
    {\fontsize{7.0}{8.0}\selectfont English $\to$ Malayalam}%
  };

  % Downward Connections
  \draw[flowarr] (start) -- (user_in);
  \draw[flowarr] (user_in) -- (d_voice);
  \draw[flowarr] (d_voice) -- node[left, font=\fontsize{7.2}{8.2}\selectfont\bfseries]{No} (d_lang);
  \draw[flowarr] (d_lang) -- node[left, font=\fontsize{7.2}{8.2}\selectfont\bfseries]{No} (embed);
  \draw[flowarr] (embed) -- (pinecone);
  \draw[flowarr] (pinecone) -- (d_match);
  \draw[flowarr] (d_match) -- node[left, font=\fontsize{7.2}{8.2}\selectfont\bfseries]{Yes} (prompt_bld);
  \draw[flowarr] (prompt_bld) -- (gemini);
  \draw[flowarr] (gemini) -- (d_trans_back);
  \draw[flowarr] (d_trans_back) -- node[left, font=\fontsize{7.2}{8.2}\selectfont\bfseries]{No} (respond);
  \draw[flowarr] (respond) -- (render);
  \draw[flowarr] (render) -- (end);

  % Branch: Voice transcription
  \draw[flowarr] (d_voice.east) -- node[above, font=\fontsize{7.2}{8.2}\selectfont\bfseries]{Yes} (voice_stt.west);
  \draw[flowarr] (voice_stt.south) |- (d_lang.north east);

  % Branch: Translation to English
  \draw[flowarr] (d_lang.east) -- node[above, font=\fontsize{7.2}{8.2}\selectfont\bfseries]{Yes} (trans_en.west);
  \draw[flowarr] (trans_en.south) |- (embed.north east);

  % Branch: Fallback when no match
  \draw[failarr] (d_match.west) -- node[above, font=\fontsize{7.2}{8.2}\selectfont\bfseries]{No} (no_match.east);
  \draw[failarr] (no_match.south) |- (respond.north west);

  % Branch: Translation to Malayalam
  \draw[flowarr] (d_trans_back.east) -- node[above, font=\fontsize{7.2}{8.2}\selectfont\bfseries]{Yes} (trans_ml.west);
  \draw[flowarr] (trans_ml.south) |- (respond.north east);

\end{tikzpicture}
\end{document}
}
    \caption{Query processing flowchart with decision branches for voice transcription, language detection, relevance thresholding, and translation.}
    \label{fig:multiagent_flowchart}
\end{figure}
\clearpage

\subsection{Formal Algorithmic Specifications of Core Functions}

To provide rigorous operational transparency, the computational pipeline of the CiviQ legal information retrieval platform is decomposed into discrete mathematical algorithms. The system functions across three tightly coupled operational tiers:
\begin{enumerate}
    \item \textbf{Deterministic Offline Ingestion and Vectorization (Algorithms~\ref{alg:chunking} and~\ref{alg:embedding_ingest}):} Converts raw statutory enactments into clean, overlapping textual windows ($c=900, o=150$) and projects them onto a unit hypersphere ($\mathbb{S}^{383} \subset \mathbb{R}^{384}$) using Layer-6 mean-pooled representations from \texttt{all-MiniLM-L6-v2}.
    \item \textbf{Linguistic Routing and Semantic Evidence Retrieval (Algorithms~\ref{alg:multilingual_router} and~\ref{alg:ann_retrieval}):} Performs $O(|q|)$ Unicode range audits for Malayalam script identification ($\mathrm{U{+}0D00}$ to $\mathrm{U{+}0D7F}$), executes asynchronous neural query translation, and conducts logarithmic-time ($O(\log |\mathcal{C}|)$) approximate nearest neighbor (ANN) search across Pinecone's HNSW graph index.
    \item \textbf{Grounded LLM Context Synthesis and Ad-Hoc Stream Ingestion (Algorithms~\ref{alg:llm_synthesis} and~\ref{alg:doc_analysis}):} Assembles prompt contexts with explicit document and section boundary headers, enforces deterministic citation generation ($T=0.3$), verifies confidence via relevance threshold gating ($\theta \ge 0.45$), and supports session-ephemeral multipart PDF document ingestion without persistent disk retention.
\end{enumerate}

\subsubsection{Execution Pipeline Characteristics and Data Contracts}
The platform enforces strict data contracts and operational invariants at every pipeline transition to guarantee deterministic outputs, statutory citation attribution, and sub-second retrieval latency:

\noindent\textbf{1. Corpus Ingestion Contract ($D \to \mathcal{C}$):} Statutory enactments are ingested via PyMuPDF with UTF-8 integrity validation. Gazette banners, running headers, and repetitive legislative footers are purged via regular expressions. Slicing with window size $c=900$ words and stride $s=750$ words ($o=150$) preserves cross-boundary legislative context, while sub-$\theta_{\min}$ ($40$ words) fragments are rejected to eliminate noise.

\noindent\textbf{2. Vector Embedding Contract ($\mathcal{C} \to \mathbb{R}^{384}$):} Text chunks are encoded using \texttt{all-MiniLM-L6-v2} with Layer-6 mean pooling and projected onto the unit hypersphere ($\|\mathbf{e}\|_2 = 1.0$), ensuring that approximate nearest neighbor graph traversal in Pinecone computes exact cosine similarities.

\noindent\textbf{3. Grounded Synthesis Contract ($q \to \mathcal{Y}$):} Queries are scanned for Malayalam Unicode codepoints ($\mathrm{U{+}0D00}$ to $\mathrm{U{+}0D7F}$) and translated to English. Candidate chunks scoring above relevance threshold $\theta \ge 0.45$ are assembled into context blocks. Gemini~2.5 Flash operates deterministically at $T=0.3$ with strict attribution prompts, returning structured citations and latency telemetry.

The complete mathematical specifications of each stage are formalized in Algorithms~\ref{alg:chunking} through~\ref{alg:doc_analysis}.

\begin{formalalgo}{Corpus Preprocessing, Sanitization, and Overlapping Semantic Chunking}{alg:chunking}
\textbf{Input:} Statute collection $\mathcal{D} = \{D_1, \dots, D_N\}$, chunk size $c = 900$, stride $o = 150$, min threshold $\theta_{\min} = 40$\\
\textbf{Output:} Preprocessed chunk repository $\mathcal{C} = \{c_1, \dots, c_M\}$ with statutory metadata tuples
\vspace{3pt}\hrule\vspace{4pt}
\begin{enumerate}
  \item[\textbf{1:}] Initialize chunk repository: $\mathcal{C} \gets \emptyset$.
  \item[\textbf{2:}] \textbf{for each} document $D_k \in \mathcal{D}$ \textbf{do}
  \item[\textbf{3:}] \quad Open document stream via PyMuPDF: $\mathrm{doc} \gets \mathrm{fitz.open}(D_k)$.
  \item[\textbf{4:}] \quad Extract text per page: $T_{\mathrm{raw}} \gets \bigcup_{p=1}^{|\mathrm{doc}|} \mathrm{extract\_text}(p)$.
  \item[\textbf{5:}] \quad Strip headers, footers, and gazette banners: $T_{\mathrm{clean}} \gets \mathrm{RegexFilter}(T_{\mathrm{raw}})$.
  \item[\textbf{6:}] \quad Normalize whitespace: $T_{\mathrm{clean}} \gets \mathrm{CollapseWhitespace}(T_{\mathrm{clean}})$.
  \item[\textbf{7:}] \quad Tokenize cleaned text into word array: $W \gets [w_1, w_2, \dots, w_K]$.
  \item[\textbf{8:}] \quad Compute stride: $s \gets c - o = 900 - 150 = 750$.
  \item[\textbf{9:}] \quad \textbf{for} $i = 0, 1, \dots, \lfloor (K - c)/s \rfloor$ \textbf{do}
  \item[\textbf{10:}] \quad \quad Slice chunk window: $\mathrm{words} \gets W[i \cdot s : i \cdot s + c]$.
  \item[\textbf{11:}] \quad \quad \textbf{if} $|\mathrm{words}| \ge \theta_{\min}$ \textbf{then}
  \item[\textbf{12:}] \quad \quad \quad Construct chunk record: $c_j \gets \mathrm{Chunk}(\mathrm{words}, \mathrm{doc}=D_k, \mathrm{idx}=i)$.
  \item[\textbf{13:}] \quad \quad \quad $\mathcal{C} \gets \mathcal{C} \cup \{c_j\}$.
  \item[\textbf{14:}] \quad \quad \textbf{end if}
  \item[\textbf{15:}] \quad \textbf{end for}
  \item[\textbf{16:}] \textbf{end for}
  \item[\textbf{17:}] \textbf{return} Chunk repository $\mathcal{C}$.
\end{enumerate}
\end{formalalgo}

\begin{formalalgo}{Dense Semantic Vector Encoding and Pinecone HNSW Vector Index Ingestion}{alg:embedding_ingest}
\textbf{Input:} Chunk repository $\mathcal{C} = \{c_1, \dots, c_M\}$, encoder $\phi$ (\texttt{all-MiniLM-L6-v2}), batch limit $B = 100$, index $\mathcal{I}$\\
\textbf{Output:} Indexed vector store populated with 384-dimensional normalized embeddings
\vspace{3pt}\hrule\vspace{4pt}
\begin{enumerate}
  \item[\textbf{1:}] Partition chunk repository into batches: $\mathcal{B} \gets \{b_1, b_2, \dots, b_{\lceil M/B \rceil}\}$.
  \item[\textbf{2:}] \textbf{for each} batch $b_u \in \mathcal{B}$ \textbf{do}
  \item[\textbf{3:}] \quad Compute dense vectors via Layer-6 pooling: $\mathbf{Z}_u \gets \phi(\mathrm{texts}(b_u)), \;\mathbf{Z}_u \in \mathbb{R}^{B \times 384}$.
  \item[\textbf{4:}] \quad Apply unit $L_2$ normalization: $\mathbf{E}_u(i) \gets \mathbf{Z}_u(i) / \|\mathbf{Z}_u(i)\|_2, \;\forall i \in [1, B]$.
  \item[\textbf{5:}] \quad Assemble Pinecone payload: $\mathcal{V}_u \gets \{(\mathrm{id}(c_i), \mathbf{E}_u(i), \mathrm{meta}(c_i))\}_{i=1}^B$.
  \item[\textbf{6:}] \quad Execute cloud upsert: $\mathcal{I}.\mathrm{upsert}(\mathrm{vectors}=\mathcal{V}_u, \mathrm{namespace}=\text{``legal-statutes''})$.
  \item[\textbf{7:}] \textbf{end for}
  \item[\textbf{8:}] Verify total indexed vectors: $\mathcal{I}.\mathrm{describe\_index\_stats}()$.
  \item[\textbf{9:}] \textbf{return} Confirmation status $\mathrm{True}$.
\end{enumerate}
\end{formalalgo}

\begin{formalalgo}{Multilingual Query Normalization and Malayalam-to-English Translation Bridge}{alg:multilingual_router}
\textbf{Input:} Raw query string $q_{\mathrm{raw}}$, neural translation service $\mathcal{M}_{\mathrm{trans}}$ (Gemini~2.5)\\
\textbf{Output:} Standardized English query $q_{\mathrm{en}}$, language routing flag $\tau_{\mathrm{lang}} \in \{\mathrm{EN}, \mathrm{ML}\}$
\vspace{3pt}\hrule\vspace{4pt}
\begin{enumerate}
  \item[\textbf{1:}] Trim whitespace: $q_{\mathrm{trim}} \gets \mathrm{Trim}(q_{\mathrm{raw}})$.
  \item[\textbf{2:}] Initialize flag: $\mathrm{has\_ml} \gets \mathrm{False}$.
  \item[\textbf{3:}] \textbf{for each} character $c \in q_{\mathrm{trim}}$ \textbf{do}
  \item[\textbf{4:}] \quad \textbf{if} $\mathrm{codepoint}(c) \in [\mathrm{0x0D00}, \mathrm{0x0D7F}]$ \textbf{then}
  \item[\textbf{5:}] \quad \quad $\mathrm{has\_ml} \gets \mathrm{True}$; \textbf{break}
  \item[\textbf{6:}] \quad \textbf{end if}
  \item[\textbf{7:}] \textbf{end for}
  \item[\textbf{8:}] \textbf{if} $\mathrm{has\_ml} = \mathrm{True}$ \textbf{then}
  \item[\textbf{9:}] \quad Set language: $\tau_{\mathrm{lang}} \gets \mathrm{ML}$.
  \item[\textbf{10:}] \quad Form translation prompt: $\mathcal{P}_{\mathrm{trans}} \gets \text{TranslatePrompt}(q_{\mathrm{trim}})$.
  \item[\textbf{11:}] \quad Execute translation: $q_{\mathrm{en}} \gets \mathcal{M}_{\mathrm{trans}}.\mathrm{generate}(\mathcal{P}_{\mathrm{trans}}, T=0.1)$.
  \item[\textbf{12:}] \textbf{else}
  \item[\textbf{13:}] \quad Set language: $\tau_{\mathrm{lang}} \gets \mathrm{EN}$; $\quad q_{\mathrm{en}} \gets q_{\mathrm{trim}}$.
  \item[\textbf{14:}] \textbf{end if}
  \item[\textbf{15:}] \textbf{return} $(q_{\mathrm{en}}, \tau_{\mathrm{lang}})$.
\end{enumerate}
\end{formalalgo}

\begin{formalalgo}{Cosine Approximate Nearest Neighbor Retrieval and Passage Scoring}{alg:ann_retrieval}
\textbf{Input:} Query $q_{\mathrm{en}}$, encoder $\phi$, Pinecone index $\mathcal{I}$, depth $k \in [1, 10]$, relevance threshold $\theta = 0.45$\\
\textbf{Output:} Top-$k$ ranked statutory chunks $\mathcal{R}_k = \{(c_{(1)}, s_{(1)}), \dots, (c_{(m)}, s_{(m)})\}$
\vspace{3pt}\hrule\vspace{4pt}
\begin{enumerate}
  \item[\textbf{1:}] Dense vector embedding: $\mathbf{z}_q \gets \phi(q_{\mathrm{en}}), \;\mathbf{z}_q \in \mathbb{R}^{384}$.
  \item[\textbf{2:}] Project onto unit sphere ($L_2$ norm): $\mathbf{e}_q \gets \mathbf{z}_q / \|\mathbf{z}_q\|_2$.
  \item[\textbf{3:}] Query Pinecone HNSW index:
  \item[\textbf{4:}] \quad $\mathcal{M}_{\mathrm{raw}} \gets \mathcal{I}.\mathrm{query}(\mathrm{vector}=\mathbf{e}_q, \mathrm{top\_k}=k, \mathrm{include\_metadata}=\mathrm{True})$.
  \item[\textbf{5:}] Compute cosine similarities: $s_i \gets \mathbf{e}_q^\top \mathbf{e}_{c_i}, \;\forall c_i \in \mathcal{M}_{\mathrm{raw}}$.
  \item[\textbf{6:}] \textbf{if} $\max_i (s_i) < \theta$ \textbf{then}
  \item[\textbf{7:}] \quad Low semantic relevance detected: $\mathcal{R}_k \gets \emptyset$.
  \item[\textbf{8:}] \textbf{else}
  \item[\textbf{9:}] \quad Filter and sort by descending similarity:
  \item[\textbf{10:}] \quad \quad $\mathcal{R}_k \gets \{(c_i, s_i) \in \mathcal{M}_{\mathrm{raw}} : s_i \ge \theta\}, \quad s_{(1)} \ge \dots \ge s_{(m)}$.
  \item[\textbf{11:}] \textbf{end if}
  \item[\textbf{12:}] \textbf{return} Ranked statutory evidence $\mathcal{R}_k$.
\end{enumerate}
\end{formalalgo}

\begin{formalalgo}{Grounded LLM Context Synthesis, Citation Attribution, and Response Packaging}{alg:llm_synthesis}
\textbf{Input:} Query $q_{\mathrm{en}}$, passages $\mathcal{R}_k$, language flag $\tau_{\mathrm{lang}}$, generator $\mathcal{G}$ (Gemini~2.5 Flash, $T=0.3$)\\
\textbf{Output:} Complete structured legal response packet $\mathcal{Y} = \{\mathrm{answer}, \mathrm{citations}, \mathrm{sources}, \mathrm{latency}\}$
\vspace{3pt}\hrule\vspace{4pt}
\begin{enumerate}
  \item[\textbf{1:}] Record start timestamp: $t_{\mathrm{start}} \gets \mathrm{time.now}()$.
  \item[\textbf{2:}] \textbf{if} $\mathcal{R}_k = \emptyset$ \textbf{then}
  \item[\textbf{3:}] \quad Assign fallback: $A_{\mathrm{en}} \gets \text{``No relevant statutory provisions found in corpus.''}$.
  \item[\textbf{4:}] \quad Assign empty citation list: $\mathcal{S} \gets \emptyset$.
  \item[\textbf{5:}] \textbf{else}
  \item[\textbf{6:}] \quad Concatenate retrieved excerpts with citation headers:
  \item[\textbf{7:}] \quad \quad $\mathcal{X}_{\mathrm{ctx}} \gets \bigoplus_{i=1}^{|\mathcal{R}_k|} [\text{Source: } c_i.\mathrm{doc} \parallel \text{Sec. } c_i.\mathrm{sec} \parallel c_i.\mathrm{text}]$.
  \item[\textbf{8:}] \quad Formulate citation-enforcing instruction prompt:
  \item[\textbf{9:}] \quad \quad $\mathcal{P}_{\mathrm{gen}} \gets \mathrm{PromptTemplate}(\mathrm{context}=\mathcal{X}_{\mathrm{ctx}}, \mathrm{question}=q_{\mathrm{en}})$.
  \item[\textbf{10:}] \quad Deterministic inference: $A_{\mathrm{en}} \gets \mathcal{G}.\mathrm{generate}(\mathcal{P}_{\mathrm{gen}}, T=0.3, \mathrm{max\_tokens}=1024)$.
  \item[\textbf{11:}] \quad Extract statutory section citations via regex: $\mathcal{S} \gets \mathrm{ExtractCitations}(A_{\mathrm{en}})$.
  \item[\textbf{12:}] \textbf{end if}
  \item[\textbf{13:}] \textbf{if} $\tau_{\mathrm{lang}} = \mathrm{ML}$ \textbf{then}
  \item[\textbf{14:}] \quad Translate answer: $A_{\mathrm{final}} \gets \mathcal{G}.\mathrm{translate}(A_{\mathrm{en}}, \text{``EN} \to \text{ML''})$.
  \item[\textbf{15:}] \textbf{else}
  \item[\textbf{16:}] \quad $A_{\mathrm{final}} \gets A_{\mathrm{en}}$.
  \item[\textbf{17:}] \textbf{end if}
  \item[\textbf{18:}] Compute latency: $\Delta t \gets \mathrm{time.now}() - t_{\mathrm{start}}$.
  \item[\textbf{19:}] Package response: $\mathcal{Y} \gets \{A_{\mathrm{final}}, \mathcal{S}, \mathcal{R}_k, \Delta t\}$.
  \item[\textbf{20:}] \textbf{return} Response packet $\mathcal{Y}$.
\end{enumerate}
\end{formalalgo}

\begin{formalalgo}{Asynchronous Multipart Document Stream Ingestion and Ad-Hoc Analysis}{alg:doc_analysis}
\textbf{Input:} Document stream $\mathcal{D}_{\mathrm{user}}$, directive mode $\delta \in \{\mathrm{summarize}, \mathrm{question}\}$, optional query $q_{\mathrm{doc}}$\\
\textbf{Output:} Document breakdown summary or grounded query analysis $\mathcal{Z}$
\vspace{3pt}\hrule\vspace{4pt}
\begin{enumerate}
  \item[\textbf{1:}] Payload check: $\mathbf{if} \;\mathrm{size}(\mathcal{D}_{\mathrm{user}}) > 50\,\mathrm{MB}\; \mathbf{then}\; \mathbf{throw} \;\mathrm{HTTP\_413}$.
  \item[\textbf{2:}] MIME audit: $\mathbf{if} \;\mathrm{mime}(\mathcal{D}_{\mathrm{user}}) \neq \text{``application/pdf''}\; \mathbf{then}\; \mathbf{throw} \;\mathrm{HTTP\_415}$.
  \item[\textbf{3:}] Stream raw bytes via PyMuPDF: $\mathrm{doc\_stream} \gets \mathrm{fitz.open}(\mathrm{stream}=\mathcal{D}_{\mathrm{user}})$.
  \item[\textbf{4:}] Extract page-labeled text: $T_{\mathrm{doc}} \gets \bigcup_{p=1}^{|\mathrm{doc\_stream}|} (\text{``[Page ''} + p + \text{``] ''} + \mathrm{extract\_text}(p))$.
  \item[\textbf{5:}] \textbf{if} $\delta = \mathrm{summarize}$ \textbf{then}
  \item[\textbf{6:}] \quad Formulate prompt: $\mathcal{P} \gets \mathrm{SummaryPrompt}(T_{\mathrm{doc}})$.
  \item[\textbf{7:}] \quad Generate summary: $\mathcal{Z} \gets \mathrm{Gemini}.\mathrm{generate}(\mathcal{P}, T=0.2)$.
  \item[\textbf{8:}] \textbf{else if} $\delta = \mathrm{question}$ \textbf{then}
  \item[\textbf{9:}] \quad Formulate prompt: $\mathcal{P} \gets \mathrm{QAPrompt}(T_{\mathrm{doc}}, q_{\mathrm{doc}})$.
  \item[\textbf{10:}] \quad Generate analysis: $\mathcal{Z} \gets \mathrm{Gemini}.\mathrm{generate}(\mathcal{P}, T=0.2)$.
  \item[\textbf{11:}] \textbf{end if}
  \item[\textbf{12:}] Release in-memory buffer immediately: $\mathrm{doc\_stream.close}()$.
  \item[\textbf{13:}] \textbf{return} Structured document response $\mathcal{Z}$.
\end{enumerate}
\end{formalalgo}

\subsection{Algorithmic Complexity, Latency Bounds, and Guarantees}
The theoretical and operational characteristics of the core algorithms are summarized in Table~\ref{tab:algo_complexity}.

\begin{table}[htbp]
\centering
\caption{Algorithmic Complexity, Memory Footprint, and Execution Bounds.}
\label{tab:algo_complexity}
\small
\begin{tabularx}{\textwidth}{l >{\centering\arraybackslash}p{2.4cm} >{\centering\arraybackslash}p{2.4cm} >{\centering\arraybackslash}p{2.2cm} X}
\toprule
\textbf{Algorithm} & \textbf{Time Complexity} & \textbf{Space Complexity} & \textbf{Typical Latency} & \textbf{Theoretical Guarantee} \\
\midrule
Alg.~\ref{alg:chunking} (Chunking) & $O(N \cdot L)$ & $O(M \cdot c)$ & $\approx 3.2$\,s (corpus) & Preserves cross-boundary statutory context \\
Alg.~\ref{alg:embedding_ingest} (Embed) & $O(M \cdot d \cdot \ell)$ & $O(B \cdot d)$ & $\approx 42$\,s (corpus) & Exact unit-norm hypersphere projection \\
Alg.~\ref{alg:multilingual_router} (Routing) & $O(|q|)$ & $O(|q|)$ & $1.8$\,ms & Deterministic Malayalam Unicode detection \\
Alg.~\ref{alg:ann_retrieval} (ANN Search) & $O(\log |\mathcal{C}|)$ & $O(k \cdot d)$ & $58$\,ms & Logarithmic HNSW query complexity \\
Alg.~\ref{alg:llm_synthesis} (Synthesis) & $O(k \cdot c + \tau_{\mathrm{out}})$ & $O(k \cdot c)$ & $1,140$\,ms & Zero-hallucination statutory grounding \\
Alg.~\ref{alg:doc_analysis} (Upload QA) & $O(P \cdot L_{\mathrm{page}})$ & $O(\mathrm{size}(\mathcal{D}))$ & $\approx 2.1$\,s & Zero data retention; session ephemeral \\
\bottomrule
\end{tabularx}
\end{table}

\subsection{RAG Pipeline Component Dependencies}
Figure~\ref{fig:dag_scheduling_walkthrough} illustrates the component dependencies in the RAG ingestion and query pipeline. The directed acyclic execution graph traces the transformation of raw PDF legal statutes into cleaned, normalized text chunks, followed by embedding generation and Pinecone vector upload. Query execution parallels this workflow: incoming text or voice queries are embedded with the identical transformer encoder, enabling cosine distance comparison against the pre-indexed legal corpus, subsequent context compilation, and Gemini-based synthesis.

\begin{figure}[htbp]
    \centering
    \adjustbox{max width=0.98\linewidth, keepaspectratio, center}{\documentclass[tikz,border=10pt]{standalone}
\usepackage{tikz}
\usepackage{fontawesome5}
\usepackage{amsmath,amssymb}
\usetikzlibrary{arrows.meta,positioning,fit,backgrounds,calc,shadows,shapes.geometric}

% Professional Palette
\definecolor{prismNavy}{RGB}{15, 44, 89}
\definecolor{prismBlue}{RGB}{29, 78, 158}
\definecolor{prismTeal}{RGB}{13, 148, 136}
\definecolor{prismMint}{RGB}{236, 253, 245}
\definecolor{prismGold}{RGB}{202, 138, 4}
\definecolor{prismDark}{RGB}{15, 23, 42}
\definecolor{prismSlate}{RGB}{71, 85, 105}
\definecolor{cInput}{RGB}{248, 250, 252}
\definecolor{cStem}{RGB}{240, 249, 255}
\definecolor{cStage1}{RGB}{240, 253, 250}
\definecolor{cStage2}{RGB}{254, 252, 232}
\definecolor{cStage3}{RGB}{254, 242, 242}
\definecolor{cStage4}{RGB}{238, 242, 255}

\begin{document}
% RAG Pipeline Component Dependency DAG
\begin{tikzpicture}[
  x=1cm, y=1cm,
  font=\sffamily,
  dagarrow/.style={-{Stealth[length=4.0pt, width=3.0pt]}, line width=0.95pt, draw=prismNavy},
  dagnode/.style={
    rectangle, rounded corners=3.5pt, line width=0.75pt,
    minimum height=1.15cm, text width=1.75cm, inner sep=3.0pt, align=center,
    font=\fontsize{6.2}{7.3}\selectfont,
    drop shadow={shadow xshift=0.6pt, shadow yshift=-0.6pt, opacity=0.08}
  },
  idxnode/.style={
    cylinder, shape border rotate=90, aspect=0.25, line width=0.85pt,
    draw=prismTeal, fill=prismMint, minimum height=1.35cm, text width=1.70cm, inner sep=2.0pt, align=center,
    font=\fontsize{6.5}{7.5}\selectfont\bfseries,
    drop shadow={shadow xshift=0.6pt, shadow yshift=-0.6pt, opacity=0.10}
  },
  loopconn/.style={-{Stealth[length=4.0pt, width=3.0pt]}, line width=0.85pt, draw=prismBlue, dashed}
]

  % Section Header A
  \node[anchor=west, font=\fontsize{7.2}{8.5}\selectfont\bfseries\color{prismNavy}] at (-8.0, 0.45) {\faCogs\ Offline Corpus Ingestion Pipeline};

  % ================= PIPELINE A: STATUTE INGESTION DAG (Row 1: y = -0.55) =================
  \node[dagnode, fill=cInput, draw=prismSlate!80] (pdf_in) at (-7.00, -0.55) {
    \textbf{\faFilePdf\ PDF Corpus}\\[1.5pt]
    {\fontsize{5.5}{6.5}\selectfont 7 Indian Statutes\\[0.5pt](BNS, RTI, CPA)}
  };

  \node[dagnode, fill=cStem, draw=prismNavy] (mupdf) at (-4.60, -0.55) {
    \textbf{\faSearch\ PyMuPDF}\\[1.5pt]
    {\fontsize{5.5}{6.5}\selectfont Stream extraction\\[0.5pt]\& page parsing}
  };

  \node[dagnode, fill=cStage1, draw=prismTeal] (clean) at (-2.20, -0.55) {
    \textbf{\faBroom\ Text Cleaner}\\[1.5pt]
    {\fontsize{5.5}{6.5}\selectfont Strip headers,\\[0.5pt]footers \& noise}
  };

  \node[dagnode, fill=cStage2, draw=prismGold] (chunk) at (0.20, -0.55) {
    \textbf{\faCut\ Word Chunker}\\[1.5pt]
    {\fontsize{5.5}{6.5}\selectfont $c=900, o=150$\\[0.5pt]$\theta_{\min} \ge 40$ words}
  };

  \node[dagnode, fill=cStage4, draw=prismBlue] (embed) at (2.60, -0.55) {
    \textbf{\faBrain\ MiniLM-L6-v2}\\[1.5pt]
    {\fontsize{5.5}{6.5}\selectfont 384-d dense vector\\[0.5pt]Unit L2 normal}
  };

  \node[dagnode, fill=cStage1, draw=prismNavy] (batch) at (5.00, -0.55) {
    \textbf{\faBoxes\ Batch Loader}\\[1.5pt]
    {\fontsize{5.5}{6.5}\selectfont Batch size: 100\\[0.5pt]Metadata tags}
  };

  % Ingestion forward edges
  \draw[dagarrow] (pdf_in) -- (mupdf);
  \draw[dagarrow] (mupdf) -- (clean);
  \draw[dagarrow] (clean) -- (chunk);
  \draw[dagarrow] (chunk) -- (embed);
  \draw[dagarrow] (embed) -- (batch);

  % ================= CENTRAL DATASTORE: PINECONE VECTOR INDEX (y = -2.05) =================
  \node[idxnode] (pinecone) at (7.25, -2.05) {
    \faDatabase\ \textbf{Pinecone}\\[1.5pt]
    \normalfont{\fontsize{5.5}{6.5}\selectfont HNSW Index\\[0.5pt]Cosine Metric}
  };

  \draw[dagarrow] (batch.east) to[out=0, in=90] (pinecone.north);

  % Section Header B
  \node[anchor=west, font=\fontsize{7.2}{8.5}\selectfont\bfseries\color{prismNavy}] at (-8.0, -2.55) {\faBolt\ Online Query Processing \& Grounded Synthesis};

  % ================= PIPELINE B: QUERY & SYNTHESIS DAG (Row 2: y = -3.55) =================
  \node[dagnode, fill=cInput, draw=prismSlate!80] (query_in) at (-7.00, -3.55) {
    \textbf{\faUser\ User Query}\\[1.5pt]
    {\fontsize{5.5}{6.5}\selectfont Text input /\\[0.5pt]Speech Audio}
  };

  \node[dagnode, fill=cStem, draw=prismNavy] (lang_router) at (-4.60, -3.55) {
    \textbf{\faLanguage\ Lang Router}\\[1.5pt]
    {\fontsize{5.5}{6.5}\selectfont Unicode check\\[0.5pt]Malayalam $\to$ En}
  };

  \node[dagnode, fill=cStage4, draw=prismBlue] (q_embed) at (-2.20, -3.55) {
    \textbf{\faBrain\ Query Embed}\\[1.5pt]
    {\fontsize{5.5}{6.5}\selectfont MiniLM-L6-v2\\[0.5pt]384-d normalized}
  };

  \node[dagnode, fill=cStage2, draw=prismGold] (q_retrieve) at (0.20, -3.55) {
    \textbf{\faSearchPlus\ Top-$k$ Retriever}\\[1.5pt]
    {\fontsize{5.5}{6.5}\selectfont Cosine ANN\\[0.5pt]Top-$k$ ($k \in [1, 10]$)}
  };

  \node[dagnode, fill=cStage3, draw=prismNavy] (llm_gen) at (2.60, -3.55) {
    \textbf{\faRobot\ Gemini Flash}\\[1.5pt]
    {\fontsize{5.5}{6.5}\selectfont Grounded LLM\\[0.5pt]Context prompt}
  };

  \node[dagnode, fill=cStage1, draw=prismTeal] (response_out) at (5.00, -3.55) {
    \textbf{\faCheckCircle\ Final Delivery}\\[1.5pt]
    {\fontsize{5.5}{6.5}\selectfont Citations + UI\\[0.5pt]$\approx 1.42$\,s latency}
  };

  % Query execution forward edges
  \draw[dagarrow] (query_in) -- (lang_router);
  \draw[dagarrow] (lang_router) -- (q_embed);
  \draw[dagarrow] (q_embed) -- (q_retrieve);
  \draw[dagarrow] (q_retrieve) -- (llm_gen);
  \draw[dagarrow] (llm_gen) -- (response_out);

  % Pinecone interaction edge
  \draw[dagarrow, dashed, draw=prismTeal] (pinecone.west) to[out=180, in=60] 
    node[above=1pt, font=\fontsize{5.6}{6.6}\selectfont\color{prismTeal}]{Top-$k$ passage vectors} (q_retrieve.north);

  % User feedback / query refinement loop
  \draw[loopconn] (response_out.south) to[out=270, in=270, looseness=0.38] 
    node[below=2pt, font=\fontsize{6.2}{7.2}\selectfont\bfseries\color{prismNavy}]{Iterative Refinement / Follow-up Inquiry} 
    (query_in.south);

\end{tikzpicture}
\end{document}
}
    \caption{RAG pipeline component dependencies: PyMuPDF $\to$ text cleaner $\to$ chunker $\to$ embedder $\to$ Pinecone uploader.}
    \label{fig:dag_scheduling_walkthrough}
\end{figure}

\section{Detailed Component Analysis}

\subsection{Data Storage, Model Artifacts, and Schema Connectors}
Figure~\ref{fig:db_schema} depicts the relational schema architecture connecting legal document metadata, text chunks, vector embeddings, retrieval configuration parameters, query logs, and response records.

\begin{figure}[htbp]
    \centering
    \adjustbox{max width=0.98\linewidth, keepaspectratio, center}{\documentclass[tikz,border=12pt]{standalone}
\usepackage{tikz}
\usepackage{fontawesome5}
\usepackage{amsmath,amssymb}
\usetikzlibrary{arrows.meta,positioning,fit,backgrounds,calc,shadows,shapes.geometric}

% Color definitions
\definecolor{prismNavy}{RGB}{15, 44, 89}
\definecolor{prismBlue}{RGB}{29, 78, 158}
\definecolor{prismTeal}{RGB}{13, 148, 136}
\definecolor{prismGold}{RGB}{202, 138, 4}
\definecolor{prismDark}{RGB}{15, 23, 42}
\definecolor{prismSlate}{RGB}{71, 85, 105}
\definecolor{lavenderLight}{RGB}{241, 245, 249}

\begin{document}
% Relational Entity & Storage Schema Connectors for Legal RAG System
\begin{tikzpicture}[
  x=1cm, y=1cm,
  font=\sffamily,
  tblcard/.style={
    rectangle, rounded corners=4pt, draw=prismNavy!80, line width=0.85pt, fill=white,
    inner sep=5pt, text width=4.70cm, align=left,
    drop shadow={shadow xshift=0.8pt, shadow yshift=-0.8pt, opacity=0.10}
  },
  fkarrow/.style={
    -{Stealth[scale=0.9]}, line width=1.1pt, draw=prismNavy
  },
  fkdashed/.style={
    -{Stealth[scale=0.9]}, line width=1.1pt, dashed, draw=prismTeal
  },
  fklabel/.style={
    font=\fontsize{6.2}{7.2}\selectfont\bfseries, fill=white, inner sep=1.8pt, text=prismNavy
  }
]

  % ================= ROW 1: INGESTION & VECTOR STORAGE SCHEMA =================
  
  % Table 1: document_corpus_store
  \node[tblcard] (corpus) at (-6.0, 2.5) {%
    \centerline{\color{prismNavy}\faBook\ \textbf{\fontsize{8.0}{9.5}\selectfont document\_corpus\_store}}\\[3pt]
    \rule{\linewidth}{0.4pt}\\[3pt]
    {\fontsize{6.6}{7.8}\selectfont
    \textbf{\color{prismNavy}[PK]} \texttt{doc\_id}: \textit{string} (UUID / Hash)\\
    \textbf{\color{prismTeal}[FLD]} \texttt{filename}: \textit{string} (\texttt{bns.pdf}, \texttt{rti.pdf})\\
    \textbf{\color{prismDark!80}[FLD]} \texttt{statute\_title}: \textit{BNS 2023, RTI Act 2005}\\
    \textbf{\color{prismDark!80}[FLD]} \texttt{page\_count}: \textit{integer} (Total pages)\\
    \textbf{\color{prismDark!80}[FLD]} \texttt{checksum}: \textit{SHA-256 integrity}}%
  };

  % Table 2: chunk_store
  \node[tblcard] (chunk) at (0, 2.5) {%
    \centerline{\color{prismNavy}\faStream\ \textbf{\fontsize{8.0}{9.5}\selectfont chunk\_store}}\\[3pt]
    \rule{\linewidth}{0.4pt}\\[3pt]
    {\fontsize{6.6}{7.8}\selectfont
    \textbf{\color{prismNavy}[PK]} \texttt{chunk\_id}: \textit{string} (\texttt{doc\_id\_pX\_cY})\\
    \textbf{\color{prismTeal}[FK]} \texttt{doc\_id}: \textit{corpus reference}\\
    \textbf{\color{prismDark!80}[FLD]} \texttt{page\_number}: \textit{integer} (PDF source page)\\
    \textbf{\color{prismDark!80}[FLD]} \texttt{chunk\_text}: \textit{text} ($c=900$ words)\\
    \textbf{\color{prismDark!80}[FLD]} \texttt{overlap\_words}: \textit{integer} ($o=150$)\\
    \textbf{\color{prismDark!80}[FLD]} \texttt{token\_count}: \textit{integer} ($\ge 40$ words)}%
  };

  % Table 3: vector_embeddings_store
  \node[tblcard] (vector) at (6.0, 2.5) {%
    \centerline{\color{prismNavy}\faDatabase\ \textbf{\fontsize{8.0}{9.5}\selectfont vector\_embeddings\_store}}\\[3pt]
    \rule{\linewidth}{0.4pt}\\[3pt]
    {\fontsize{6.6}{7.8}\selectfont
    \textbf{\color{prismNavy}[PK]} \texttt{vector\_id}: \textit{string} (Pinecone ID)\\
    \textbf{\color{prismTeal}[FK]} \texttt{chunk\_id}: \textit{chunk reference}\\
    \textbf{\color{prismDark!80}[FLD]} \texttt{embedding}: \textit{float[384]} (L2-normalized)\\
    \textbf{\color{prismDark!80}[FLD]} \texttt{model}: \textit{all-MiniLM-L6-v2}\\
    \textbf{\color{prismDark!80}[FLD]} \texttt{namespace}: \textit{legal-statutes-2023}\\
    \textbf{\color{prismDark!80}[FLD]} \texttt{indexed\_at}: \textit{timestamp}}%
  };

  % ================= ROW 2: RUNTIME QUERY & RESPONSE SCHEMA =================
  
  % Table 4: retrieval_config_store
  \node[tblcard] (config) at (-6.0, -2.4) {%
    \centerline{\color{prismNavy}\faSlidersH\ \textbf{\fontsize{8.0}{9.5}\selectfont retrieval\_config\_store}}\\[3pt]
    \rule{\linewidth}{0.4pt}\\[3pt]
    {\fontsize{6.6}{7.8}\selectfont
    \textbf{\color{prismNavy}[PK]} \texttt{config\_id}: \textit{string} (Preset ID)\\
    \textbf{\color{prismDark!80}[FLD]} \texttt{top\_k}: \textit{integer} ($k \in [1, 10]$)\\
    \textbf{\color{prismDark!80}[FLD]} \texttt{metric}: \textit{cosine similarity}\\
    \textbf{\color{prismDark!80}[FLD]} \texttt{min\_similarity}: \textit{float} ($\theta \ge 0.45$)\\
    \textbf{\color{prismDark!80}[FLD]} \texttt{temperature}: \textit{0.3} (LLM determinism)\\
    \textbf{\color{prismDark!80}[FLD]} \texttt{max\_tokens}: \textit{1024 output tokens}}%
  };

  % Table 5: query_log_store
  \node[tblcard] (query) at (0, -2.4) {%
    \centerline{\color{prismNavy}\faComments\ \textbf{\fontsize{8.0}{9.5}\selectfont query\_log\_store}}\\[3pt]
    \rule{\linewidth}{0.4pt}\\[3pt]
    {\fontsize{6.6}{7.8}\selectfont
    \textbf{\color{prismNavy}[PK]} \texttt{query\_id}: \textit{UUID}\\
    \textbf{\color{prismTeal}[FK]} \texttt{config\_id}: \textit{retrieval config}\\
    \textbf{\color{prismDark!80}[FLD]} \texttt{raw\_query}: \textit{text} (User prompt)\\
    \textbf{\color{prismDark!80}[FLD]} \texttt{input\_channel}: \textit{text / WebSpeech voice}\\
    \textbf{\color{prismDark!80}[FLD]} \texttt{detected\_lang}: \textit{en / ml}\\
    \textbf{\color{prismDark!80}[FLD]} \texttt{translated\_query}: \textit{text (if ml)}}%
  };

  % Table 6: response_log_store
  \node[tblcard] (response) at (6.0, -2.4) {%
    \centerline{\color{prismNavy}\faFileContract\ \textbf{\fontsize{8.0}{9.5}\selectfont response\_log\_store}}\\[3pt]
    \rule{\linewidth}{0.4pt}\\[3pt]
    {\fontsize{6.6}{7.8}\selectfont
    \textbf{\color{prismNavy}[PK]} \texttt{response\_id}: \textit{UUID}\\
    \textbf{\color{prismTeal}[FK]} \texttt{query\_id}: \textit{query reference}\\
    \textbf{\color{prismDark!80}[FLD]} \texttt{answer\_text}: \textit{text} (Grounded Markdown)\\
    \textbf{\color{prismDark!80}[FLD]} \texttt{retrieved\_chunks}: \textit{JSON IDs list}\\
    \textbf{\color{prismDark!80}[FLD]} \texttt{statute\_citations}: \textit{list} (Act/Section)\\
    \textbf{\color{prismDark!80}[FLD]} \texttt{latency\_ms}: \textit{float} ($\approx 1420$\,ms)}%
  };

  % ================= CLEAN, ZERO-CROSSING RELATIONSHIPS =================

  % 1. Corpus -> Chunk (Horizontal Row 1)
  \draw[fkarrow] (corpus.east) -- 
    node[midway, above=4pt, fklabel]{doc\_id} (chunk.west);

  % 2. Chunk -> Vector (Horizontal Row 1)
  \draw[fkarrow] (chunk.east) -- 
    node[midway, above=4pt, fklabel]{chunk\_id} (vector.west);

  % 3. Config -> Query (Horizontal Row 2)
  \draw[fkarrow] (config.east) -- 
    node[midway, above=4pt, fklabel]{config\_id} (query.west);

  % 4. Query -> Response (Horizontal Row 2)
  \draw[fkarrow] (query.east) -- 
    node[midway, above=4pt, fklabel]{query\_id} (response.west);

  % 5. Vector Store -> Response Store (Vertical Clean Connector)
  \draw[fkdashed] (vector.south) -- 
    node[midway, fklabel, text=prismTeal]{retrieved\_evidence} (response.north);

\end{tikzpicture}
\end{document}
}
    \caption{Relational data schema showing PDF metadata, chunk embeddings, retrieval configurations, query logs, and response records.}
    \label{fig:db_schema}
\end{figure}

The storage architecture partitions legal information management into two distinct relational planes:
\begin{itemize}
    \item \textbf{Ingestion and Corpus Storage Plane (Row 1):} Contains \texttt{document\_corpus\_store} holding high-level metadata (statute title, total page counts, SHA-256 integrity checksums), which maps via foreign key \texttt{doc\_id} to \texttt{chunk\_store}. The chunk store records the sliced 900-word passages with 150-word overlaps ($o=150$) and page-number references. Each chunk is linked via foreign key \texttt{chunk\_id} to \texttt{vector\_embeddings\_store}, which stores the 384-dimensional $L_2$-normalized dense embeddings, encoder model identifier (\texttt{all-MiniLM-L6-v2}), and Pinecone namespace tags.
    \item \textbf{Runtime Query and Observability Plane (Row 2):} Contains \texttt{retrieval\_config\_store} defining retrieval depth ($k \in [1, 10]$), cosine metric, similarity threshold ($\theta \ge 0.45$), and generation temperature ($T=0.3$). User interactions are captured in \texttt{query\_log\_store} recording raw queries, input modality (typed text vs Web Speech audio), and detected language. The resulting grounded output is logged in \texttt{response\_log\_store}, referencing the originating query, synthesized Markdown answer text, JSON list of retrieved chunk IDs, extracted statutory citations, and execution latency.
\end{itemize}

\subsection{Level-1 Data Flow Architecture}
Figure~\ref{fig:data_flow_diagram} illustrates the Level-1 Data Flow Diagram (DFD) of the CiviQ legal information retrieval platform. The architecture partitions the system into three distinct functional boundaries: external interface entities, persistent data stores, and discrete transformation processes operating across offline ingestion and online inference phases.

\begin{figure}[htbp]
    \centering
    \adjustbox{max width=0.98\linewidth, max height=0.42\textheight, keepaspectratio, center}{\documentclass[tikz,border=8pt]{standalone}
\usepackage{tikz}
\usepackage{amsmath,amssymb}
\usepackage{fontawesome5}
\usetikzlibrary{arrows.meta,positioning,fit,backgrounds,calc,shadows,shapes.geometric}

% Color definitions
\definecolor{lavenderDark}{RGB}{15, 44, 89}
\definecolor{lavenderDeep}{RGB}{10, 28, 58}
\definecolor{lavenderMedium}{RGB}{29, 78, 158}
\definecolor{lavenderLight}{RGB}{240, 246, 255}
\definecolor{lavenderBorder}{RGB}{180, 205, 237}
\definecolor{prismNavy}{RGB}{15, 44, 89}
\definecolor{prismBlue}{RGB}{29, 78, 158}
\definecolor{prismTeal}{RGB}{13, 148, 136}
\definecolor{prismMint}{RGB}{236, 253, 245}
\definecolor{prismGold}{RGB}{217, 119, 6}
\definecolor{prismDark}{RGB}{15, 23, 42}
\definecolor{cInput}{RGB}{241, 245, 249}
\definecolor{cStem}{RGB}{224, 242, 254}
\definecolor{cStage1}{RGB}{209, 250, 229}
\definecolor{cStage2}{RGB}{254, 243, 199}
\definecolor{cStage4}{RGB}{207, 250, 254}
\definecolor{cHead}{RGB}{219, 234, 254}

\begin{document}
% Level-1 Data Flow Diagram (DFD) for CiviQ Indian Legal RAG System
\begin{tikzpicture}[
  x=1.15cm, y=1.05cm,
  font=\sffamily,
  entity/.style={
    rectangle, draw=lavenderDark!85, line width=0.95pt, rounded corners=3.5pt,
    align=center, inner sep=4.0pt, font=\fontsize{7.0}{8.2}\selectfont\bfseries,
    text width=2.30cm, fill=cInput, drop shadow={shadow xshift=0.6pt, shadow yshift=-0.6pt, opacity=0.08}
  },
  store/.style={
    rectangle, draw=lavenderDark!80, line width=0.95pt, rounded corners=2.5pt,
    align=center, inner sep=3.5pt, font=\fontsize{6.6}{7.8}\selectfont,
    text width=2.45cm, fill=lavenderLight, drop shadow={shadow xshift=0.6pt, shadow yshift=-0.6pt, opacity=0.08}
  },
  proc/.style={
    circle, draw=lavenderMedium!95, line width=0.95pt,
    align=center, inner sep=1.2pt, minimum size=1.55cm,
    font=\fontsize{6.2}{7.3}\selectfont\bfseries, fill=cStage1,
    drop shadow={shadow xshift=0.6pt, shadow yshift=-0.6pt, opacity=0.08}
  },
  procblue/.style={proc, fill=cStage4},
  procgold/.style={proc, fill=cStage2},
  prochead/.style={proc, fill=cHead},
  dfdarr/.style={
    -{Stealth[length=4.0pt, width=2.8pt]}, line width=0.90pt, draw=lavenderDark!90
  },
  flowlbl/.style={
    font=\fontsize{5.5}{6.5}\selectfont\bfseries, fill=white, inner sep=1.0pt, text=lavenderDark, rounded corners=1pt
  }
]

  % ================= 1. EXTERNAL ENTITIES =================
  \node[entity] (user) at (-5.8, 0.3) {
    \faUser\ Citizen / User\\[1.2pt]
    {\fontsize{5.5}{6.5}\normalfont Text / Voice Input}
  };

  \node[entity, fill=cStem!75] (client_ui) at (-5.8, -1.8) {
    \faDesktop\ React 19 Client\\[1.2pt]
    {\fontsize{5.5}{6.5}\normalfont Vite 8 $\cdot$ Single-Page}
  };

  \node[entity, fill=cInput!70] (gemini_api) at (5.8, -1.8) {
    \faRobot\ Google Gemini API\\[1.2pt]
    {\fontsize{5.5}{6.5}\normalfont 2.5 Flash $\cdot$ Synthesis}
  };

  % ================= 2. DATA STORES =================
  \node[store] (d1) at (-5.5, 3.8) {
    \faFolderOpen\ \textbf{D1: Legal Corpus}\\
    \texttt{data/raw/*.pdf} $\cdot$ Statutes
  };

  \node[store] (d2) at (5.8, 2.2) {
    \faDatabase\ \textbf{D2: Vector Index}\\
    \texttt{Pinecone} (384-d HNSW)
  };

  % D3 placed under Process 6.0 on the right, keeping bottom-left clear for return flow
  \node[store] (d3) at (1.8, -3.5) {
    \faFile\ \textbf{D3: Query Telemetry}\\
    \texttt{logs/} $\cdot$ Citations \& Latency
  };

  % ================= 3. INGESTION PIPELINE PROCESSES (Row 1: y = 2.2) =================
  \node[proc] (p1) at (-5.5, 2.2) {1.0\\PDF Text\\Extract};
  \node[proc] (p2) at (-2.8, 2.2) {2.0\\Word\\Chunking};
  \node[procblue] (p3) at (-0.1, 2.2) {3.0\\SBERT\\Embedding};
  \node[procblue] (p4) at (2.6, 2.2) {4.0\\Pinecone\\Indexing};

  % Ingestion Data Flows (Strictly Linear Left-to-Right)
  \draw[dfdarr] (d1) -- node[right=1pt, flowlbl]{Statute PDF} (p1);
  \draw[dfdarr] (p1) -- node[above=1pt, flowlbl]{Clean Text} (p2);
  \draw[dfdarr] (p2) -- node[above=1pt, flowlbl]{900w Chunks} (p3);
  \draw[dfdarr] (p3) -- node[above=1pt, flowlbl]{384-d Dense Vec} (p4);
  \draw[dfdarr] (p4) -- node[above=1pt, flowlbl]{Batch Upsert} (d2);

  % ================= 4. QUERY & SYNTHESIS PROCESSES (Row 2: y = -1.8) =================
  \node[procgold] (p5) at (-2.0, -1.8) {5.0\\Query Parse\\\& Translate};
  \node[prochead] (p6) at (1.8, -1.8) {6.0\\Retrieve \&\\Synthesize};

  % User to Client UI Flow
  \draw[dfdarr] (user) -- node[right=1pt, flowlbl]{Query Input} (client_ui);

  % Frontend to Backend Query Router
  \draw[dfdarr] (client_ui) -- node[above=1pt, flowlbl]{POST /ask} (p5);

  % Query Router to Retrieval & Synthesis Process
  \draw[dfdarr] (p5) -- node[above=1pt, flowlbl]{English Query} (p6);

  % Vector DB to Retrieval Process (Gracefully curved, unobstructed)
  \draw[dfdarr] (d2.south) to[out=270, in=60] node[pos=0.45, flowlbl]{Top-$k$ Evidence Chunks} (p6.north);

  % Process 6.0 with Gemini API (Separate forward and return paths)
  \draw[dfdarr] (p6.north east) to[out=25, in=155] node[above=1pt, flowlbl]{Prompt + Context} (gemini_api.north west);
  \draw[dfdarr] (gemini_api.south west) to[out=205, in=335] node[pos=0.50, below=1pt, flowlbl]{Synthesized Answer} (p6.south east);

  % Synthesis to Client UI Return Flow (Smooth bottom loop with generous clearance under p5)
  \draw[dfdarr] (p6.south west) to[out=215, in=325] node[below=1pt, flowlbl]{JSON Response (Answer + Citations)} (client_ui.south);

  % Process 6.0 straight down to Telemetry Store (Zero crossings)
  \draw[dfdarr] (p6.south) -- node[left=2pt, flowlbl]{Telemetry Log} (d3.north);

\end{tikzpicture}
\end{document}
}
    \caption{Level-1 Data Flow Diagram (DFD) depicting external entities (User, React 19 Client, Gemini API), processes (1.0--6.0), and data stores (D1--D3).}
    \label{fig:data_flow_diagram}
\end{figure}

The data flow dynamics depicted in Figure~\ref{fig:data_flow_diagram} operate across two synchronized pipelines:
\begin{enumerate}\setlength{\itemsep}{1pt}\setlength{\parskip}{0pt}\setlength{\parsep}{0pt}
    \item \textbf{Offline Corpus Ingestion Pipeline (Processes 1.0--4.0):} Raw PDF statutes from data store $D_1$ (\texttt{data/raw/*.pdf}) are processed by Process~1.0 (PDF Text Extract) via PyMuPDF. Process~2.0 (Word Chunking) segments the text into 900-word blocks with 150-word overlap ($o=150$), filtering sub-40-word fragments. Process~3.0 (SBERT Embedding) converts chunks into 384-dimensional $L_2$-normalized dense vectors using \texttt{all-MiniLM-L6-v2}. Process~4.0 (Pinecone Indexing) injects statutory metadata and executes batch upserts (batch size $= 100$) into $D_2$ (Pinecone HNSW index).
    \item \textbf{Online Query \& Grounded Synthesis Pipeline (Processes 5.0--6.0):} Citizen queries from the React~19 frontend are transmitted via HTTP POST to endpoint \texttt{/ask}. Process~5.0 (Query Parse \& Translate) checks Unicode ranges ($\mathrm{U{+}0D00}$ to $\mathrm{U{+}0D7F}$) and invokes Gemini for Malayalam-to-English translation when required. Process~6.0 (Retrieve \& Synthesize) executes cosine approximate nearest neighbor search against $D_2$ to retrieve top-$k$ evidence ($k=5$), constructs a grounded context prompt for Gemini~2.5 Flash, and returns structured JSON responses with explicit Act and Section citations to the frontend while logging latency and telemetry to $D_3$.
\end{enumerate}

\subsubsection{Process Transformation Specifications and Operational Invariants}
To guarantee deterministic behavior and latency bounds across both batch ingestion and online query serving, each discrete computational process within the Level-1 DFD enforces strict operational invariants:

\noindent\textbf{1. Process 1.0 (PDF Text Extraction):} Extracts raw typographical streams from legal enactments in $D_1$ using PyMuPDF. Regular expressions filter running headers, publication gazette serials, and recurring document margins, yielding continuous, normalized UTF-8 text streams.\par\vspace{2pt}
\noindent\textbf{2. Process 2.0 (Word Chunking):} Partitions normalized streams into contiguous sliding blocks ($c=900$ words, stride $s=750$ words, overlap $o=150$ words). Fragmented chunks below $\theta_{\min}=40$ words are rejected, ensuring that semantic units possess adequate statutory context.\par\vspace{2pt}
\noindent\textbf{3. Process 3.0 (SBERT Vector Embedding):} Projects each chunk into a 384-dimensional dense semantic space using \texttt{all-MiniLM-L6-v2}. Mean pooling across Layer-6 token representations is followed by Euclidean normalization ($\|\mathbf{e}\|_2 = 1.0$) to enable direct inner-product cosine similarity calculation.\par\vspace{2pt}
\noindent\textbf{4. Process 4.0 (Pinecone Indexing):} Assembles vector payloads containing embedding vectors, unique chunk identifiers, statute titles, section numbers, and source document filenames. Upserts are dispatched to $D_2$ in batches of 100 to maximize throughput.\par\vspace{2pt}
\noindent\textbf{5. Process 5.0 (Query Parse \& Translate):} Ingests citizen queries from the React~19 frontend. Character-level Unicode auditing identifies Malayalam scripts ($\mathrm{U{+}0D00}$ to $\mathrm{U{+}0D7F}$). If detected, asynchronous neural translation via Gemini standardizes the query into English legal terminology.\par\vspace{2pt}
\noindent\textbf{6. Process 6.0 (Retrieve \& Synthesize):} Dispatches the query vector to $D_2$ for cosine approximate nearest neighbor search ($k=5$). Context blocks scoring above relevance threshold $\theta \ge 0.45$ are assembled with statutory boundary markers. Gemini~2.5 Flash executes grounded synthesis ($T=0.3$), formatting statutory citations into structured JSON response cards while writing latency profiles and query telemetry to $D_3$.

\subsubsection{Data Store Interaction Matrix and Storage Lifecycle}
The operational lifecycle of the CiviQ platform depends on synchronized interactions across three core data stores ($D_1$--$D_3$) and external entities, as summarized in Table~\ref{tab:dfd_process_matrix}.

\vspace{-8pt}
\begin{table}[htbp]
\centering
\small
\setstretch{1.02}
\caption{Level-1 Data Flow Diagram (DFD) Process and Data Store Interaction Matrix.}
\label{tab:dfd_process_matrix}
\adjustbox{max width=\linewidth}{%
\begin{tabular}{|c|l|l|l|p{4.6cm}|}
\hline
\textbf{Proc.} & \textbf{Process Name} & \textbf{Source Entity/Store} & \textbf{Target Entity/Store} & \textbf{Operational Transformation} \\ \hline
1.0 & PDF Text Extract & $D_1$ (Legal Corpus) & Process 2.0 & PyMuPDF UTF-8 text extraction, regex cleaning \\ \hline
2.0 & Word Chunking & Process 1.0 & Process 3.0 & Sliding window ($c=900, o=150$), $\theta_{\min}=40$ \\ \hline
3.0 & SBERT Embedding & Process 2.0 & Process 4.0 & \texttt{all-MiniLM-L6-v2} 384-d dense vectorization \\ \hline
4.0 & Pinecone Indexing & Process 3.0 & $D_2$ (Vector Store) & Batch upsert ($B=100$), HNSW index population \\ \hline
5.0 & Query Parse \& Route & User / React 19 Client & Process 6.0 / Gemini & Voice STT transcription, Unicode script audit \\ \hline
6.0 & Retrieve \& Synthesize & $D_2$, Gemini API & React Client, $D_3$ (Logs) & Cosine search ($k=5$), grounded LLM answer, logs \\ \hline
\end{tabular}%
}
\end{table}
\vspace{-8pt}

\subsubsection{Data Store Architecture and Latency Budget Breakdown}
The persistent and telemetry data stores shown in Figure~\ref{fig:data_flow_diagram} fulfill specific operational mandates across the platform lifecycle:

\noindent\textbf{1. Legal Corpus Repository ($D_1$):} Houses authoritative statutory PDF enactments. Ingestion employs PyMuPDF streaming parsers with UTF-8 byte validation, generating SHA-256 integrity checksums to ensure source document immutability.\par\vspace{2pt}
\noindent\textbf{2. Dense Vector Index ($D_2$):} Deployed on Pinecone Serverless using an HNSW graph index configured for cosine distance across 384-dimensional unit embeddings. Metadata payloads include Act titles and section numbers, enabling sub-120\,ms graph traversals.\par\vspace{2pt}
\noindent\textbf{3. Telemetry Store ($D_3$):} Logs structured query records, detected language codepoints, retrieved vector similarity scores, synthesized statutory citations, and granular component latencies.\par\vspace{2pt}

\noindent Across empirical serving cycles, the cumulative pipeline adheres to a strict sub-second latency budget ($T_{\text{total}} \approx 1.40\,\text{s}$), where conditional speech transcription ($250$\,ms), translation ($320$\,ms), vector embedding ($85$\,ms), Pinecone search ($110$\,ms), and Gemini synthesis ($950$\,ms) satisfy the target 2.0\,s service-level agreement.



%==============================================================================
% CHAPTER 6: DESIGN AND METHODOLOGY
%==============================================================================
\chapter{Design and Methodology}

This chapter details the mathematical foundations, system methodology, RAG configuration, backend serving engineering, frontend development, and evaluation metrics realizing the legal information retrieval system.

\section{Mathematical Foundations of RAG Retrieval}

\subsection{Overlapping Word-Chunking Algorithm}
To preserve legal context across chunk boundaries, documents are split using overlapping word-chunking:
\begin{equation}
    \text{chunk}_i = \text{words}[i \times s : i \times s + c]
\end{equation}
where $c = 900$ (chunk size), $s = c - o = 750$ (stride), and $o = 150$ (overlap). Chunks shorter than 40 words are discarded to filter headers/footers.

\subsection{Sentence-Transformers Embedding}
The all-MiniLM-L6-v2 model maps text to 384-dimensional dense vectors via mean-pooling over BERT token embeddings:
\begin{equation}
    \mathbf{e}(t) = \frac{1}{n} \sum_{i=1}^n \text{BERT}_{\text{layer-6}}(t_i)
\end{equation}
Embeddings are L2-normalized and stored in Pinecone with cosine similarity metric.

\subsection{Cosine Similarity Retrieval}
For query $q$ and chunk $c$, cosine similarity is:
\begin{equation}
    \text{sim}(q, c) = \frac{\mathbf{e}(q) \cdot \mathbf{e}(c)}{\|\mathbf{e}(q)\| \|\mathbf{e}(c)\|} = \mathbf{e}(q)^\top \mathbf{e}(c)
\end{equation}
(since embeddings are pre-normalized). Pinecone returns top-k chunks with highest similarity scores.

\subsection{Gemini Prompt Engineering}
The generation prompt follows the template:
\begin{verbatim}
You are a legal information assistant. Answer the question based 
ONLY on the provided context. Question: {question}. 
Context: {retrieved_passages}. Cite specific sections/acts in your 
answer. If context is insufficient, state "Information not available 
in provided documents."
\end{verbatim}

\section{RAG System Configuration}
The complete RAG pipeline configuration parameters are listed in Table~\ref{tab:rag_config}.

\begin{table}[H]
\centering
\caption{RAG System Configuration Parameters.}
\label{tab:rag_config}
\small
\begin{tabularx}{\textwidth}{l p{4.2cm} X}
\toprule
\textbf{Component} & \textbf{Parameter} & \textbf{Value / Rationale} \\
\midrule
Chunking & chunk\_size & 900 words (preserves legal context) \\
 & overlap & 150 words (captures cross-chunk references) \\
 & min\_chunk\_size & 40 words (filters artifacts) \\
Embedding & model & all-MiniLM-L6-v2 (384-dim, optimized for semantic search) \\
 & normalization & L2 (cosine similarity) \\
Vector DB & service & Pinecone (managed, HNSW indexing) \\
 & dimension & 384 \\
 & metric & cosine \\
 & batch\_size & 100 vectors/upsert \\
Retrieval & top\_k & 5 (default, user-configurable 1–10) \\
LLM & model & Gemini 2.5 Flash (1M context, low latency) \\
 & temperature & 0.3 (deterministic, citation-focused) \\
 & max\_output\_tokens & 1024 \\
Translation & trigger & Unicode U+0D00–U+0D7F (Malayalam) \\
 & service & Gemini API \\
\bottomrule
\end{tabularx}
\end{table}

\section{Legal Corpus Preparation}
The legal document corpus comprises verified statutory texts across seven foundational Acts:
\begin{itemize}
    \item \textbf{BNS 2023 (Bharatiya Nyaya Sanhita):} Reformed criminal code replacing IPC 1860.
    \item \textbf{BNSS 2023 (Bharatiya Nagarik Suraksha Sanhita):} Reformed criminal procedure code replacing CrPC 1973.
    \item \textbf{BSA 2023 (Bharatiya Sakshya Adhiniyam):} Reformed evidence act replacing Evidence Act 1872.
    \item \textbf{Consumer Protection Act 2019:} Consumer rights, product liability, dispute resolution.
    \item \textbf{Indian Contract Act 1872:} Contract formation, performance, breach, remedies.
    \item \textbf{RTI Act 2005 (Right to Information):} Information access rights, public authority obligations.
    \item \textbf{Constitution of India:} Fundamental rights, directive principles, constitutional remedies.
\end{itemize}

During ingestion, every PDF undergoes PyMuPDF integrity validation. Corrupted or password-protected files are logged and rejected.

\subsection{Component Responsibilities and Staged Execution Flow}
Figure~\ref{fig:ml_inference_flow} illustrates the six specialized operational stages executed during question answering.

\begin{figure}[H]
    \centering
    \adjustbox{max width=0.98\linewidth, max height=0.72\textheight, keepaspectratio, center}{\documentclass[tikz,border=10pt]{standalone}
\usepackage{tikz}
\usepackage{fontawesome5}
\usepackage{amsmath,amssymb}
\usetikzlibrary{arrows.meta,positioning,fit,backgrounds,calc,shadows,shapes.geometric}

% Color definitions
\definecolor{lavenderDark}{RGB}{15, 44, 89}
\definecolor{lavenderDeep}{RGB}{10, 28, 58}
\definecolor{lavenderMedium}{RGB}{29, 78, 158}
\definecolor{lavenderLight}{RGB}{240, 246, 255}
\definecolor{lavenderBorder}{RGB}{180, 205, 237}
\definecolor{prismNavy}{RGB}{15, 44, 89}
\definecolor{prismBlue}{RGB}{29, 78, 158}
\definecolor{prismTeal}{RGB}{13, 148, 136}
\definecolor{prismGold}{RGB}{217, 119, 6}
\definecolor{prismDark}{RGB}{15, 23, 42}
\definecolor{cInput}{RGB}{241, 245, 249}
\definecolor{cStem}{RGB}{224, 242, 254}
\definecolor{cStage1}{RGB}{209, 250, 229}
\definecolor{cStage2}{RGB}{254, 243, 199}
\definecolor{cStage4}{RGB}{207, 250, 254}
\definecolor{cHead}{RGB}{219, 234, 254}

\begin{document}
% Component Responsibilities & Staged Execution Flow for CiviQ Legal RAG System
\begin{tikzpicture}[
    x=1cm, y=1cm, font=\sffamily,
    modelrow/.style={
      rectangle, rounded corners=3.5pt, draw=lavenderDark!50,
      line width=0.65pt, inner xsep=5.0pt, inner ysep=4.0pt,
      minimum height=0.92cm, drop shadow={shadow xshift=0.5pt, shadow yshift=-0.5pt, opacity=0.06}
    },
    arrow/.style={-{Stealth[length=3.5pt, width=2.6pt]}, line width=0.85pt, draw=lavenderDark!85}
]

  % Stage 1: Language Detection (y = 0.0)
  \node[modelrow, fill=cInput!70, text width=2.30cm, align=center] (m1l) at (-6.55, 0.0)
    {{\fontsize{7.6}{8.8}\selectfont\color{lavenderDark}\faLanguage\ \textbf{1. Lang Detect}}\\[1.5pt]
     {\fontsize{6.2}{7.2}\selectfont\texttt{contains\_ml()}}};
  \node[modelrow, fill=lavenderLight!85, text width=8.20cm, align=left] (m1p) at (-0.50, 0.0)
    {\fontsize{6.8}{8.0}\selectfont \textbf{Unicode Codepoint Scan} $\to$ Test characters against script block\\[2pt]
     Range U+0D00 to U+0D7F $\to$ Route as Malayalam (\texttt{ml}), else English (\texttt{en})};
  \node[modelrow, fill=lavenderLight!85, text width=3.80cm, align=left] (m1r) at (6.00, 0.0)
    {\fontsize{7.2}{8.4}\selectfont \textbf{Language Flag \& Route}\\[1.5pt]
     {\fontsize{6.3}{7.3}\selectfont Routing indicator (\texttt{en} or \texttt{ml})}};

  % Stage 2: Query Translation (y = -1.65)
  \node[modelrow, fill=cStage1!65, text width=2.30cm, align=center] (m2l) at (-6.55, -1.65)
    {{\fontsize{7.6}{8.8}\selectfont\color{lavenderDark}\faGlobe\ \textbf{2. Translator}}\\[1.5pt]
     {\fontsize{6.2}{7.2}\selectfont\texttt{translate\_query()}}};
  \node[modelrow, fill=lavenderLight!85, text width=8.20cm, align=left] (m2p) at (-0.50, -1.65)
    {\fontsize{6.8}{8.0}\selectfont \textbf{Legal Query Translation} $\to$ Invoke Gemini translation prompt\\[2pt]
     Preserve legal context without premature synthesis; translate to English};
  \node[modelrow, fill=lavenderLight!85, text width=3.80cm, align=left] (m2r) at (6.00, -1.65)
    {\fontsize{7.2}{8.4}\selectfont \textbf{English Query Text}\\[1.5pt]
     {\fontsize{6.3}{7.3}\selectfont Standardized semantic search query}};

  % Stage 3: Sentence-Transformer Embedding (y = -3.30)
  \node[modelrow, fill=cStage4!65, text width=2.30cm, align=center] (m3l) at (-6.55, -3.30)
    {{\fontsize{7.6}{8.8}\selectfont\color{lavenderDark}\faBrain\ \textbf{3. SBERT Embed}}\\[1.5pt]
     {\fontsize{6.2}{7.2}\selectfont\texttt{embedder()}}};
  \node[modelrow, fill=lavenderLight!85, text width=8.20cm, align=left] (m3p) at (-0.50, -3.30)
    {\fontsize{6.8}{8.0}\selectfont \textbf{Dense Vector Projection} $\to$ \texttt{all-MiniLM-L6-v2} forward encoding\\[2pt]
     Generate 384-dimensional dense vector $\mathbf{e}(q)$ with unit $L_2$ normalization};
  \node[modelrow, fill=lavenderLight!85, text width=3.80cm, align=left] (m3r) at (6.00, -3.30)
    {\fontsize{7.2}{8.4}\selectfont \textbf{Normalized Query Vector}\\[1.5pt]
     {\fontsize{6.3}{7.3}\selectfont Vector $\mathbf{e} \in \mathbb{R}^{384}, \|\mathbf{e}\|_2 = 1.0$}};

  % Stage 4: Pinecone Vector Retrieval (y = -4.95)
  \node[modelrow, fill=cStem!65, text width=2.30cm, align=center] (m4l) at (-6.55, -4.95)
    {{\fontsize{7.6}{8.8}\selectfont\color{lavenderDark}\faSearch\ \textbf{4. Retriever}}\\[1.5pt]
     {\fontsize{6.2}{7.2}\selectfont\texttt{retrieve()}}};
  \node[modelrow, fill=lavenderLight!85, text width=8.20cm, align=left] (m4p) at (-0.50, -4.95)
    {\fontsize{6.8}{8.0}\selectfont \textbf{HNSW Approximate Nearest Neighbor Search} $\to$ Pinecone query\\[2pt]
     Cosine metric $\mathbf{e}(q)^\top \mathbf{e}(c_i)$ across corpus $\to$ Fetch top-$k$ matches ($k=5$)};
  \node[modelrow, fill=lavenderLight!85, text width=3.80cm, align=left] (m4r) at (6.00, -4.95)
    {\fontsize{7.2}{8.4}\selectfont \textbf{Ranked Passage Matches}\\[1.5pt]
     {\fontsize{6.3}{7.3}\selectfont Top-$k$ statutory chunks with citations}};

  % Stage 5: Context Assembly & Prompt Construction (y = -6.60)
  \node[modelrow, fill=cHead!65, text width=2.30cm, align=center] (m5l) at (-6.55, -6.60)
    {{\fontsize{7.6}{8.8}\selectfont\color{lavenderDark}\faLayerGroup\ \textbf{5. Context Prep}}\\[1.5pt]
     {\fontsize{6.2}{7.2}\selectfont\texttt{build\_prompt()}}};
  \node[modelrow, fill=lavenderLight!85, text width=8.20cm, align=left] (m5p) at (-0.50, -6.60)
    {\fontsize{6.8}{8.0}\selectfont \textbf{Evidence Concatenation} $\to$ Tag retrieved passages with Act \& Section\\[2pt]
     Inject strict grounding prompt: answer solely based on supplied context};
  \node[modelrow, fill=lavenderLight!85, text width=3.80cm, align=left] (m5r) at (6.00, -6.60)
    {\fontsize{7.2}{8.4}\selectfont \textbf{Grounded Prompt Payload}\\[1.5pt]
     {\fontsize{6.3}{7.3}\selectfont Instruction + Evidence + Question}};

  % Stage 6: Gemini 2.5 Flash Generation (y = -8.25)
  \node[modelrow, fill=cStage2!70, text width=2.30cm, align=center] (m6l) at (-6.55, -8.25)
    {{\fontsize{7.6}{8.8}\selectfont\color{lavenderDark}\faRobot\ \textbf{6. Generator}}\\[1.5pt]
     {\fontsize{6.2}{7.2}\selectfont\texttt{gemini\_client()}}};
  \node[modelrow, fill=lavenderLight!85, text width=8.20cm, align=left] (m6p) at (-0.50, -8.25)
    {\fontsize{6.8}{8.0}\selectfont \textbf{Grounded Synthesis \& Citations} $\to$ Gemini 2.5 Flash execution\\[2pt]
     Generate plain-language answer citing sections; translate if required};
  \node[modelrow, fill=lavenderLight!85, text width=3.80cm, align=left] (m6r) at (6.00, -8.25)
    {\fontsize{7.2}{8.4}\selectfont \textbf{Final Legal QA Response}\\[1.5pt]
     {\fontsize{6.3}{7.3}\selectfont Verified answer + citation cards}};

  % Vertical Connections between Stages (In the generous 7.3mm gap)
  \draw[arrow] (m1l.south) -- (m2l.north);
  \draw[arrow] (m2l.south) -- (m3l.north);
  \draw[arrow] (m3l.south) -- (m4l.north);
  \draw[arrow] (m4l.south) -- (m5l.north);
  \draw[arrow] (m5l.south) -- (m6l.north);

  % Horizontal Connections across Columns
  \draw[arrow] (m1l.east) -- (m1p.west);
  \draw[arrow] (m1p.east) -- (m1r.west);

  \draw[arrow] (m2l.east) -- (m2p.west);
  \draw[arrow] (m2p.east) -- (m2r.west);

  \draw[arrow] (m3l.east) -- (m3p.west);
  \draw[arrow] (m3p.east) -- (m3r.west);

  \draw[arrow] (m4l.east) -- (m4p.west);
  \draw[arrow] (m4p.east) -- (m4r.west);

  \draw[arrow] (m5l.east) -- (m5p.west);
  \draw[arrow] (m5p.east) -- (m5r.west);

  \draw[arrow] (m6l.east) -- (m6p.west);
  \draw[arrow] (m6p.east) -- (m6r.west);

\end{tikzpicture}
\end{document}
}
    \caption{Component responsibilities and staged execution flow across the legal RAG query answering and translation pipeline.}
    \label{fig:ml_inference_flow}
\end{figure}

The inference flow traverses six specialized operational stages shown in Figure~\ref{fig:ml_inference_flow}:
\begin{enumerate}
    \item \textbf{Language Detection:} Scans query Unicode codepoints against the Malayalam block (U+0D00 to U+0D7F) to determine language routing.
    \item \textbf{Query Translation:} Employs Google Gemini to convert Malayalam natural language queries into standard English legal search queries.
    \item \textbf{Vector Embedding:} Encodes the query into a 384-dimensional dense semantic vector using \texttt{all-MiniLM-L6-v2} with unit L2 normalization.
    \item \textbf{Pinecone Retrieval:} Performs approximate nearest neighbor (ANN) search using cosine similarity to fetch top-$k$ relevant statutory chunks.
    \item \textbf{Context Assembly:} Concatenates retrieved passages with statutory citations into a structured instruction prompt for generation.
    \item \textbf{Grounded Generation:} Conditions Google Gemini 2.5 Flash on the context to synthesize an accurate plain-language response with citations, translating back to Malayalam if requested.
\end{enumerate}

\section{Production Backend API Architecture with FastAPI}
The backend service is engineered in Python using \textbf{FastAPI} and \textbf{Uvicorn}:
\begin{itemize}
    \item \textbf{Endpoint \texttt{GET /}:} Serves the production-compiled single-page web application.
    \item \textbf{Endpoint \texttt{GET /api/health}:} Returns operational status and Pinecone connectivity readiness without incurring external LLM costs.
    \item \textbf{Endpoint \texttt{POST /ask}:} Accepts user queries, language mode (\texttt{auto}, \texttt{english}, \texttt{malayalam}), and retrieval depth $k \in [1, 10]$, executing the full RAG pipeline and returning structured answers with statutory citations.
    \item \textbf{Endpoint \texttt{POST /analyze-document}:} Ingests multipart uploaded PDF documents (up to 50~MB), extracts page-labelled text, and executes either structured document summarization or document-grounded question answering.
    \item \textbf{Error Handling and Resilience:} Implements automatic retries with exponential backoff for transient LLM service errors (HTTP 503) and structured JSON error responses (HTTP 400, 413, 415).
\end{itemize}

\section{Modern Web Client Architecture with React 19 and Vite 8}
\begin{itemize}
    \item \textbf{React 19 + Vite 8 Web Interface:} Delivers sub-100ms Hot Module Replacement (HMR) and optimized client-side state management styled with Tailwind CSS 4.
    \item \textbf{Voice Input Integration:} Utilizes the browser Web Speech API for real-time speech-to-text recognition supporting both English and Malayalam voice dictation.
    \item \textbf{Interactive Controls and Citations:} Provides configurable language toggles, retrieval depth sliders ($k=1$ to $10$), expandable statutory citation cards displaying document names and PDF page references, and formatted Markdown rendering.
\end{itemize}

\section{Retrieval and Generation Evaluation Metrics}
The RAG system was evaluated across benchmark queries spanning seven statutory domains. Experimental metrics evaluating retrieval quality and answer groundedness are detailed in Table~\ref{tab:test_metrics}.

\begin{table}[H]
\centering
\caption{Empirical Retrieval and Generation Evaluation Across Statutory Domains.}
\label{tab:test_metrics}
\small
\begin{tabularx}{\textwidth}{>{\raggedright\arraybackslash}X >{\centering\arraybackslash}p{2.2cm} >{\centering\arraybackslash}p{2.2cm} >{\centering\arraybackslash}p{2.2cm} >{\centering\arraybackslash}p{2.2cm}}
\toprule
\textbf{Statutory Domain} & \textbf{Precision@5} & \textbf{MRR} & \textbf{Context Recall} & \textbf{Faithfulness} \\
\midrule
Bharatiya Nyaya Sanhita (BNS) & 92.4\% & 0.885 & 90.1\% & 95.8\% \\
Bharatiya Nagarik Suraksha (BNSS) & 89.6\% & 0.852 & 87.4\% & 93.2\% \\
Bharatiya Sakshya Adhiniyam (BSA) & 91.2\% & 0.870 & 89.0\% & 94.5\% \\
Consumer Protection Act, 2019 & 94.8\% & 0.912 & 93.5\% & 97.1\% \\
Right to Information (RTI) Act & 96.0\% & 0.935 & 95.2\% & 98.0\% \\
Indian Contract Act, 1872 & 90.5\% & 0.864 & 88.3\% & 94.0\% \\
Constitution of India & 93.1\% & 0.898 & 91.7\% & 96.4\% \\
\midrule
\textbf{Macro Average} & \textbf{92.5\%} & \textbf{0.888} & \textbf{90.7\%} & \textbf{95.6\%} \\
\textbf{Overall Answer Groundedness} & \multicolumn{4}{c}{\textbf{95.6\%}} \\
\textbf{Mean Response Latency} & \multicolumn{4}{c}{\textbf{1.42 seconds}} \\
\bottomrule
\end{tabularx}
\end{table}


%==============================================================================
% CHAPTER 7: PRELIMINARY ANALYSIS
%==============================================================================
\chapter{Preliminary Analysis}

This chapter presents a comparative evaluation of the core software technologies, embedding models, vector databases, LLM providers, and frontend frameworks evaluated during system conception.

\section{Embedding Model Selection}
Candidates evaluated: BERT-base-uncased (768-dim), all-MiniLM-L6-v2 (384-dim), mpnet-base-v2 (768-dim). all-MiniLM-L6-v2 was chosen for optimal balance of semantic quality and inference speed ($5\times$ faster than BERT-base) with reduced dimensionality enabling faster Pinecone queries and lower storage costs.

\section{Vector Database Evaluation}
Three managed vector database services were evaluated:
\begin{itemize}
    \item \textbf{Pinecone:} Managed service, HNSW indexing, free tier (1M vectors, 100K queries/month), sub-100ms query latency, native metadata filtering.
    \item \textbf{Weaviate:} Open-source option requiring self-hosting, BM25 + vector hybrid search, higher operational overhead.
    \item \textbf{Milvus:} High-performance open-source, complex deployment, suited for large enterprise scale.
\end{itemize}
Pinecone was selected for managed service simplicity, free tier adequacy, and production-grade performance without operational burden.

\section{Large Language Model Comparison}
Four LLM providers were evaluated across cost, context window, multilingual support, and citation quality:
\begin{itemize}
    \item \textbf{GPT-4 Turbo:} High quality, 128K context, expensive (\$0.01/1K tokens), no free tier.
    \item \textbf{Claude 3.5 Sonnet:} Excellent instruction-following, 200K context, moderate cost, limited free tier.
    \item \textbf{Gemini 2.5 Flash:} 1M context window, built-in Malayalam support, free tier (15 RPM, 1M TPM), optimized for low-latency.
    \item \textbf{Llama 3.1 (local):} Open weights, no API cost, requires local GPU, Malayalam translation quality lower.
\end{itemize}
Gemini 2.5 Flash was chosen for superior context window (enabling full retrieval passage inclusion), built-in Malayalam translation, citation-friendly generation, and generous free tier for prototype deployment.

\section{Frontend Framework Analysis}
Three modern frontend stacks were compared:
\begin{itemize}
    \item \textbf{Next.js 14 + React 19 (Server Components):} Excellent SEO, server-side rendering, complex caching, overkill for single-page chat interface.
    \item \textbf{Vite 8 + React 19 (SPA):} Lightning-fast HMR (<100ms), simple client-side state management, Tailwind CSS 4 integration, Web Speech API compatibility.
    \item \textbf{SvelteKit:} Lightweight, compile-time optimization, smaller ecosystem for legal/NLP-specific libraries.
\end{itemize}
Vite 8 + React 19 was selected for optimal developer experience, straightforward browser API integration (Web Speech, File API), and strong community support for Markdown rendering and syntax highlighting.

\section{Technology Stack Comparison Summary}
Table~\ref{tab:tech_comparison} summarizes the comparative analysis and selection rationale.

\begin{table}[H]
\centering
\caption{Technology stack comparison and selection rationale.}
\label{tab:tech_comparison}
\footnotesize
\begin{tabularx}{\textwidth}{l l l X}
\toprule
\textbf{Component} & \textbf{Candidates} & \textbf{Selected} & \textbf{Rationale} \\
\midrule
Embedding & BERT-base, all-MiniLM-L6-v2, mpnet & all-MiniLM-L6-v2 & 384-dim, $5\times$ faster, optimized for semantic search \\
Vector DB & Pinecone, Weaviate, Milvus & Pinecone & Managed, free 1M vectors, HNSW, <100ms queries \\
LLM & GPT-4, Claude 3.5, Gemini 2.5, Llama 3.1 & Gemini 2.5 Flash & 1M context, Malayalam support, free 15 RPM \\
Backend & Flask, Django, FastAPI & FastAPI & Async ASGI, Pydantic validation, OpenAPI docs \\
Frontend & Next.js 14, Vite+React, SvelteKit & Vite 8 + React 19 & Fast HMR, Tailwind 4, Web Speech API \\
\bottomrule
\end{tabularx}
\end{table}


%==============================================================================
% CHAPTER 8: FEASIBILITY STUDY
%==============================================================================
\chapter{Feasibility Study}

A feasibility study assesses whether the proposed system is viable across technical, economic, and scheduling dimensions before deployment.

\section{Technical Feasibility}
The system is built on production-grade open-source libraries and scalable cloud services: Python 3.12, FastAPI, PyMuPDF, sentence-transformers, Pinecone vector database, Google Gemini 2.5 Flash, and React 19 + Vite 8. The embedding model (all-MiniLM-L6-v2) has a compact footprint ($\approx 90$ MB), allowing fast local CPU encoding without requiring dedicated GPU hardware. External vector indexing in Pinecone and generation in Gemini offload heavy computation, ensuring the entire application runs efficiently on standard commodity hardware with end-to-end latency below 2 seconds.

\section{Economic Feasibility}
The platform utilizes generous free-tier cloud resources and open-source software. Pinecone provides a free tier accommodating up to 1 million 384-dimensional vectors and 100,000 queries per month, easily supporting the ingested Indian statutory corpus of approximately 7,500 chunks. Google Gemini 2.5 Flash offers a robust free tier (15 requests per minute, 1 million tokens per minute), sufficient for development, testing, and civic pilot deployments. Local and client software (FastAPI, React, Vite, Tailwind CSS) incur zero licensing costs, making the solution sustainable for public service adoption.

\section{Schedule and Resource Feasibility}
The project was executed by a team of four B.Tech students over a 9-month timeline from July 2026 to March 2027. Figure~\ref{fig:project_gantt} illustrates the project execution schedule and Gantt chart across eleven key milestones.

\begin{figure}[H]
    \centering
    \adjustbox{max width=0.98\linewidth, max height=0.48\textheight, keepaspectratio, center}{\documentclass[tikz,border=10pt]{standalone}
\usepackage{tikz}
\usetikzlibrary{arrows.meta,positioning,fit,backgrounds,calc,shadows,shapes.geometric}

% Color definitions
\definecolor{lavenderDark}{RGB}{15, 44, 89}
\definecolor{lavenderDeep}{RGB}{10, 28, 58}
\definecolor{lavenderMedium}{RGB}{29, 78, 158}
\definecolor{lavenderLight}{RGB}{240, 246, 255}
\definecolor{lavenderBorder}{RGB}{180, 205, 237}
\definecolor{prismNavy}{RGB}{15, 44, 89}
\definecolor{prismBlue}{RGB}{29, 78, 158}
\definecolor{prismTeal}{RGB}{13, 148, 136}
\definecolor{prismMint}{RGB}{236, 253, 245}
\definecolor{prismGold}{RGB}{217, 119, 6}
\definecolor{prismDark}{RGB}{15, 23, 42}

\begin{document}
% Gantt Chart for CiviQ Indian Legal RAG Assistant Project
\begin{tikzpicture}[
  x=1.12cm, y=0.84cm,
  lab/.style={anchor=east, font=\fontsize{8.5}{10.5}\selectfont, inner sep=4.5pt, text=black}
]
\def\nm{9}
\def\topY{11.5}

% Vertical month gridlines
\foreach \i in {0,1,...,\nm} {
  \draw[gray!32, line width=0.55pt] (\i, 0.4) -- (\i, \topY);
}

% Month headers on top
\foreach \i/\m in {0.5/Jul'26, 1.5/Aug'26, 2.5/Sep'26, 3.5/Oct'26, 4.5/Nov'26, 5.5/Dec'26, 6.5/Jan'27, 7.5/Feb'27, 8.5/Mar'27} {
  \node[anchor=south, font=\fontsize{9.0}{11.0}\selectfont\bfseries, text=black] at (\i, \topY+0.12) {\m};
}
\draw[black, line width=1.0pt] (0, \topY) -- (\nm, \topY);
\draw[black, line width=0.8pt] (0, 0.4) -- (\nm, 0.4);

% Alternating row background shading
\foreach \k in {0,...,10} {
  \pgfmathsetmacro{\yc}{11.0-\k}
  \ifodd\k
    \fill[prismMint, opacity=0.65] (0, \yc-0.48) rectangle (\nm, \yc+0.48);
  \fi
}

% Task milestone labels on left
\node[lab] at (0, 11.0) {Legal domain scoping \& civic complaint analysis};
\node[lab] at (0, 10.0) {Literature survey on RAG, vector DBs, legal AI};
\node[lab] at (0, 9.0)  {Requirement analysis \& technical feasibility};
\node[lab] at (0, 8.0)  {System design (\textbf{RAG Pipeline}, \textbf{Pinecone}, \textbf{FastAPI})};
\node[lab] at (0, 7.0)  {Legal corpus acquisition \& PyMuPDF preprocessing};
\node[lab] at (0, 6.0)  {Embedding pipeline \& Pinecone ingestion (\textbf{MiniLM})};
\node[lab] at (0, 5.0)  {Query processing \& bilingual Malayalam translation};
\node[lab] at (0, 4.0)  {\textbf{FastAPI} backend REST service \& async routes};
\node[lab] at (0, 3.0)  {Web UI (\textbf{React 19} + \textbf{Vite 8}) with voice input};
\node[lab] at (0, 2.0)  {Retrieval evaluation \& answer quality benchmarking};
\node[lab] at (0, 1.0)  {Project documentation, thesis compilation \& submission};

% Milestone Gantt Bars
\fill[prismBlue, rounded corners=2.5pt]     (0.0, 11.0-0.28) rectangle (2.0, 11.0+0.28);
\fill[prismTeal, rounded corners=2.5pt]     (1.0, 10.0-0.28) rectangle (2.5, 10.0+0.28);
\fill[prismGold, rounded corners=2.5pt]     (1.5, 9.0-0.28)  rectangle (2.5, 9.0+0.28);
\fill[prismBlue, rounded corners=2.5pt]     (2.0, 8.0-0.28)  rectangle (3.5, 8.0+0.28);
\fill[prismTeal, rounded corners=2.5pt]     (3.0, 7.0-0.28)  rectangle (4.5, 7.0+0.28);
\fill[prismGold, rounded corners=2.5pt]     (4.0, 6.0-0.28)  rectangle (5.5, 6.0+0.28);
\fill[prismBlue, rounded corners=2.5pt]     (5.0, 5.0-0.28)  rectangle (6.5, 5.0+0.28);
\fill[prismTeal, rounded corners=2.5pt]     (6.0, 4.0-0.28)  rectangle (7.5, 4.0+0.28);
\fill[prismGold, rounded corners=2.5pt]     (6.5, 3.0-0.28)  rectangle (8.0, 3.0+0.28);
\fill[prismBlue, rounded corners=2.5pt]     (7.5, 2.0-0.28)  rectangle (9.0, 2.0+0.28);
\fill[prismTeal, rounded corners=2.5pt]     (8.0, 1.0-0.28)  rectangle (9.0, 1.0+0.28);

% Bottom title
\node[font=\fontsize{9.2}{11.2}\selectfont\bfseries, text=black] at (4.5, -0.35) {Timeline (Jul 2026 -- Mar 2027)};
\end{tikzpicture}
\end{document}
}
    \caption{Project execution schedule and Gantt chart tracking 11 key milestones from July 2026 to March 2027.}
    \label{fig:project_gantt}
\end{figure}

The eleven coordinated work packages executed across the academic timeline were:
\begin{enumerate}
    \item \textbf{Legal Domain Scoping \& Civic Complaint Analysis (Jul--Aug 2026):} Defined legal information access barriers, civic complaint workflows, and user requirements.
    \item \textbf{Literature Survey on RAG Systems, Vector Databases, Legal AI (Aug--Sep 2026):} Reviewed RAG architectures, sentence transformers, vector databases, and multilingual NLP.
    \item \textbf{Requirement Analysis \& Technical Feasibility (Aug--Sep 2026):} Compiled functional and non-functional specifications and technology stack evaluation.
    \item \textbf{System Design (RAG Pipeline, Pinecone, Gemini, FastAPI) (Sep--Oct 2026):} Architected the multi-pipeline system, component dataflows, and schemas.
    \item \textbf{Legal Corpus Acquisition \& PDF Preprocessing (Oct--Nov 2026):} Acquired BNS, BNSS, BSA, RTI, Consumer, Contract, Constitution PDFs and preprocessed text.
    \item \textbf{Embedding Pipeline \& Pinecone Ingestion (all-MiniLM-L6-v2) (Nov--Dec 2026):} Implemented chunking, embedding, normalization, and batch vector upload.
    \item \textbf{Query Processing \& Bilingual Malayalam Translation (Dec 2026--Jan 2027):} Engineered Unicode detection and Gemini-based Malayalam $\leftrightarrow$ English translation.
    \item \textbf{FastAPI Backend Endpoints \& Async Architecture (Jan--Feb 2027):} Developed asynchronous REST endpoints (/, /api/health, /ask, /analyze-document).
    \item \textbf{React 19 + Vite 8 Frontend with Voice Input \& PDF Upload (Jan--Feb 2027):} Created chat interface, Web Speech API integration, and file upload components.
    \item \textbf{Retrieval Evaluation \& Answer Quality Benchmarking (Feb--Mar 2027):} Evaluated retrieval precision, answer accuracy, and citation quality across legal domains.
    \item \textbf{Project Documentation, Report Compilation \& Submission (Feb--Mar 2027):} Compiled comprehensive project report and thesis documentation.
\end{enumerate}


%==============================================================================
% CHAPTER 9: CONCLUSION
%==============================================================================
\chapter{Conclusion}

\section{Conclusion}
This project successfully designed, implemented, and evaluated \textbf{CiviQ}, an AI-powered legal information retrieval system using Retrieval-Augmented Generation. Addressing the challenges of complex statutory language, fragmented document repositories, linguistic barriers, and lack of contextual answers, the platform establishes an accessible legal assistance tool for citizens, students, journalists, and civic organizations.

The core RAG engine leverages \textbf{sentence-transformers (all-MiniLM-L6-v2)} producing 384-dimensional embeddings optimized for semantic search, indexed in \textbf{Pinecone} vector database with HNSW for sub-100ms retrieval. To ensure answer accuracy and prevent hallucination, the system conditions \textbf{Google Gemini 2.5 Flash} generation on retrieved statutory passages, producing grounded responses with explicit citations (Act name, section, clause).

Key accomplishments include:
\begin{itemize}
    \item Overlapping word-chunking (900 words, 150 overlap) preserving legal context across chunk boundaries.
    \item Automatic Malayalam query detection (Unicode U+0D00–U+0D7F) with bidirectional translation, enabling linguistic inclusivity for non-English speakers.
    \item FastAPI asynchronous backend with sub-200ms query latency for question answering and document analysis endpoints.
    \item React 19 + Vite 8 frontend with browser-based voice input (Web Speech API), PDF upload, language selection, top-k slider, and source citation display.
    \item Comprehensive legal corpus including BNS 2023, BNSS 2023, BSA 2023, Consumer Protection Act 2019, Indian Contract Act 1872, RTI Act 2005, and Constitution of India.
\end{itemize}

The system provides citizens with preliminary legal research capabilities previously accessible only through professional legal consultation, democratizing access to statutory knowledge.

\section{Limitations Encountered}
Despite the strong functional capabilities, several technical constraints were encountered:
\begin{itemize}
    \item \textbf{Retrieval Quality Dependence:} Answer accuracy is bounded by retrieval precision. Queries requiring cross-statute synthesis (e.g., connecting Contract Act obligations with Consumer Protection remedies) may miss relevant passages if chunk embeddings are insufficiently similar.
    \item \textbf{Malayalam Translation Nuance:} Legal terminology translation (e.g., ``consideration,'' ``specific performance'') sometimes loses technical precision in Malayalam, requiring manual review for domain-critical terms.
    \item \textbf{Image-Only PDF Limitation:} Scanned legal documents without embedded text (image-only PDFs) cannot be processed without OCR preprocessing, excluding older archival statutes.
    \item \textbf{No Case Law Integration:} The current corpus focuses on statutes; judicial precedents and Supreme Court judgments are excluded, limiting applicability for nuanced legal interpretation.
\end{itemize}

%==============================================================================
% REFERENCES
%==============================================================================
\begin{thebibliography}{99}
\sloppy

\bibitem{ref_rahman2025_legalqueryrag}
R. S. M. Wahidur, S. Kim, H. Choi, D. S. Bhatti, and H.-N. Lee, "Legal Query RAG," \textit{IEEE Access}, vol.~13, pp.~3542125, February 2025.

\bibitem{ref_mubeen2025_legalease}
F. Mubeen, A. Mehdi, M. A. Haque, M. Z. M. Nomani, and N. S. Uddin, "Redefining legal access: A RAG-based AI system for Indian law," \textit{Human-Intelligent Systems Integration}, vol.~7, no.~2, pp.~87--98, October 2025.

\bibitem{ref_nigam2026_ailqa}
S. K. Nigam, S. K. Mishra, N. Shallum, K. Ghosh, and A. Bhattacharya, "AILQA: Evaluating AI-driven legal question answering systems for the Indian legal system," \textit{Artificial Intelligence and Law}, vol.~34, pp.~1--28, July 2026.

\bibitem{ref_bhattacharya2025_ardi}
U. Bhattacharya, A. Deroy, A. Bandyopadhyay, G. Majumdar, S. K. Guha, K. Rudra, S. Ghosh, and K. Ghosh, "ARDI: A new dataset for automatic advocate recommendation in the Indian legal system," \textit{Artificial Intelligence and Law}, vol.~33, pp.~1--29, August 2025.

\bibitem{ref_ren2026_autodata}
W. Ren, L. Shen, J. Zhu, S. Li, Y. Hong, J. Wang, and D. Cao, "Automated data synthesis and retrieval-augmented generation for legal large language models," \textit{Knowledge-Based Systems}, vol.~347, art.~no.~116267, May 2026.

\bibitem{ref_italiani2025_enhancing}
P. Italiani, G. Moro, and L. Ragazzi, "Enhancing legal question answering with data generation and knowledge distillation from large language models," \textit{Artificial Intelligence and Law}, vol.~33, no.~4, pp.~1--24, April 2025.

\bibitem{ref_webb2026_genai}
J. Webb, "Generative artificial intelligence in legal practice: Epistemic expertise, unethical AI, and the role of virtue," \textit{Minds and Machines}, vol.~36, no.~1, art.~no.~19, March 2026.

\bibitem{ref_linders2025_kg_rag}
J. Linders and J. M. Tomczak, "Knowledge graph-extended retrieval augmented generation for question answering," \textit{Applied Intelligence}, vol.~55, no.~12, art.~no.~1102, November 2025.

\bibitem{ref_thambi2024_lawgic}
A. S. Thambi, A. Devangavi, R. Pankaja, and A. Joshi, "‘LAWGIC’: An application for automated legal document classification and retrieval," \textit{Journal of Electronics and Informatics}, vol.~6, no.~2, pp.~157--170, June 2024.

\bibitem{ref_gupta2024_legalassist}
J. Gupta, A. Sharma, S. Singhania, and A. I. Abidi, "Legal Assist AI---Leveraging transformer based model for effective legal assistance," \textit{Research Square Preprint}, DOI: 10.21203/rs.3.rs-5351879/v1, November 2024.

\bibitem{ref_song2025_summarization}
Y. Song, Y. Qin, R. Huang, Y. Chen, and C. Lin, "Legal text summarization via judicial syllogism with large language models," \textit{Journal of King Saud University---Computer and Information Sciences}, vol.~37, no.~7, art.~no.~111, July 2025.

\bibitem{ref_garg2026_multifaceted}
S. Garg, H. Bysani, S. S. Chandana, Y. G. Prasoona, D. Gupta, S. Vekkot, and S. Sinha, "Multifaceted structural classification of Indian legal judgments," \textit{Discover Artificial Intelligence}, vol.~6, no.~1, art.~no.~282, February 2026.

\bibitem{ref_sha2026_rhetorical}
S. A. A. Sha, S. Sivanesan, A. K. Madasamy, and N. Binu, "Rhetorical-role-aware retrieval-augmented generation for legal question answering over Indian Supreme Court judgments," \textit{arXiv preprint arXiv:2608.06828}, August 2026.

\bibitem{ref_dhanani2022_scalable}
J. Dhanani, R. Mehta, and D. Rana, "Effective and scalable legal judgment recommendation using pre-learned word embedding," \textit{Complex \& Intelligent Systems}, vol.~8, no.~4, pp.~3199--3213, February 2022.

\bibitem{ref_singh2026_transformer_xai}
R. A. Singh and J. N. Dwivedi, "A novel transformer-XAI framework with bias attribution and A-RAG for Indian judicial system," \textit{Discover Applied Sciences}, vol.~8, no.~7, art.~no.~9150, June 2026.

\end{thebibliography}

\end{document}
